% !TEX program = lualatex
% ---------------------------------------------------------------------------------------------
%  Readback note 31: Columns and rows as Galois representations: labels at the chosen places, anticyclotomic characters at new primes, families of L-values and the crystalline loss at one place
%
%  Build:  python3 tools/rb.py note <slug>
%
%  One unit of a Lean artifact, rewritten as ordinary mathematics. The prose is hand-written. Lean names,
%  line ranges and counts come from data.tex; quoted Lean statements are copied from the artifact by the build.
%
%  Rules of the text: standard terminology; every symbol is explained where it first appears and listed again
%  in the notation at the end; no letter has two meanings; a proof is complete or says what it takes for granted;
%  Lean names appear only in the LEAN lines and in the concordance.
%
%  The signature rule (the build enforces it): a function, relation or predicate that the artifact defines is
%  written only through its notation macro, which the build makes from the spec (`notation`) with one slot for
%  each explicit argument of the Lean definition. No argument is fixed and dropped, also not "for brevity" and
%  not for one section. Every argument and every hypothesis of a stated declaration stands inside its
%  statement, or in a Setting that the statement cites by number; only implicit and instance arguments (the
%  ambient objects) may stand in the conventions.
% ---------------------------------------------------------------------------------------------
\documentclass[11pt,a4paper]{article}
\ifdefined\pdfvariable\pdfvariable suppressoptionalinfo \numexpr 1+2+4+8\relax\fi
\usepackage[a4paper,hmargin=35mm,top=22mm,bottom=27mm,footskip=12mm,nohead]{geometry}

\usepackage{fontspec}
\directlua{
  luaotfload.add_fallback("gfallback", {"STIXTwoMath-Regular.otf:mode=node;", "DejaVuSans.ttf:mode=node;"})
  luaotfload.add_fallback("gmonofallback", {"DejaVuSansMono.ttf:mode=node;", "STIXTwoMath-Regular.otf:mode=node;"})
}
\setmainfont{NewCM10}[Extension=.otf,UprightFont=*-Book,ItalicFont=*-BookItalic,
  BoldFont=*-Bold,BoldItalicFont=*-BoldItalic,SmallCapsFont=*-Book,SmallCapsFeatures={Letters=SmallCaps},
  RawFeature={fallback=gfallback}]
\setsansfont{NewCMSans10}[Extension=.otf,UprightFont=*-Book,ItalicFont=*-BookOblique,
  BoldFont=*-Bold,BoldItalicFont=*-BoldOblique,RawFeature={fallback=gfallback}]
\setmonofont{JetBrainsMono}[Extension=.otf,UprightFont=*-Regular,ItalicFont=*-Italic,
  BoldFont=*-SemiBold,BoldItalicFont=*-SemiBoldItalic,Scale=MatchLowercase,
  RawFeature={-calt;-liga;fallback=gmonofallback}]
\usepackage{amsmath}
\usepackage{amsthm}
\usepackage[math-style=TeX]{unicode-math}
\setmathfont{NewCMMath-Book.otf}
\emergencystretch=2.5em\relax

\usepackage{xcolor}
\usepackage{tikz}
\usetikzlibrary{calc,arrows.meta,positioning,backgrounds}
\usepackage[protrusion=true,expansion=true]{microtype}
\usepackage{enumitem}
\usepackage{needspace}
\usepackage{booktabs}
\usepackage{array}
\usepackage{longtable}
\usepackage{fancyhdr}
\usepackage{lastpage}
\usepackage{multicol}
\usepackage[hidelinks,pdfusetitle]{hyperref}

\definecolor{ink}{HTML}{0F172A}
\definecolor{inksoft}{HTML}{475569}
\definecolor{inkfaint}{HTML}{94A3B8}
\definecolor{rulegrey}{HTML}{CBD5E1}
\definecolor{arith}{HTML}{C2410C}
\definecolor{est}{HTML}{7C3AED}
\definecolor{ana}{HTML}{0891B2}
\definecolor{ok}{HTML}{047857}
\definecolor{leancol}{HTML}{3730A3}

\title{Columns and rows as Galois representations: labels at the chosen places, anticyclotomic characters at new primes, families of L-values and the crystalline loss at one place}
\author{grh-compact}
\date{}

% ------------------------------------------------------------------------------------- macros
\newcommand{\leanfont}{\ttfamily\fontsize{8.6}{10.5}\selectfont\color{leancol}}
\newcommand{\lean}[1]{\csname LN@#1\endcsname}
\newcommand{\fillin}[1]{{\color{arith}\itshape[#1]}}  % a place of the template that is not written yet; the build reports every one that is left
\newcommand{\leanquote}[1]{\csname LQ@#1\endcsname}
% generated by rb.py note; do not edit
\newcommand{\ConcDecls}{116}
\newcommand{\Pin}{72fead10}
\newcommand{\PinText}{that of commit {\leanfont 72fead10}}
\newcommand{\ArtifactName}{grh-compact}
\expandafter\newcommand\csname LN@IsConjSelfDual\endcsname{\href{https://github.com/lsernau/grh-lean/blob/72fead102f663d5c7a407964a13dc0cf9d977207/grh-compact/src/New/Selmer/AptTupleDefs.lean\#L21-L24}{{\leanfont IsConjSelfDual}}}
\expandafter\newcommand\csname LN@crudeUnit\endcsname{\href{https://github.com/lsernau/grh-lean/blob/72fead102f663d5c7a407964a13dc0cf9d977207/grh-compact/src/New/Selmer/AptTupleDefs.lean\#L27-L27}{{\leanfont crudeUnit}}}
\expandafter\newcommand\csname LN@SelmerLengthLe\endcsname{\href{https://github.com/lsernau/grh-lean/blob/72fead102f663d5c7a407964a13dc0cf9d977207/grh-compact/src/New/Selmer/AptTupleDefs.lean\#L34-L38}{{\leanfont SelmerLengthLe}}}
\expandafter\newcommand\csname LQ@SelmerLengthLe\endcsname{\leanquoteline{0}{def SelmerLengthLe Q m γ (z : intQpBar ℓ) := ∀ᶠ O : CoeffRing ℓ in .atTop, ∃ u ψ' y, y.1 = ↑z ∧}\leanquoteline{2}{IsLatticeIn R u ∧ (∀ g, ψ' g = ψ g) ∧ let ρ := twistRep (.id O) ψ' u; ρ γ = (y : ↑O) • 1 ∧}\leanquoteline{4}{IsUnit (y - 1) ∧ cycChar Ω ℓ γ = 1 ∧ O.normLength (Module.length O ↥(selmerModule ρ}\leanquoteline{4}{(placesDividing \_ ℓ ∪ placesAboveG D ∪ .conjClosure Q) (crysRels ρ) (QuotSpan (ℓ \textasciicircum{} m : ↑O)) ⊓}\leanquoteline{4}{(LinearMap.proj γ).ker)) ≤ 10 * D.f * Fintype.card n * crudeUnit ℓ D + crudeUnit ℓ D}}
\expandafter\newcommand\csname LN@IsColumn\endcsname{\href{https://github.com/lsernau/grh-lean/blob/72fead102f663d5c7a407964a13dc0cf9d977207/grh-compact/src/New/Selmer/AptTupleDefs.lean\#L43-L45}{{\leanfont IsColumn}}}
\expandafter\newcommand\csname LN@IsRow\endcsname{\href{https://github.com/lsernau/grh-lean/blob/72fead102f663d5c7a407964a13dc0cf9d977207/grh-compact/src/New/Selmer/AptTupleDefs.lean\#L50-L59}{{\leanfont IsRow}}}
\expandafter\newcommand\csname LN@chiOneInv\endcsname{\href{https://github.com/lsernau/grh-lean/blob/72fead102f663d5c7a407964a13dc0cf9d977207/grh-compact/src/New/Selmer/AptTupleDefs.lean\#L68-L68}{{\leanfont chiOneInv}}}
\expandafter\newcommand\csname LN@chiOneInvspec\endcsname{\href{https://github.com/lsernau/grh-lean/blob/72fead102f663d5c7a407964a13dc0cf9d977207/grh-compact/src/New/Selmer/AptTupleDefs.lean\#L73-L79}{{\leanfont chiOneInv\_spec}}}
\expandafter\newcommand\csname LN@normalisedRep\endcsname{\href{https://github.com/lsernau/grh-lean/blob/72fead102f663d5c7a407964a13dc0cf9d977207/grh-compact/src/New/Selmer/AptTupleDefs.lean\#L82-L82}{{\leanfont normalisedRep}}}
\expandafter\newcommand\csname LN@twistedLabel\endcsname{\href{https://github.com/lsernau/grh-lean/blob/72fead102f663d5c7a407964a13dc0cf9d977207/grh-compact/src/New/Selmer/AptTupleDefs.lean\#L86-L87}{{\leanfont twistedLabel}}}
\expandafter\newcommand\csname LN@IsApt\endcsname{\href{https://github.com/lsernau/grh-lean/blob/72fead102f663d5c7a407964a13dc0cf9d977207/grh-compact/src/New/Selmer/AptTupleDefs.lean\#L92-L100}{{\leanfont IsApt}}}
\expandafter\newcommand\csname LN@tupleProj\endcsname{\href{https://github.com/lsernau/grh-lean/blob/72fead102f663d5c7a407964a13dc0cf9d977207/grh-compact/src/New/Selmer/AptTupleDefs.lean\#L104-L104}{{\leanfont tupleProj}}}
\expandafter\newcommand\csname LN@HasRow\endcsname{\href{https://github.com/lsernau/grh-lean/blob/72fead102f663d5c7a407964a13dc0cf9d977207/grh-compact/src/New/Selmer/AptTupleDefs.lean\#L108-L110}{{\leanfont HasRow}}}
\expandafter\newcommand\csname LN@IsLabelledColumn\endcsname{\href{https://github.com/lsernau/grh-lean/blob/72fead102f663d5c7a407964a13dc0cf9d977207/grh-compact/src/New/Selmer/AptTupleDefs.lean\#L114-L116}{{\leanfont IsLabelledColumn}}}
\expandafter\newcommand\csname LN@gpowsubonedvdchosenGroupOrder\endcsname{\href{https://github.com/lsernau/grh-lean/blob/72fead102f663d5c7a407964a13dc0cf9d977207/grh-compact/src/New/Selmer/AptTuples.lean\#L15-L20}{{\leanfont g\_pow\_sub\_one\_dvd\_chosenGroupOrder}}}
\expandafter\newcommand\csname LN@ellVallecrudeUnit\endcsname{\href{https://github.com/lsernau/grh-lean/blob/72fead102f663d5c7a407964a13dc0cf9d977207/grh-compact/src/New/Selmer/AptTuples.lean\#L24-L26}{{\leanfont ellVal\_le\_crudeUnit}}}
\expandafter\newcommand\csname LN@cardleofsubsetplacesAboveG\endcsname{\href{https://github.com/lsernau/grh-lean/blob/72fead102f663d5c7a407964a13dc0cf9d977207/grh-compact/src/New/Selmer/AptTuples.lean\#L28-L31}{{\leanfont card\_le\_of\_subset\_placesAboveG}}}
\expandafter\newcommand\csname LN@existschosennotMemplacesDividing\endcsname{\href{https://github.com/lsernau/grh-lean/blob/72fead102f663d5c7a407964a13dc0cf9d977207/grh-compact/src/New/Selmer/AptTuples.lean\#L33-L36}{{\leanfont exists\_chosen\_notMem\_placesDividing}}}
\expandafter\newcommand\csname LN@existsalgHom\endcsname{\href{https://github.com/lsernau/grh-lean/blob/72fead102f663d5c7a407964a13dc0cf9d977207/grh-compact/src/New/Selmer/AptTuples.lean\#L44-L48}{{\leanfont IsLatticeIn.\allowbreak{}exists\_algHom}}}
\expandafter\newcommand\csname LN@continuous\endcsname{\href{https://github.com/lsernau/grh-lean/blob/72fead102f663d5c7a407964a13dc0cf9d977207/grh-compact/src/New/Selmer/AptTuples.lean\#L50-L51}{{\leanfont IsLatticeIn.\allowbreak{}continuous}}}
\expandafter\newcommand\csname LN@eqone\endcsname{\href{https://github.com/lsernau/grh-lean/blob/72fead102f663d5c7a407964a13dc0cf9d977207/grh-compact/src/New/Selmer/AptTuples.lean\#L53-L54}{{\leanfont IsLatticeIn.\allowbreak{}eq\_one}}}
\expandafter\newcommand\csname LN@traceeq\endcsname{\href{https://github.com/lsernau/grh-lean/blob/72fead102f663d5c7a407964a13dc0cf9d977207/grh-compact/src/New/Selmer/AptTuples.lean\#L56-L57}{{\leanfont IsLatticeIn.\allowbreak{}trace\_eq}}}
\expandafter\newcommand\csname LN@cardle\endcsname{\href{https://github.com/lsernau/grh-lean/blob/72fead102f663d5c7a407964a13dc0cf9d977207/grh-compact/src/New/Selmer/AptTuples.lean\#L64-L65}{{\leanfont IsApt.\allowbreak{}card\_le}}}
\expandafter\newcommand\csname LN@isSymmetrictupleProj\endcsname{\href{https://github.com/lsernau/grh-lean/blob/72fead102f663d5c7a407964a13dc0cf9d977207/grh-compact/src/New/Selmer/AptTuples.lean\#L67-L68}{{\leanfont IsApt.\allowbreak{}isSymmetric\_tupleProj}}}
\expandafter\newcommand\csname LN@tupleProjeqselfofmemrange\endcsname{\href{https://github.com/lsernau/grh-lean/blob/72fead102f663d5c7a407964a13dc0cf9d977207/grh-compact/src/New/Selmer/AptTuples.lean\#L70-L71}{{\leanfont IsApt.\allowbreak{}tupleProj\_eq\_self\_of\_mem\_range}}}
\expandafter\newcommand\csname LN@rangetupleProjmeminvtSubmodule\endcsname{\href{https://github.com/lsernau/grh-lean/blob/72fead102f663d5c7a407964a13dc0cf9d977207/grh-compact/src/New/Selmer/AptTuples.lean\#L73-L77}{{\leanfont IsApt.\allowbreak{}range\_tupleProj\_mem\_invtSubmodule}}}
\expandafter\newcommand\csname LN@frequentlyisSplitnotMemplaces\endcsname{\href{https://github.com/lsernau/grh-lean/blob/72fead102f663d5c7a407964a13dc0cf9d977207/grh-compact/src/New/Selmer/AptTuples.lean\#L79-L81}{{\leanfont frequently\_isSplit\_notMem\_places}}}
\expandafter\newcommand\csname LN@existsbaseLevel\endcsname{\href{https://github.com/lsernau/grh-lean/blob/72fead102f663d5c7a407964a13dc0cf9d977207/grh-compact/src/New/Selmer/BaseLevelOfLabel.lean\#L23-L74}{{\leanfont exists\_baseLevel}}}
\expandafter\newcommand\csname LN@obeyingCocycles\endcsname{\href{https://github.com/lsernau/grh-lean/blob/72fead102f663d5c7a407964a13dc0cf9d977207/grh-compact/src/New/Selmer/BlocksAndCocycles.lean\#L16-L17}{{\leanfont obeyingCocycles}}}
\expandafter\newcommand\csname LN@existsgraphsubrepresentation\endcsname{\href{https://github.com/lsernau/grh-lean/blob/72fead102f663d5c7a407964a13dc0cf9d977207/grh-compact/src/New/Selmer/BlocksAndCocycles.lean\#L22-L37}{{\leanfont obeyingCocycles.\allowbreak{}exists\_graph\_subrepresentation}}}
\expandafter\newcommand\csname LN@RowBlockIndex\endcsname{\href{https://github.com/lsernau/grh-lean/blob/72fead102f663d5c7a407964a13dc0cf9d977207/grh-compact/src/New/Selmer/BlocksAndCocycles.lean\#L41-L41}{{\leanfont RowBlockIndex}}}
\expandafter\newcommand\csname LN@blockRep\endcsname{\href{https://github.com/lsernau/grh-lean/blob/72fead102f663d5c7a407964a13dc0cf9d977207/grh-compact/src/New/Selmer/BlocksAndCocycles.lean\#L50-L51}{{\leanfont blockRep}}}
\expandafter\newcommand\csname LN@AuxPrimeData\endcsname{\href{https://github.com/lsernau/grh-lean/blob/72fead102f663d5c7a407964a13dc0cf9d977207/grh-compact/src/New/Selmer/Corepresentable.lean\#L38-L40}{{\leanfont AuxPrimeData}}}
\expandafter\newcommand\csname LN@IsAnticyclotomicAt\endcsname{\href{https://github.com/lsernau/grh-lean/blob/72fead102f663d5c7a407964a13dc0cf9d977207/grh-compact/src/New/Selmer/Corepresentable.lean\#L46-L51}{{\leanfont IsAnticyclotomicAt}}}
\expandafter\newcommand\csname LN@IsAuxiliary\endcsname{\href{https://github.com/lsernau/grh-lean/blob/72fead102f663d5c7a407964a13dc0cf9d977207/grh-compact/src/New/Selmer/Corepresentable.lean\#L55-L55}{{\leanfont IsAuxiliary}}}
\expandafter\newcommand\csname LN@IsFamily\endcsname{\href{https://github.com/lsernau/grh-lean/blob/72fead102f663d5c7a407964a13dc0cf9d977207/grh-compact/src/New/Selmer/Corepresentable.lean\#L60-L68}{{\leanfont IsFamily}}}
\expandafter\newcommand\csname LN@existsisAnticyclotomicAt\endcsname{\href{https://github.com/lsernau/grh-lean/blob/72fead102f663d5c7a407964a13dc0cf9d977207/grh-compact/src/New/Selmer/Corepresentable.lean\#L77-L132}{{\leanfont exists\_isAnticyclotomicAt}}}
\expandafter\newcommand\csname LQ@existsisAnticyclotomicAt\endcsname{\leanquoteline{0}{lemma exists\_isAnticyclotomicAt [ThePrime ℓ'] N (S : Finset (FinPlace E₁)) \{R\} (h₀ : R.Finite) \{z W\}}\leanquoteline{2}{(h : z ∈ normalClosure (inertiaAbove (R \textbackslash{} (.conjClosure S ∪ placesDividing E₁ ℓ')))) (hW :}\leanquoteline{2}{W ∈ nhds z) : ∃ P ∉ R.conjClosure, P.q.Prime ∧ ∃ l : AuxPrimeData E₁ (ℓ' \textasciicircum{} N), l.φ ∈ W ∧}\leanquoteline{2}{l.h l.φ = 1 ∧ ∀ T ⊆ ↑S ∪ placesDividing E₁ ℓ', IsAnticyclotomicAt T P l}}
\expandafter\newcommand\csname LN@evenClasses\endcsname{\href{https://github.com/lsernau/grh-lean/blob/72fead102f663d5c7a407964a13dc0cf9d977207/grh-compact/src/New/Selmer/Corepresentable.lean\#L136-L138}{{\leanfont evenClasses}}}
\expandafter\newcommand\csname LN@IsNewPlaces\endcsname{\href{https://github.com/lsernau/grh-lean/blob/72fead102f663d5c7a407964a13dc0cf9d977207/grh-compact/src/New/Selmer/Corepresentable.lean\#L142-L143}{{\leanfont IsNewPlaces}}}
\expandafter\newcommand\csname LN@withConjInv\endcsname{\href{https://github.com/lsernau/grh-lean/blob/72fead102f663d5c7a407964a13dc0cf9d977207/grh-compact/src/New/Selmer/Corepresentable.lean\#L147-L147}{{\leanfont withConjInv}}}
\expandafter\newcommand\csname LN@IsPrimeTriangle\endcsname{\href{https://github.com/lsernau/grh-lean/blob/72fead102f663d5c7a407964a13dc0cf9d977207/grh-compact/src/New/Selmer/Corepresentable.lean\#L152-L162}{{\leanfont IsPrimeTriangle}}}
\expandafter\newcommand\csname LN@IsAuxTwist\endcsname{\href{https://github.com/lsernau/grh-lean/blob/72fead102f663d5c7a407964a13dc0cf9d977207/grh-compact/src/New/Selmer/Corepresentable.lean\#L166-L167}{{\leanfont IsAuxTwist}}}
\expandafter\newcommand\csname LN@fourdvdcycCharsubone\endcsname{\href{https://github.com/lsernau/grh-lean/blob/72fead102f663d5c7a407964a13dc0cf9d977207/grh-compact/src/New/Selmer/Corepresentable.lean\#L169-L173}{{\leanfont four\_dvd\_cycChar\_sub\_one}}}
\expandafter\newcommand\csname LN@cycInvwithConjInv\endcsname{\href{https://github.com/lsernau/grh-lean/blob/72fead102f663d5c7a407964a13dc0cf9d977207/grh-compact/src/New/Selmer/Corepresentable.lean\#L178-L185}{{\leanfont IsNewPlaces.\allowbreak{}cycInv\_withConjInv}}}
\expandafter\newcommand\csname LN@memselmerModuleoftameeqzero\endcsname{\href{https://github.com/lsernau/grh-lean/blob/72fead102f663d5c7a407964a13dc0cf9d977207/grh-compact/src/New/Selmer/Corepresentable.lean\#L190-L199}{{\leanfont IsNewPlaces.\allowbreak{}mem\_selmerModule\_of\_tame\_eq\_zero}}}
\expandafter\newcommand\csname LN@unionWithConj\endcsname{\href{https://github.com/lsernau/grh-lean/blob/72fead102f663d5c7a407964a13dc0cf9d977207/grh-compact/src/New/Selmer/Corepresentable.lean\#L206-L206}{{\leanfont unionWithConj}}}
\expandafter\newcommand\csname LN@pointChar\endcsname{\href{https://github.com/lsernau/grh-lean/blob/72fead102f663d5c7a407964a13dc0cf9d977207/grh-compact/src/New/Selmer/Corepresentable.lean\#L210-L211}{{\leanfont pointChar}}}
\expandafter\newcommand\csname LN@twistAtPoint\endcsname{\href{https://github.com/lsernau/grh-lean/blob/72fead102f663d5c7a407964a13dc0cf9d977207/grh-compact/src/New/Selmer/Corepresentable.lean\#L215-L216}{{\leanfont twistAtPoint}}}
\expandafter\newcommand\csname LN@reduceModPow\endcsname{\href{https://github.com/lsernau/grh-lean/blob/72fead102f663d5c7a407964a13dc0cf9d977207/grh-compact/src/New/Selmer/Corepresentable.lean\#L219-L220}{{\leanfont reduceModPow}}}
\expandafter\newcommand\csname LN@IsBalanced\endcsname{\href{https://github.com/lsernau/grh-lean/blob/72fead102f663d5c7a407964a13dc0cf9d977207/grh-compact/src/New/Selmer/Corepresentable.lean\#L227-L238}{{\leanfont IsBalanced}}}
\expandafter\newcommand\csname LN@IsPencil\endcsname{\href{https://github.com/lsernau/grh-lean/blob/72fead102f663d5c7a407964a13dc0cf9d977207/grh-compact/src/New/Selmer/Corepresentable.lean\#L243-L250}{{\leanfont IsPencil}}}
\expandafter\newcommand\csname LN@tieRelations\endcsname{\href{https://github.com/lsernau/grh-lean/blob/72fead102f663d5c7a407964a13dc0cf9d977207/grh-compact/src/New/Selmer/Corepresentable.lean\#L254-L256}{{\leanfont tieRelations}}}
\expandafter\newcommand\csname LN@idOrConjInvmeminertiaAbove\endcsname{\href{https://github.com/lsernau/grh-lean/blob/72fead102f663d5c7a407964a13dc0cf9d977207/grh-compact/src/New/Selmer/Corepresentable.lean\#L258-L260}{{\leanfont idOrConjInv\_mem\_inertiaAbove}}}
\expandafter\newcommand\csname LN@existscorepresentingcoords\endcsname{\href{https://github.com/lsernau/grh-lean/blob/72fead102f663d5c7a407964a13dc0cf9d977207/grh-compact/src/New/Selmer/Corepresentable.lean\#L267-L286}{{\leanfont exists\_corepresenting\_coords}}}
\expandafter\newcommand\csname LN@crystallinedefectle\endcsname{\href{https://github.com/lsernau/grh-lean/blob/72fead102f663d5c7a407964a13dc0cf9d977207/grh-compact/src/New/Selmer/CrystallineBound.lean\#L17-L74}{{\leanfont crystalline\_defect\_le}}}
\expandafter\newcommand\csname LN@matrixLattice\endcsname{\href{https://github.com/lsernau/grh-lean/blob/72fead102f663d5c7a407964a13dc0cf9d977207/grh-compact/src/New/Selmer/CrystallineRank.lean\#L22-L23}{{\leanfont matrixLattice}}}
\expandafter\newcommand\csname LN@finrankcrysCocycleseqaddnegjumps\endcsname{\href{https://github.com/lsernau/grh-lean/blob/72fead102f663d5c7a407964a13dc0cf9d977207/grh-compact/src/New/Selmer/CrystallineRank.lean\#L28-L42}{{\leanfont finrank\_crysCocycles\_eq\_add\_neg\_jumps}}}
\expandafter\newcommand\csname LN@normLengthdefectle\endcsname{\href{https://github.com/lsernau/grh-lean/blob/72fead102f663d5c7a407964a13dc0cf9d977207/grh-compact/src/New/Selmer/CrystallineRank.lean\#L47-L56}{{\leanfont normLength\_defect\_le}}}
\expandafter\newcommand\csname LN@localChart\endcsname{\href{https://github.com/lsernau/grh-lean/blob/72fead102f663d5c7a407964a13dc0cf9d977207/grh-compact/src/New/Selmer/DeterminantBound.lean\#L22-L26}{{\leanfont localChart}}}
\expandafter\newcommand\csname LN@coeffEmb\endcsname{\href{https://github.com/lsernau/grh-lean/blob/72fead102f663d5c7a407964a13dc0cf9d977207/grh-compact/src/New/Selmer/DeterminantBound.lean\#L34-L35}{{\leanfont coeffEmb}}}
\expandafter\newcommand\csname LN@ZfAt\endcsname{\href{https://github.com/lsernau/grh-lean/blob/72fead102f663d5c7a407964a13dc0cf9d977207/grh-compact/src/New/Selmer/DeterminantBound.lean\#L39-L39}{{\leanfont ZfAt}}}
\expandafter\newcommand\csname LN@crysRels\endcsname{\href{https://github.com/lsernau/grh-lean/blob/72fead102f663d5c7a407964a13dc0cf9d977207/grh-compact/src/New/Selmer/DeterminantBound.lean\#L43-L43}{{\leanfont crysRels}}}
\expandafter\newcommand\csname LN@blockLoss\endcsname{\href{https://github.com/lsernau/grh-lean/blob/72fead102f663d5c7a407964a13dc0cf9d977207/grh-compact/src/New/Selmer/DeterminantBound.lean\#L47-L48}{{\leanfont blockLoss}}}
\expandafter\newcommand\csname LN@obeyingcrysRels\endcsname{\href{https://github.com/lsernau/grh-lean/blob/72fead102f663d5c7a407964a13dc0cf9d977207/grh-compact/src/New/Selmer/DeterminantBound.lean\#L50-L52}{{\leanfont obeying\_crysRels}}}
\expandafter\newcommand\csname LN@cmSign\endcsname{\href{https://github.com/lsernau/grh-lean/blob/72fead102f663d5c7a407964a13dc0cf9d977207/grh-compact/src/New/Selmer/DeterminantBound.lean\#L56-L56}{{\leanfont cmSign}}}
\expandafter\newcommand\csname LN@complexEmbOf\endcsname{\href{https://github.com/lsernau/grh-lean/blob/72fead102f663d5c7a407964a13dc0cf9d977207/grh-compact/src/New/Selmer/DeterminantBound.lean\#L59-L60}{{\leanfont complexEmbOf}}}
\expandafter\newcommand\csname LN@IsRowType\endcsname{\href{https://github.com/lsernau/grh-lean/blob/72fead102f663d5c7a407964a13dc0cf9d977207/grh-compact/src/New/Selmer/DeterminantBound.lean\#L64-L70}{{\leanfont IsRowType}}}
\expandafter\newcommand\csname LN@lambdaLoss\endcsname{\href{https://github.com/lsernau/grh-lean/blob/72fead102f663d5c7a407964a13dc0cf9d977207/grh-compact/src/New/Selmer/DeterminantBound.lean\#L74-L75}{{\leanfont lambdaLoss}}}
\expandafter\newcommand\csname LN@loss\endcsname{\href{https://github.com/lsernau/grh-lean/blob/72fead102f663d5c7a407964a13dc0cf9d977207/grh-compact/src/New/Selmer/DeterminantBound.lean\#L78-L79}{{\leanfont loss}}}
\expandafter\newcommand\csname LN@coeffEmbsurjective\endcsname{\href{https://github.com/lsernau/grh-lean/blob/72fead102f663d5c7a407964a13dc0cf9d977207/grh-compact/src/New/Selmer/DeterminantBound.lean\#L81-L85}{{\leanfont coeffEmb\_surjective}}}
\expandafter\newcommand\csname LN@tateTwistsmul\endcsname{\href{https://github.com/lsernau/grh-lean/blob/72fead102f663d5c7a407964a13dc0cf9d977207/grh-compact/src/New/Selmer/DeterminantBound.lean\#L87-L89}{{\leanfont tateTwist\_smul}}}
\expandafter\newcommand\csname LN@complexEmbOfeqiff\endcsname{\href{https://github.com/lsernau/grh-lean/blob/72fead102f663d5c7a407964a13dc0cf9d977207/grh-compact/src/New/Selmer/DeterminantBound.lean\#L91-L94}{{\leanfont IsCompletionGroup.\allowbreak{}complexEmbOf\_eq\_iff}}}
\expandafter\newcommand\csname LN@indexeqcardrange\endcsname{\href{https://github.com/lsernau/grh-lean/blob/72fead102f663d5c7a407964a13dc0cf9d977207/grh-compact/src/New/Selmer/DeterminantBound.lean\#L98-L100}{{\leanfont IsCompletionGroup.\allowbreak{}index\_eq\_card\_range}}}
\expandafter\newcommand\csname LN@sumcomplexEmbOfle\endcsname{\href{https://github.com/lsernau/grh-lean/blob/72fead102f663d5c7a407964a13dc0cf9d977207/grh-compact/src/New/Selmer/DeterminantBound.lean\#L102-L106}{{\leanfont sum\_complexEmbOf\_le}}}
\expandafter\newcommand\csname LN@normgramdetleofoperators\endcsname{\href{https://github.com/lsernau/grh-lean/blob/72fead102f663d5c7a407964a13dc0cf9d977207/grh-compact/src/New/Selmer/DeterminantBound.lean\#L111-L127}{{\leanfont norm\_gram\_det\_le\_of\_operators}}}
\expandafter\newcommand\csname LN@isAptauxPair\endcsname{\href{https://github.com/lsernau/grh-lean/blob/72fead102f663d5c7a407964a13dc0cf9d977207/grh-compact/src/New/Selmer/HeckeStableSubspaces.lean\#L22-L30}{{\leanfont isApt\_auxPair}}}
\expandafter\newcommand\csname LN@isAptrankone\endcsname{\href{https://github.com/lsernau/grh-lean/blob/72fead102f663d5c7a407964a13dc0cf9d977207/grh-compact/src/New/Selmer/HeckeStableSubspaces.lean\#L34-L38}{{\leanfont isApt\_rank\_one}}}
\expandafter\newcommand\csname LN@trivialLiftmemrangetupleProjrankone\endcsname{\href{https://github.com/lsernau/grh-lean/blob/72fead102f663d5c7a407964a13dc0cf9d977207/grh-compact/src/New/Selmer/HeckeStableSubspaces.lean\#L42-L54}{{\leanfont trivialLift\_mem\_range\_tupleProj\_rank\_one}}}
\expandafter\newcommand\csname LN@isConjSelfDualofattached\endcsname{\href{https://github.com/lsernau/grh-lean/blob/72fead102f663d5c7a407964a13dc0cf9d977207/grh-compact/src/New/Selmer/HeckeStableSubspaces.lean\#L59-L77}{{\leanfont isConjSelfDual\_of\_attached}}}
\expandafter\newcommand\csname LN@localclausesofconjPlace\endcsname{\href{https://github.com/lsernau/grh-lean/blob/72fead102f663d5c7a407964a13dc0cf9d977207/grh-compact/src/New/Selmer/HeckeStableSubspaces.lean\#L83-L110}{{\leanfont local\_clauses\_of\_conjPlace}}}
\expandafter\newcommand\csname LN@orthogonalmeminvtSubmodulehecke\endcsname{\href{https://github.com/lsernau/grh-lean/blob/72fead102f663d5c7a407964a13dc0cf9d977207/grh-compact/src/New/Selmer/HeckeStableSubspaces.lean\#L112-L115}{{\leanfont orthogonal\_mem\_invtSubmodule\_hecke}}}
\expandafter\newcommand\csname LN@leofisotypicmem\endcsname{\href{https://github.com/lsernau/grh-lean/blob/72fead102f663d5c7a407964a13dc0cf9d977207/grh-compact/src/New/Selmer/HeckeStableSubspaces.lean\#L119-L126}{{\leanfont le\_of\_isotypic\_mem}}}
\expandafter\newcommand\csname LN@existsintegralmultipleevalone\endcsname{\href{https://github.com/lsernau/grh-lean/blob/72fead102f663d5c7a407964a13dc0cf9d977207/grh-compact/src/New/Selmer/HeckeString.lean\#L15-L26}{{\leanfont exists\_integral\_multiple\_eval\_one}}}
\expandafter\newcommand\csname LN@existseventuallynoeigenvalueone\endcsname{\href{https://github.com/lsernau/grh-lean/blob/72fead102f663d5c7a407964a13dc0cf9d977207/grh-compact/src/New/Selmer/HeckeString.lean\#L30-L42}{{\leanfont exists\_eventually\_no\_eigenvalue\_one}}}
\expandafter\newcommand\csname LN@isRootheckestring\endcsname{\href{https://github.com/lsernau/grh-lean/blob/72fead102f663d5c7a407964a13dc0cf9d977207/grh-compact/src/New/Selmer/HeckeString.lean\#L52-L62}{{\leanfont isRoot\_hecke\_string}}}
\expandafter\newcommand\csname LN@eventuallyheckePolyeqstringmul\endcsname{\href{https://github.com/lsernau/grh-lean/blob/72fead102f663d5c7a407964a13dc0cf9d977207/grh-compact/src/New/Selmer/HeckeString.lean\#L68-L87}{{\leanfont eventually\_heckePoly\_eq\_string\_mul}}}
\expandafter\newcommand\csname LN@normpowsubpoweqone\endcsname{\href{https://github.com/lsernau/grh-lean/blob/72fead102f663d5c7a407964a13dc0cf9d977207/grh-compact/src/New/Selmer/LabelledColumns.lean\#L25-L32}{{\leanfont norm\_pow\_sub\_pow\_eq\_one}}}
\expandafter\newcommand\csname LN@existsinertiaeigenvalues\endcsname{\href{https://github.com/lsernau/grh-lean/blob/72fead102f663d5c7a407964a13dc0cf9d977207/grh-compact/src/New/Selmer/LabelledColumns.lean\#L40-L62}{{\leanfont exists\_inertia\_eigenvalues}}}
\expandafter\newcommand\csname LN@isRow\endcsname{\href{https://github.com/lsernau/grh-lean/blob/72fead102f663d5c7a407964a13dc0cf9d977207/grh-compact/src/New/Selmer/LabelledColumns.lean\#L67-L125}{{\leanfont IsLabelledColumn.\allowbreak{}isRow}}}
\expandafter\newcommand\csname LQ@isRow\endcsname{\leanquoteline{0}{lemma IsLabelledColumn.isRow (hc : IsApt O c ϑ) \{R k J\} (h :}\leanquoteline{2}{IsLabelledColumn (m := r) O k J (twistedLabel O c) ∅ R)}\leanquoteline{2}{(hJ : IsRowType D (card r) k J) : IsRow D k J R}}
\expandafter\newcommand\csname LN@normalisedRepquadraticoninertia\endcsname{\href{https://github.com/lsernau/grh-lean/blob/72fead102f663d5c7a407964a13dc0cf9d977207/grh-compact/src/New/Selmer/LabelledColumns.lean\#L130-L155}{{\leanfont normalisedRep\_quadratic\_on\_inertia}}}
\expandafter\newcommand\csname LN@twistavatar\endcsname{\href{https://github.com/lsernau/grh-lean/blob/72fead102f663d5c7a407964a13dc0cf9d977207/grh-compact/src/New/Selmer/LabelledColumns.lean\#L167-L182}{{\leanfont IsLabelledColumn.\allowbreak{}twist\_avatar}}}
\expandafter\newcommand\csname LN@isLabelledColumnofhasProductLabel\endcsname{\href{https://github.com/lsernau/grh-lean/blob/72fead102f663d5c7a407964a13dc0cf9d977207/grh-compact/src/New/Selmer/LabelledColumns.lean\#L187-L197}{{\leanfont isLabelledColumn\_of\_hasProductLabel}}}
\expandafter\newcommand\csname LN@algebraMapmemcycloField\endcsname{\href{https://github.com/lsernau/grh-lean/blob/72fead102f663d5c7a407964a13dc0cf9d977207/grh-compact/src/New/Selmer/LevelPlacesChosen.lean\#L20-L21}{{\leanfont Data.\allowbreak{}algebraMap\_mem\_cycloField}}}
\expandafter\newcommand\csname LN@AbsGal\endcsname{\href{https://github.com/lsernau/grh-lean/blob/72fead102f663d5c7a407964a13dc0cf9d977207/grh-compact/src/New/Selmer/LevelPlacesChosen.lean\#L25-L25}{{\leanfont AbsGal}}}
\expandafter\newcommand\csname LN@placesAboveG\endcsname{\href{https://github.com/lsernau/grh-lean/blob/72fead102f663d5c7a407964a13dc0cf9d977207/grh-compact/src/New/Selmer/LevelPlacesChosen.lean\#L28-L28}{{\leanfont placesAboveG}}}
\expandafter\newcommand\csname LN@badPlaces\endcsname{\href{https://github.com/lsernau/grh-lean/blob/72fead102f663d5c7a407964a13dc0cf9d977207/grh-compact/src/New/Selmer/LevelPlacesChosen.lean\#L32-L32}{{\leanfont badPlaces}}}
\expandafter\newcommand\csname LN@placesDividingunionsubsetbadPlaces\endcsname{\href{https://github.com/lsernau/grh-lean/blob/72fead102f663d5c7a407964a13dc0cf9d977207/grh-compact/src/New/Selmer/LevelPlacesChosen.lean\#L34-L36}{{\leanfont placesDividing\_union\_subset\_badPlaces}}}
\expandafter\newcommand\csname LN@starmemplacesAboveG\endcsname{\href{https://github.com/lsernau/grh-lean/blob/72fead102f663d5c7a407964a13dc0cf9d977207/grh-compact/src/New/Selmer/LevelPlacesChosen.lean\#L40-L41}{{\leanfont star\_mem\_placesAboveG}}}
\expandafter\newcommand\csname LN@normnatCastlenormsubofpoweqone\endcsname{\href{https://github.com/lsernau/grh-lean/blob/72fead102f663d5c7a407964a13dc0cf9d977207/grh-compact/src/New/Selmer/LevelPlacesChosen.lean\#L44-L46}{{\leanfont norm\_natCast\_le\_norm\_sub\_of\_pow\_eq\_one}}}
\expandafter\newcommand\csname LN@continuouscycInv\endcsname{\href{https://github.com/lsernau/grh-lean/blob/72fead102f663d5c7a407964a13dc0cf9d977207/grh-compact/src/New/Selmer/LevelPlacesChosen.lean\#L48-L49}{{\leanfont continuous\_cycInv}}}
\expandafter\newcommand\csname LN@cycInveqofisFrob\endcsname{\href{https://github.com/lsernau/grh-lean/blob/72fead102f663d5c7a407964a13dc0cf9d977207/grh-compact/src/New/Selmer/LevelPlacesChosen.lean\#L51-L53}{{\leanfont cycInv\_eq\_of\_isFrob}}}
\expandafter\newcommand\csname LN@invAvatar\endcsname{\href{https://github.com/lsernau/grh-lean/blob/72fead102f663d5c7a407964a13dc0cf9d977207/grh-compact/src/New/Selmer/LevelPlacesChosen.lean\#L59-L60}{{\leanfont Data.\allowbreak{}TameChar.\allowbreak{}invAvatar}}}
\expandafter\newcommand\csname LN@avatarconjGalmul\endcsname{\href{https://github.com/lsernau/grh-lean/blob/72fead102f663d5c7a407964a13dc0cf9d977207/grh-compact/src/New/Selmer/LevelPlacesChosen.lean\#L64-L73}{{\leanfont avatar\_conjGal\_mul}}}
\expandafter\newcommand\csname LN@isLabelledAtavatar\endcsname{\href{https://github.com/lsernau/grh-lean/blob/72fead102f663d5c7a407964a13dc0cf9d977207/grh-compact/src/New/Selmer/LevelPlacesChosen.lean\#L78-L84}{{\leanfont isLabelledAt\_avatar}}}
\expandafter\newcommand\csname LN@isChosenchosenPlace\endcsname{\href{https://github.com/lsernau/grh-lean/blob/72fead102f663d5c7a407964a13dc0cf9d977207/grh-compact/src/New/Selmer/LevelPlacesChosen.lean\#L86-L89}{{\leanfont isChosen\_chosenPlace}}}
\expandafter\newcommand\csname LN@blockdefectleatplace\endcsname{\href{https://github.com/lsernau/grh-lean/blob/72fead102f663d5c7a407964a13dc0cf9d977207/grh-compact/src/New/Selmer/OnePlaceDefect.lean\#L15-L84}{{\leanfont block\_defect\_le\_at\_place}}}
\expandafter\newcommand\csname LN@fixesemboffixesrootsOfUnity\endcsname{\href{https://github.com/lsernau/grh-lean/blob/72fead102f663d5c7a407964a13dc0cf9d977207/grh-compact/src/New/Selmer/ThickSpacing.lean\#L12-L16}{{\leanfont fixes\_emb\_of\_fixes\_rootsOfUnity}}}
\expandafter\newcommand\csname LN@cmSignspec\endcsname{\href{https://github.com/lsernau/grh-lean/blob/72fead102f663d5c7a407964a13dc0cf9d977207/grh-compact/src/New/Selmer/ThickSpacing.lean\#L18-L20}{{\leanfont cmSign\_spec}}}
\expandafter\newcommand\csname LN@placeClassratio\endcsname{\href{https://github.com/lsernau/grh-lean/blob/72fead102f663d5c7a407964a13dc0cf9d977207/grh-compact/src/New/Selmer/ThickSpacing.lean\#L30-L39}{{\leanfont placeClass\_ratio}}}
\expandafter\newcommand\csname LN@notthickofshortarc\endcsname{\href{https://github.com/lsernau/grh-lean/blob/72fead102f663d5c7a407964a13dc0cf9d977207/grh-compact/src/New/Selmer/ThickSpacing.lean\#L45-L57}{{\leanfont not\_thick\_of\_short\_arc}}}
\expandafter\newcommand\csname LN@cmSigneqofthick\endcsname{\href{https://github.com/lsernau/grh-lean/blob/72fead102f663d5c7a407964a13dc0cf9d977207/grh-compact/src/New/Selmer/ThickSpacing.lean\#L64-L75}{{\leanfont cmSign\_eq\_of\_thick}}}
\expandafter\newcommand\csname LN@existsisIntegralMeasure\endcsname{\href{https://github.com/lsernau/grh-lean/blob/72fead102f663d5c7a407964a13dc0cf9d977207/grh-compact/src/New/Selmer/ValueFamilies.lean\#L30-L80}{{\leanfont Data.\allowbreak{}Pair.\allowbreak{}exists\_isIntegralMeasure}}}
\expandafter\newcommand\csname LN@coprimeqofisAuxiliary\endcsname{\href{https://github.com/lsernau/grh-lean/blob/72fead102f663d5c7a407964a13dc0cf9d977207/grh-compact/src/New/Selmer/ValueFamilies.lean\#L82-L85}{{\leanfont coprime\_q\_of\_isAuxiliary}}}
\expandafter\newcommand\csname LN@tameCharOfAuxiliary\endcsname{\href{https://github.com/lsernau/grh-lean/blob/72fead102f663d5c7a407964a13dc0cf9d977207/grh-compact/src/New/Selmer/ValueFamilies.lean\#L89-L90}{{\leanfont tameCharOfAuxiliary}}}
\expandafter\newcommand\csname LN@invAvatareqofartinMapeq\endcsname{\href{https://github.com/lsernau/grh-lean/blob/72fead102f663d5c7a407964a13dc0cf9d977207/grh-compact/src/New/Selmer/ValueFamilies.lean\#L92-L95}{{\leanfont invAvatar\_eq\_of\_artinMap\_eq}}}
\expandafter\newcommand\csname LN@existstameCharinvAvatareq\endcsname{\href{https://github.com/lsernau/grh-lean/blob/72fead102f663d5c7a407964a13dc0cf9d977207/grh-compact/src/New/Selmer/ValueFamilies.lean\#L97-L106}{{\leanfont exists\_tameChar\_invAvatar\_eq}}}
\expandafter\newcommand\csname LN@existsisFamily\endcsname{\href{https://github.com/lsernau/grh-lean/blob/72fead102f663d5c7a407964a13dc0cf9d977207/grh-compact/src/New/Selmer/ValueFamilies.lean\#L111-L153}{{\leanfont exists\_isFamily}}}
\newcommand{\ConcordanceRows}{\crow{IsConjSelfDual}{\href{https://github.com/lsernau/grh-lean/blob/72fead102f663d5c7a407964a13dc0cf9d977207/grh-compact/src/New/Selmer/AptTupleDefs.lean\#L21-L24}{{\leanfont IsConjSelfDual}}}{AptTupleDefs.lean}{21--24}
\crow{crudeUnit}{\href{https://github.com/lsernau/grh-lean/blob/72fead102f663d5c7a407964a13dc0cf9d977207/grh-compact/src/New/Selmer/AptTupleDefs.lean\#L27-L27}{{\leanfont crudeUnit}}}{AptTupleDefs.lean}{27--27}
\crow{SelmerLengthLe}{\href{https://github.com/lsernau/grh-lean/blob/72fead102f663d5c7a407964a13dc0cf9d977207/grh-compact/src/New/Selmer/AptTupleDefs.lean\#L34-L38}{{\leanfont SelmerLengthLe}}}{AptTupleDefs.lean}{34--38}
\crow{IsColumn}{\href{https://github.com/lsernau/grh-lean/blob/72fead102f663d5c7a407964a13dc0cf9d977207/grh-compact/src/New/Selmer/AptTupleDefs.lean\#L43-L45}{{\leanfont IsColumn}}}{AptTupleDefs.lean}{43--45}
\crow{IsRow}{\href{https://github.com/lsernau/grh-lean/blob/72fead102f663d5c7a407964a13dc0cf9d977207/grh-compact/src/New/Selmer/AptTupleDefs.lean\#L50-L59}{{\leanfont IsRow}}}{AptTupleDefs.lean}{50--59}
\crow{chiOneInv}{\href{https://github.com/lsernau/grh-lean/blob/72fead102f663d5c7a407964a13dc0cf9d977207/grh-compact/src/New/Selmer/AptTupleDefs.lean\#L68-L68}{{\leanfont chiOneInv}}}{AptTupleDefs.lean}{68--68}
\crow{chiOneInvspec}{\href{https://github.com/lsernau/grh-lean/blob/72fead102f663d5c7a407964a13dc0cf9d977207/grh-compact/src/New/Selmer/AptTupleDefs.lean\#L73-L79}{{\leanfont chiOneInv\_spec}}}{AptTupleDefs.lean}{73--79}
\crow{normalisedRep}{\href{https://github.com/lsernau/grh-lean/blob/72fead102f663d5c7a407964a13dc0cf9d977207/grh-compact/src/New/Selmer/AptTupleDefs.lean\#L82-L82}{{\leanfont normalisedRep}}}{AptTupleDefs.lean}{82--82}
\crow{twistedLabel}{\href{https://github.com/lsernau/grh-lean/blob/72fead102f663d5c7a407964a13dc0cf9d977207/grh-compact/src/New/Selmer/AptTupleDefs.lean\#L86-L87}{{\leanfont twistedLabel}}}{AptTupleDefs.lean}{86--87}
\crow{IsApt}{\href{https://github.com/lsernau/grh-lean/blob/72fead102f663d5c7a407964a13dc0cf9d977207/grh-compact/src/New/Selmer/AptTupleDefs.lean\#L92-L100}{{\leanfont IsApt}}}{AptTupleDefs.lean}{92--100}
\crow{tupleProj}{\href{https://github.com/lsernau/grh-lean/blob/72fead102f663d5c7a407964a13dc0cf9d977207/grh-compact/src/New/Selmer/AptTupleDefs.lean\#L104-L104}{{\leanfont tupleProj}}}{AptTupleDefs.lean}{104--104}
\crow{HasRow}{\href{https://github.com/lsernau/grh-lean/blob/72fead102f663d5c7a407964a13dc0cf9d977207/grh-compact/src/New/Selmer/AptTupleDefs.lean\#L108-L110}{{\leanfont HasRow}}}{AptTupleDefs.lean}{108--110}
\crow{IsLabelledColumn}{\href{https://github.com/lsernau/grh-lean/blob/72fead102f663d5c7a407964a13dc0cf9d977207/grh-compact/src/New/Selmer/AptTupleDefs.lean\#L114-L116}{{\leanfont IsLabelledColumn}}}{AptTupleDefs.lean}{114--116}
\crow{gpowsubonedvdchosenGroupOrder}{\href{https://github.com/lsernau/grh-lean/blob/72fead102f663d5c7a407964a13dc0cf9d977207/grh-compact/src/New/Selmer/AptTuples.lean\#L15-L20}{{\leanfont g\_pow\_sub\_one\_dvd\_chosenGroupOrder}}}{AptTuples.lean}{15--20}
\crow{ellVallecrudeUnit}{\href{https://github.com/lsernau/grh-lean/blob/72fead102f663d5c7a407964a13dc0cf9d977207/grh-compact/src/New/Selmer/AptTuples.lean\#L24-L26}{{\leanfont ellVal\_le\_crudeUnit}}}{AptTuples.lean}{24--26}
\crow{cardleofsubsetplacesAboveG}{\href{https://github.com/lsernau/grh-lean/blob/72fead102f663d5c7a407964a13dc0cf9d977207/grh-compact/src/New/Selmer/AptTuples.lean\#L28-L31}{{\leanfont card\_le\_of\_subset\_placesAboveG}}}{AptTuples.lean}{28--31}
\crow{existschosennotMemplacesDividing}{\href{https://github.com/lsernau/grh-lean/blob/72fead102f663d5c7a407964a13dc0cf9d977207/grh-compact/src/New/Selmer/AptTuples.lean\#L33-L36}{{\leanfont exists\_chosen\_notMem\_placesDividing}}}{AptTuples.lean}{33--36}
\crow{existsalgHom}{\href{https://github.com/lsernau/grh-lean/blob/72fead102f663d5c7a407964a13dc0cf9d977207/grh-compact/src/New/Selmer/AptTuples.lean\#L44-L48}{{\leanfont IsLatticeIn.\allowbreak{}exists\_algHom}}}{AptTuples.lean}{44--48}
\crow{continuous}{\href{https://github.com/lsernau/grh-lean/blob/72fead102f663d5c7a407964a13dc0cf9d977207/grh-compact/src/New/Selmer/AptTuples.lean\#L50-L51}{{\leanfont IsLatticeIn.\allowbreak{}continuous}}}{AptTuples.lean}{50--51}
\crow{eqone}{\href{https://github.com/lsernau/grh-lean/blob/72fead102f663d5c7a407964a13dc0cf9d977207/grh-compact/src/New/Selmer/AptTuples.lean\#L53-L54}{{\leanfont IsLatticeIn.\allowbreak{}eq\_one}}}{AptTuples.lean}{53--54}
\crow{traceeq}{\href{https://github.com/lsernau/grh-lean/blob/72fead102f663d5c7a407964a13dc0cf9d977207/grh-compact/src/New/Selmer/AptTuples.lean\#L56-L57}{{\leanfont IsLatticeIn.\allowbreak{}trace\_eq}}}{AptTuples.lean}{56--57}
\crow{cardle}{\href{https://github.com/lsernau/grh-lean/blob/72fead102f663d5c7a407964a13dc0cf9d977207/grh-compact/src/New/Selmer/AptTuples.lean\#L64-L65}{{\leanfont IsApt.\allowbreak{}card\_le}}}{AptTuples.lean}{64--65}
\crow{isSymmetrictupleProj}{\href{https://github.com/lsernau/grh-lean/blob/72fead102f663d5c7a407964a13dc0cf9d977207/grh-compact/src/New/Selmer/AptTuples.lean\#L67-L68}{{\leanfont IsApt.\allowbreak{}isSymmetric\_tupleProj}}}{AptTuples.lean}{67--68}
\crow{tupleProjeqselfofmemrange}{\href{https://github.com/lsernau/grh-lean/blob/72fead102f663d5c7a407964a13dc0cf9d977207/grh-compact/src/New/Selmer/AptTuples.lean\#L70-L71}{{\leanfont IsApt.\allowbreak{}tupleProj\_eq\_self\_of\_mem\_range}}}{AptTuples.lean}{70--71}
\crow{rangetupleProjmeminvtSubmodule}{\href{https://github.com/lsernau/grh-lean/blob/72fead102f663d5c7a407964a13dc0cf9d977207/grh-compact/src/New/Selmer/AptTuples.lean\#L73-L77}{{\leanfont IsApt.\allowbreak{}range\_tupleProj\_mem\_invtSubmodule}}}{AptTuples.lean}{73--77}
\crow{frequentlyisSplitnotMemplaces}{\href{https://github.com/lsernau/grh-lean/blob/72fead102f663d5c7a407964a13dc0cf9d977207/grh-compact/src/New/Selmer/AptTuples.lean\#L79-L81}{{\leanfont frequently\_isSplit\_notMem\_places}}}{AptTuples.lean}{79--81}
\crow{existsbaseLevel}{\href{https://github.com/lsernau/grh-lean/blob/72fead102f663d5c7a407964a13dc0cf9d977207/grh-compact/src/New/Selmer/BaseLevelOfLabel.lean\#L23-L74}{{\leanfont exists\_baseLevel}}}{BaseLevelOfLabel.lean}{23--74}
\crow{obeyingCocycles}{\href{https://github.com/lsernau/grh-lean/blob/72fead102f663d5c7a407964a13dc0cf9d977207/grh-compact/src/New/Selmer/BlocksAndCocycles.lean\#L16-L17}{{\leanfont obeyingCocycles}}}{BlocksAndCocycles.lean}{16--17}
\crow{existsgraphsubrepresentation}{\href{https://github.com/lsernau/grh-lean/blob/72fead102f663d5c7a407964a13dc0cf9d977207/grh-compact/src/New/Selmer/BlocksAndCocycles.lean\#L22-L37}{{\leanfont obeyingCocycles.\allowbreak{}exists\_graph\_subrepresentation}}}{BlocksAndCocycles.lean}{22--37}
\crow{RowBlockIndex}{\href{https://github.com/lsernau/grh-lean/blob/72fead102f663d5c7a407964a13dc0cf9d977207/grh-compact/src/New/Selmer/BlocksAndCocycles.lean\#L41-L41}{{\leanfont RowBlockIndex}}}{BlocksAndCocycles.lean}{41--41}
\crow{blockRep}{\href{https://github.com/lsernau/grh-lean/blob/72fead102f663d5c7a407964a13dc0cf9d977207/grh-compact/src/New/Selmer/BlocksAndCocycles.lean\#L50-L51}{{\leanfont blockRep}}}{BlocksAndCocycles.lean}{50--51}
\crow{AuxPrimeData}{\href{https://github.com/lsernau/grh-lean/blob/72fead102f663d5c7a407964a13dc0cf9d977207/grh-compact/src/New/Selmer/Corepresentable.lean\#L38-L40}{{\leanfont AuxPrimeData}}}{Corepresentable.lean}{38--40}
\crow{IsAnticyclotomicAt}{\href{https://github.com/lsernau/grh-lean/blob/72fead102f663d5c7a407964a13dc0cf9d977207/grh-compact/src/New/Selmer/Corepresentable.lean\#L46-L51}{{\leanfont IsAnticyclotomicAt}}}{Corepresentable.lean}{46--51}
\crow{IsAuxiliary}{\href{https://github.com/lsernau/grh-lean/blob/72fead102f663d5c7a407964a13dc0cf9d977207/grh-compact/src/New/Selmer/Corepresentable.lean\#L55-L55}{{\leanfont IsAuxiliary}}}{Corepresentable.lean}{55--55}
\crow{IsFamily}{\href{https://github.com/lsernau/grh-lean/blob/72fead102f663d5c7a407964a13dc0cf9d977207/grh-compact/src/New/Selmer/Corepresentable.lean\#L60-L68}{{\leanfont IsFamily}}}{Corepresentable.lean}{60--68}
\crow{existsisAnticyclotomicAt}{\href{https://github.com/lsernau/grh-lean/blob/72fead102f663d5c7a407964a13dc0cf9d977207/grh-compact/src/New/Selmer/Corepresentable.lean\#L77-L132}{{\leanfont exists\_isAnticyclotomicAt}}}{Corepresentable.lean}{77--132}
\crow{evenClasses}{\href{https://github.com/lsernau/grh-lean/blob/72fead102f663d5c7a407964a13dc0cf9d977207/grh-compact/src/New/Selmer/Corepresentable.lean\#L136-L138}{{\leanfont evenClasses}}}{Corepresentable.lean}{136--138}
\crow{IsNewPlaces}{\href{https://github.com/lsernau/grh-lean/blob/72fead102f663d5c7a407964a13dc0cf9d977207/grh-compact/src/New/Selmer/Corepresentable.lean\#L142-L143}{{\leanfont IsNewPlaces}}}{Corepresentable.lean}{142--143}
\crow{withConjInv}{\href{https://github.com/lsernau/grh-lean/blob/72fead102f663d5c7a407964a13dc0cf9d977207/grh-compact/src/New/Selmer/Corepresentable.lean\#L147-L147}{{\leanfont withConjInv}}}{Corepresentable.lean}{147--147}
\crow{IsPrimeTriangle}{\href{https://github.com/lsernau/grh-lean/blob/72fead102f663d5c7a407964a13dc0cf9d977207/grh-compact/src/New/Selmer/Corepresentable.lean\#L152-L162}{{\leanfont IsPrimeTriangle}}}{Corepresentable.lean}{152--162}
\crow{IsAuxTwist}{\href{https://github.com/lsernau/grh-lean/blob/72fead102f663d5c7a407964a13dc0cf9d977207/grh-compact/src/New/Selmer/Corepresentable.lean\#L166-L167}{{\leanfont IsAuxTwist}}}{Corepresentable.lean}{166--167}
\crow{fourdvdcycCharsubone}{\href{https://github.com/lsernau/grh-lean/blob/72fead102f663d5c7a407964a13dc0cf9d977207/grh-compact/src/New/Selmer/Corepresentable.lean\#L169-L173}{{\leanfont four\_dvd\_cycChar\_sub\_one}}}{Corepresentable.lean}{169--173}
\crow{cycInvwithConjInv}{\href{https://github.com/lsernau/grh-lean/blob/72fead102f663d5c7a407964a13dc0cf9d977207/grh-compact/src/New/Selmer/Corepresentable.lean\#L178-L185}{{\leanfont IsNewPlaces.\allowbreak{}cycInv\_withConjInv}}}{Corepresentable.lean}{178--185}
\crow{memselmerModuleoftameeqzero}{\href{https://github.com/lsernau/grh-lean/blob/72fead102f663d5c7a407964a13dc0cf9d977207/grh-compact/src/New/Selmer/Corepresentable.lean\#L190-L199}{{\leanfont IsNewPlaces.\allowbreak{}mem\_selmerModule\_of\_tame\_eq\_zero}}}{Corepresentable.lean}{190--199}
\crow{unionWithConj}{\href{https://github.com/lsernau/grh-lean/blob/72fead102f663d5c7a407964a13dc0cf9d977207/grh-compact/src/New/Selmer/Corepresentable.lean\#L206-L206}{{\leanfont unionWithConj}}}{Corepresentable.lean}{206--206}
\crow{pointChar}{\href{https://github.com/lsernau/grh-lean/blob/72fead102f663d5c7a407964a13dc0cf9d977207/grh-compact/src/New/Selmer/Corepresentable.lean\#L210-L211}{{\leanfont pointChar}}}{Corepresentable.lean}{210--211}
\crow{twistAtPoint}{\href{https://github.com/lsernau/grh-lean/blob/72fead102f663d5c7a407964a13dc0cf9d977207/grh-compact/src/New/Selmer/Corepresentable.lean\#L215-L216}{{\leanfont twistAtPoint}}}{Corepresentable.lean}{215--216}
\crow{reduceModPow}{\href{https://github.com/lsernau/grh-lean/blob/72fead102f663d5c7a407964a13dc0cf9d977207/grh-compact/src/New/Selmer/Corepresentable.lean\#L219-L220}{{\leanfont reduceModPow}}}{Corepresentable.lean}{219--220}
\crow{IsBalanced}{\href{https://github.com/lsernau/grh-lean/blob/72fead102f663d5c7a407964a13dc0cf9d977207/grh-compact/src/New/Selmer/Corepresentable.lean\#L227-L238}{{\leanfont IsBalanced}}}{Corepresentable.lean}{227--238}
\crow{IsPencil}{\href{https://github.com/lsernau/grh-lean/blob/72fead102f663d5c7a407964a13dc0cf9d977207/grh-compact/src/New/Selmer/Corepresentable.lean\#L243-L250}{{\leanfont IsPencil}}}{Corepresentable.lean}{243--250}
\crow{tieRelations}{\href{https://github.com/lsernau/grh-lean/blob/72fead102f663d5c7a407964a13dc0cf9d977207/grh-compact/src/New/Selmer/Corepresentable.lean\#L254-L256}{{\leanfont tieRelations}}}{Corepresentable.lean}{254--256}
\crow{idOrConjInvmeminertiaAbove}{\href{https://github.com/lsernau/grh-lean/blob/72fead102f663d5c7a407964a13dc0cf9d977207/grh-compact/src/New/Selmer/Corepresentable.lean\#L258-L260}{{\leanfont idOrConjInv\_mem\_inertiaAbove}}}{Corepresentable.lean}{258--260}
\crow{existscorepresentingcoords}{\href{https://github.com/lsernau/grh-lean/blob/72fead102f663d5c7a407964a13dc0cf9d977207/grh-compact/src/New/Selmer/Corepresentable.lean\#L267-L286}{{\leanfont exists\_corepresenting\_coords}}}{Corepresentable.lean}{267--286}
\crow{crystallinedefectle}{\href{https://github.com/lsernau/grh-lean/blob/72fead102f663d5c7a407964a13dc0cf9d977207/grh-compact/src/New/Selmer/CrystallineBound.lean\#L17-L74}{{\leanfont crystalline\_defect\_le}}}{CrystallineBound.lean}{17--74}
\crow{matrixLattice}{\href{https://github.com/lsernau/grh-lean/blob/72fead102f663d5c7a407964a13dc0cf9d977207/grh-compact/src/New/Selmer/CrystallineRank.lean\#L22-L23}{{\leanfont matrixLattice}}}{CrystallineRank.lean}{22--23}
\crow{finrankcrysCocycleseqaddnegjumps}{\href{https://github.com/lsernau/grh-lean/blob/72fead102f663d5c7a407964a13dc0cf9d977207/grh-compact/src/New/Selmer/CrystallineRank.lean\#L28-L42}{{\leanfont finrank\_crysCocycles\_eq\_add\_neg\_jumps}}}{CrystallineRank.lean}{28--42}
\crow{normLengthdefectle}{\href{https://github.com/lsernau/grh-lean/blob/72fead102f663d5c7a407964a13dc0cf9d977207/grh-compact/src/New/Selmer/CrystallineRank.lean\#L47-L56}{{\leanfont normLength\_defect\_le}}}{CrystallineRank.lean}{47--56}
\crow{localChart}{\href{https://github.com/lsernau/grh-lean/blob/72fead102f663d5c7a407964a13dc0cf9d977207/grh-compact/src/New/Selmer/DeterminantBound.lean\#L22-L26}{{\leanfont localChart}}}{DeterminantBound.lean}{22--26}
\crow{coeffEmb}{\href{https://github.com/lsernau/grh-lean/blob/72fead102f663d5c7a407964a13dc0cf9d977207/grh-compact/src/New/Selmer/DeterminantBound.lean\#L34-L35}{{\leanfont coeffEmb}}}{DeterminantBound.lean}{34--35}
\crow{ZfAt}{\href{https://github.com/lsernau/grh-lean/blob/72fead102f663d5c7a407964a13dc0cf9d977207/grh-compact/src/New/Selmer/DeterminantBound.lean\#L39-L39}{{\leanfont ZfAt}}}{DeterminantBound.lean}{39--39}
\crow{crysRels}{\href{https://github.com/lsernau/grh-lean/blob/72fead102f663d5c7a407964a13dc0cf9d977207/grh-compact/src/New/Selmer/DeterminantBound.lean\#L43-L43}{{\leanfont crysRels}}}{DeterminantBound.lean}{43--43}
\crow{blockLoss}{\href{https://github.com/lsernau/grh-lean/blob/72fead102f663d5c7a407964a13dc0cf9d977207/grh-compact/src/New/Selmer/DeterminantBound.lean\#L47-L48}{{\leanfont blockLoss}}}{DeterminantBound.lean}{47--48}
\crow{obeyingcrysRels}{\href{https://github.com/lsernau/grh-lean/blob/72fead102f663d5c7a407964a13dc0cf9d977207/grh-compact/src/New/Selmer/DeterminantBound.lean\#L50-L52}{{\leanfont obeying\_crysRels}}}{DeterminantBound.lean}{50--52}
\crow{cmSign}{\href{https://github.com/lsernau/grh-lean/blob/72fead102f663d5c7a407964a13dc0cf9d977207/grh-compact/src/New/Selmer/DeterminantBound.lean\#L56-L56}{{\leanfont cmSign}}}{DeterminantBound.lean}{56--56}
\crow{complexEmbOf}{\href{https://github.com/lsernau/grh-lean/blob/72fead102f663d5c7a407964a13dc0cf9d977207/grh-compact/src/New/Selmer/DeterminantBound.lean\#L59-L60}{{\leanfont complexEmbOf}}}{DeterminantBound.lean}{59--60}
\crow{IsRowType}{\href{https://github.com/lsernau/grh-lean/blob/72fead102f663d5c7a407964a13dc0cf9d977207/grh-compact/src/New/Selmer/DeterminantBound.lean\#L64-L70}{{\leanfont IsRowType}}}{DeterminantBound.lean}{64--70}
\crow{lambdaLoss}{\href{https://github.com/lsernau/grh-lean/blob/72fead102f663d5c7a407964a13dc0cf9d977207/grh-compact/src/New/Selmer/DeterminantBound.lean\#L74-L75}{{\leanfont lambdaLoss}}}{DeterminantBound.lean}{74--75}
\crow{loss}{\href{https://github.com/lsernau/grh-lean/blob/72fead102f663d5c7a407964a13dc0cf9d977207/grh-compact/src/New/Selmer/DeterminantBound.lean\#L78-L79}{{\leanfont loss}}}{DeterminantBound.lean}{78--79}
\crow{coeffEmbsurjective}{\href{https://github.com/lsernau/grh-lean/blob/72fead102f663d5c7a407964a13dc0cf9d977207/grh-compact/src/New/Selmer/DeterminantBound.lean\#L81-L85}{{\leanfont coeffEmb\_surjective}}}{DeterminantBound.lean}{81--85}
\crow{tateTwistsmul}{\href{https://github.com/lsernau/grh-lean/blob/72fead102f663d5c7a407964a13dc0cf9d977207/grh-compact/src/New/Selmer/DeterminantBound.lean\#L87-L89}{{\leanfont tateTwist\_smul}}}{DeterminantBound.lean}{87--89}
\crow{complexEmbOfeqiff}{\href{https://github.com/lsernau/grh-lean/blob/72fead102f663d5c7a407964a13dc0cf9d977207/grh-compact/src/New/Selmer/DeterminantBound.lean\#L91-L94}{{\leanfont IsCompletionGroup.\allowbreak{}complexEmbOf\_eq\_iff}}}{DeterminantBound.lean}{91--94}
\crow{indexeqcardrange}{\href{https://github.com/lsernau/grh-lean/blob/72fead102f663d5c7a407964a13dc0cf9d977207/grh-compact/src/New/Selmer/DeterminantBound.lean\#L98-L100}{{\leanfont IsCompletionGroup.\allowbreak{}index\_eq\_card\_range}}}{DeterminantBound.lean}{98--100}
\crow{sumcomplexEmbOfle}{\href{https://github.com/lsernau/grh-lean/blob/72fead102f663d5c7a407964a13dc0cf9d977207/grh-compact/src/New/Selmer/DeterminantBound.lean\#L102-L106}{{\leanfont sum\_complexEmbOf\_le}}}{DeterminantBound.lean}{102--106}
\crow{normgramdetleofoperators}{\href{https://github.com/lsernau/grh-lean/blob/72fead102f663d5c7a407964a13dc0cf9d977207/grh-compact/src/New/Selmer/DeterminantBound.lean\#L111-L127}{{\leanfont norm\_gram\_det\_le\_of\_operators}}}{DeterminantBound.lean}{111--127}
\crow{isAptauxPair}{\href{https://github.com/lsernau/grh-lean/blob/72fead102f663d5c7a407964a13dc0cf9d977207/grh-compact/src/New/Selmer/HeckeStableSubspaces.lean\#L22-L30}{{\leanfont isApt\_auxPair}}}{HeckeStableSubspaces.lean}{22--30}
\crow{isAptrankone}{\href{https://github.com/lsernau/grh-lean/blob/72fead102f663d5c7a407964a13dc0cf9d977207/grh-compact/src/New/Selmer/HeckeStableSubspaces.lean\#L34-L38}{{\leanfont isApt\_rank\_one}}}{HeckeStableSubspaces.lean}{34--38}
\crow{trivialLiftmemrangetupleProjrankone}{\href{https://github.com/lsernau/grh-lean/blob/72fead102f663d5c7a407964a13dc0cf9d977207/grh-compact/src/New/Selmer/HeckeStableSubspaces.lean\#L42-L54}{{\leanfont trivialLift\_mem\_range\_tupleProj\_rank\_one}}}{HeckeStableSubspaces.lean}{42--54}
\crow{isConjSelfDualofattached}{\href{https://github.com/lsernau/grh-lean/blob/72fead102f663d5c7a407964a13dc0cf9d977207/grh-compact/src/New/Selmer/HeckeStableSubspaces.lean\#L59-L77}{{\leanfont isConjSelfDual\_of\_attached}}}{HeckeStableSubspaces.lean}{59--77}
\crow{localclausesofconjPlace}{\href{https://github.com/lsernau/grh-lean/blob/72fead102f663d5c7a407964a13dc0cf9d977207/grh-compact/src/New/Selmer/HeckeStableSubspaces.lean\#L83-L110}{{\leanfont local\_clauses\_of\_conjPlace}}}{HeckeStableSubspaces.lean}{83--110}
\crow{orthogonalmeminvtSubmodulehecke}{\href{https://github.com/lsernau/grh-lean/blob/72fead102f663d5c7a407964a13dc0cf9d977207/grh-compact/src/New/Selmer/HeckeStableSubspaces.lean\#L112-L115}{{\leanfont orthogonal\_mem\_invtSubmodule\_hecke}}}{HeckeStableSubspaces.lean}{112--115}
\crow{leofisotypicmem}{\href{https://github.com/lsernau/grh-lean/blob/72fead102f663d5c7a407964a13dc0cf9d977207/grh-compact/src/New/Selmer/HeckeStableSubspaces.lean\#L119-L126}{{\leanfont le\_of\_isotypic\_mem}}}{HeckeStableSubspaces.lean}{119--126}
\crow{existsintegralmultipleevalone}{\href{https://github.com/lsernau/grh-lean/blob/72fead102f663d5c7a407964a13dc0cf9d977207/grh-compact/src/New/Selmer/HeckeString.lean\#L15-L26}{{\leanfont exists\_integral\_multiple\_eval\_one}}}{HeckeString.lean}{15--26}
\crow{existseventuallynoeigenvalueone}{\href{https://github.com/lsernau/grh-lean/blob/72fead102f663d5c7a407964a13dc0cf9d977207/grh-compact/src/New/Selmer/HeckeString.lean\#L30-L42}{{\leanfont exists\_eventually\_no\_eigenvalue\_one}}}{HeckeString.lean}{30--42}
\crow{isRootheckestring}{\href{https://github.com/lsernau/grh-lean/blob/72fead102f663d5c7a407964a13dc0cf9d977207/grh-compact/src/New/Selmer/HeckeString.lean\#L52-L62}{{\leanfont isRoot\_hecke\_string}}}{HeckeString.lean}{52--62}
\crow{eventuallyheckePolyeqstringmul}{\href{https://github.com/lsernau/grh-lean/blob/72fead102f663d5c7a407964a13dc0cf9d977207/grh-compact/src/New/Selmer/HeckeString.lean\#L68-L87}{{\leanfont eventually\_heckePoly\_eq\_string\_mul}}}{HeckeString.lean}{68--87}
\crow{normpowsubpoweqone}{\href{https://github.com/lsernau/grh-lean/blob/72fead102f663d5c7a407964a13dc0cf9d977207/grh-compact/src/New/Selmer/LabelledColumns.lean\#L25-L32}{{\leanfont norm\_pow\_sub\_pow\_eq\_one}}}{LabelledColumns.lean}{25--32}
\crow{existsinertiaeigenvalues}{\href{https://github.com/lsernau/grh-lean/blob/72fead102f663d5c7a407964a13dc0cf9d977207/grh-compact/src/New/Selmer/LabelledColumns.lean\#L40-L62}{{\leanfont exists\_inertia\_eigenvalues}}}{LabelledColumns.lean}{40--62}
\crow{isRow}{\href{https://github.com/lsernau/grh-lean/blob/72fead102f663d5c7a407964a13dc0cf9d977207/grh-compact/src/New/Selmer/LabelledColumns.lean\#L67-L125}{{\leanfont IsLabelledColumn.\allowbreak{}isRow}}}{LabelledColumns.lean}{67--125}
\crow{normalisedRepquadraticoninertia}{\href{https://github.com/lsernau/grh-lean/blob/72fead102f663d5c7a407964a13dc0cf9d977207/grh-compact/src/New/Selmer/LabelledColumns.lean\#L130-L155}{{\leanfont normalisedRep\_quadratic\_on\_inertia}}}{LabelledColumns.lean}{130--155}
\crow{twistavatar}{\href{https://github.com/lsernau/grh-lean/blob/72fead102f663d5c7a407964a13dc0cf9d977207/grh-compact/src/New/Selmer/LabelledColumns.lean\#L167-L182}{{\leanfont IsLabelledColumn.\allowbreak{}twist\_avatar}}}{LabelledColumns.lean}{167--182}
\crow{isLabelledColumnofhasProductLabel}{\href{https://github.com/lsernau/grh-lean/blob/72fead102f663d5c7a407964a13dc0cf9d977207/grh-compact/src/New/Selmer/LabelledColumns.lean\#L187-L197}{{\leanfont isLabelledColumn\_of\_hasProductLabel}}}{LabelledColumns.lean}{187--197}
\crow{algebraMapmemcycloField}{\href{https://github.com/lsernau/grh-lean/blob/72fead102f663d5c7a407964a13dc0cf9d977207/grh-compact/src/New/Selmer/LevelPlacesChosen.lean\#L20-L21}{{\leanfont Data.\allowbreak{}algebraMap\_mem\_cycloField}}}{LevelPlacesChosen.lean}{20--21}
\crow{AbsGal}{\href{https://github.com/lsernau/grh-lean/blob/72fead102f663d5c7a407964a13dc0cf9d977207/grh-compact/src/New/Selmer/LevelPlacesChosen.lean\#L25-L25}{{\leanfont AbsGal}}}{LevelPlacesChosen.lean}{25--25}
\crow{placesAboveG}{\href{https://github.com/lsernau/grh-lean/blob/72fead102f663d5c7a407964a13dc0cf9d977207/grh-compact/src/New/Selmer/LevelPlacesChosen.lean\#L28-L28}{{\leanfont placesAboveG}}}{LevelPlacesChosen.lean}{28--28}
\crow{badPlaces}{\href{https://github.com/lsernau/grh-lean/blob/72fead102f663d5c7a407964a13dc0cf9d977207/grh-compact/src/New/Selmer/LevelPlacesChosen.lean\#L32-L32}{{\leanfont badPlaces}}}{LevelPlacesChosen.lean}{32--32}
\crow{placesDividingunionsubsetbadPlaces}{\href{https://github.com/lsernau/grh-lean/blob/72fead102f663d5c7a407964a13dc0cf9d977207/grh-compact/src/New/Selmer/LevelPlacesChosen.lean\#L34-L36}{{\leanfont placesDividing\_union\_subset\_badPlaces}}}{LevelPlacesChosen.lean}{34--36}
\crow{starmemplacesAboveG}{\href{https://github.com/lsernau/grh-lean/blob/72fead102f663d5c7a407964a13dc0cf9d977207/grh-compact/src/New/Selmer/LevelPlacesChosen.lean\#L40-L41}{{\leanfont star\_mem\_placesAboveG}}}{LevelPlacesChosen.lean}{40--41}
\crow{normnatCastlenormsubofpoweqone}{\href{https://github.com/lsernau/grh-lean/blob/72fead102f663d5c7a407964a13dc0cf9d977207/grh-compact/src/New/Selmer/LevelPlacesChosen.lean\#L44-L46}{{\leanfont norm\_natCast\_le\_norm\_sub\_of\_pow\_eq\_one}}}{LevelPlacesChosen.lean}{44--46}
\crow{continuouscycInv}{\href{https://github.com/lsernau/grh-lean/blob/72fead102f663d5c7a407964a13dc0cf9d977207/grh-compact/src/New/Selmer/LevelPlacesChosen.lean\#L48-L49}{{\leanfont continuous\_cycInv}}}{LevelPlacesChosen.lean}{48--49}
\crow{cycInveqofisFrob}{\href{https://github.com/lsernau/grh-lean/blob/72fead102f663d5c7a407964a13dc0cf9d977207/grh-compact/src/New/Selmer/LevelPlacesChosen.lean\#L51-L53}{{\leanfont cycInv\_eq\_of\_isFrob}}}{LevelPlacesChosen.lean}{51--53}
\crow{invAvatar}{\href{https://github.com/lsernau/grh-lean/blob/72fead102f663d5c7a407964a13dc0cf9d977207/grh-compact/src/New/Selmer/LevelPlacesChosen.lean\#L59-L60}{{\leanfont Data.\allowbreak{}TameChar.\allowbreak{}invAvatar}}}{LevelPlacesChosen.lean}{59--60}
\crow{avatarconjGalmul}{\href{https://github.com/lsernau/grh-lean/blob/72fead102f663d5c7a407964a13dc0cf9d977207/grh-compact/src/New/Selmer/LevelPlacesChosen.lean\#L64-L73}{{\leanfont avatar\_conjGal\_mul}}}{LevelPlacesChosen.lean}{64--73}
\crow{isLabelledAtavatar}{\href{https://github.com/lsernau/grh-lean/blob/72fead102f663d5c7a407964a13dc0cf9d977207/grh-compact/src/New/Selmer/LevelPlacesChosen.lean\#L78-L84}{{\leanfont isLabelledAt\_avatar}}}{LevelPlacesChosen.lean}{78--84}
\crow{isChosenchosenPlace}{\href{https://github.com/lsernau/grh-lean/blob/72fead102f663d5c7a407964a13dc0cf9d977207/grh-compact/src/New/Selmer/LevelPlacesChosen.lean\#L86-L89}{{\leanfont isChosen\_chosenPlace}}}{LevelPlacesChosen.lean}{86--89}
\crow{blockdefectleatplace}{\href{https://github.com/lsernau/grh-lean/blob/72fead102f663d5c7a407964a13dc0cf9d977207/grh-compact/src/New/Selmer/OnePlaceDefect.lean\#L15-L84}{{\leanfont block\_defect\_le\_at\_place}}}{OnePlaceDefect.lean}{15--84}
\crow{fixesemboffixesrootsOfUnity}{\href{https://github.com/lsernau/grh-lean/blob/72fead102f663d5c7a407964a13dc0cf9d977207/grh-compact/src/New/Selmer/ThickSpacing.lean\#L12-L16}{{\leanfont fixes\_emb\_of\_fixes\_rootsOfUnity}}}{ThickSpacing.lean}{12--16}
\crow{cmSignspec}{\href{https://github.com/lsernau/grh-lean/blob/72fead102f663d5c7a407964a13dc0cf9d977207/grh-compact/src/New/Selmer/ThickSpacing.lean\#L18-L20}{{\leanfont cmSign\_spec}}}{ThickSpacing.lean}{18--20}
\crow{placeClassratio}{\href{https://github.com/lsernau/grh-lean/blob/72fead102f663d5c7a407964a13dc0cf9d977207/grh-compact/src/New/Selmer/ThickSpacing.lean\#L30-L39}{{\leanfont placeClass\_ratio}}}{ThickSpacing.lean}{30--39}
\crow{notthickofshortarc}{\href{https://github.com/lsernau/grh-lean/blob/72fead102f663d5c7a407964a13dc0cf9d977207/grh-compact/src/New/Selmer/ThickSpacing.lean\#L45-L57}{{\leanfont not\_thick\_of\_short\_arc}}}{ThickSpacing.lean}{45--57}
\crow{cmSigneqofthick}{\href{https://github.com/lsernau/grh-lean/blob/72fead102f663d5c7a407964a13dc0cf9d977207/grh-compact/src/New/Selmer/ThickSpacing.lean\#L64-L75}{{\leanfont cmSign\_eq\_of\_thick}}}{ThickSpacing.lean}{64--75}
\crow{existsisIntegralMeasure}{\href{https://github.com/lsernau/grh-lean/blob/72fead102f663d5c7a407964a13dc0cf9d977207/grh-compact/src/New/Selmer/ValueFamilies.lean\#L30-L80}{{\leanfont Data.\allowbreak{}Pair.\allowbreak{}exists\_isIntegralMeasure}}}{ValueFamilies.lean}{30--80}
\crow{coprimeqofisAuxiliary}{\href{https://github.com/lsernau/grh-lean/blob/72fead102f663d5c7a407964a13dc0cf9d977207/grh-compact/src/New/Selmer/ValueFamilies.lean\#L82-L85}{{\leanfont coprime\_q\_of\_isAuxiliary}}}{ValueFamilies.lean}{82--85}
\crow{tameCharOfAuxiliary}{\href{https://github.com/lsernau/grh-lean/blob/72fead102f663d5c7a407964a13dc0cf9d977207/grh-compact/src/New/Selmer/ValueFamilies.lean\#L89-L90}{{\leanfont tameCharOfAuxiliary}}}{ValueFamilies.lean}{89--90}
\crow{invAvatareqofartinMapeq}{\href{https://github.com/lsernau/grh-lean/blob/72fead102f663d5c7a407964a13dc0cf9d977207/grh-compact/src/New/Selmer/ValueFamilies.lean\#L92-L95}{{\leanfont invAvatar\_eq\_of\_artinMap\_eq}}}{ValueFamilies.lean}{92--95}
\crow{existstameCharinvAvatareq}{\href{https://github.com/lsernau/grh-lean/blob/72fead102f663d5c7a407964a13dc0cf9d977207/grh-compact/src/New/Selmer/ValueFamilies.lean\#L97-L106}{{\leanfont exists\_tameChar\_invAvatar\_eq}}}{ValueFamilies.lean}{97--106}
\crow{existsisFamily}{\href{https://github.com/lsernau/grh-lean/blob/72fead102f663d5c7a407964a13dc0cf9d977207/grh-compact/src/New/Selmer/ValueFamilies.lean\#L111-L153}{{\leanfont exists\_isFamily}}}{ValueFamilies.lean}{111--153}}
% the notations of the defined functions: one slot for each explicit argument of the Lean declaration
\newcommand{\IsConjSelfDual}[2]{\mathsf{CSD}_{#1}(#2)}
\newcommand{\crudeUnit}[2]{\mathrm{cr}_{#1}(#2)}
\newcommand{\SelmerLengthLe}[6]{\mathsf{Short}(#1,#2;\,#3,#4;\,#5,#6)}
\newcommand{\IsColumn}[5]{\mathsf{ColRep}_{#1}(#2,#3,#4;\,#5)}
\newcommand{\IsRow}[4]{\mathsf{RowRep}_{#1}(#2,#3;\,#4)}
\newcommand{\chiOneInv}[1]{\chi_{1,#1}^{-1}}
\newcommand{\normalisedRep}[2]{\operatorname{nm}_{#1}(#2)}
\newcommand{\twistedLabel}[3]{\operatorname{tl}_{#1}(#2,#3)}
\newcommand{\IsApt}[3]{\mathsf{Apt}(#1,#2;\,#3)}
\newcommand{\tupleProj}[2]{\operatorname{pr}_{#1}(#2)}
\newcommand{\HasRow}[3]{\mathsf{HasRow}(#1,#2;\,#3)}
\newcommand{\IsLabelledColumn}[6]{\mathsf{LabCol}(#1;\,#2,#3,#5;\,#4;\,#6)}
\newcommand{\obeyingCocycles}[3]{\mathcal{Z}_{#1}(#2;\,#3)}
\newcommand{\RowBlockIndex}[2]{\mathrm{I}_{#1}(#2)}
\newcommand{\blockRep}[3]{\operatorname{blk}_{#1}(#2)(#3)}
\newcommand{\AuxPrimeData}[2]{\mathrm{Aux}(#1,#2)}
\newcommand{\IsAnticyclotomicAt}[3]{\mathsf{Acyc}(#1,#2;\,#3)}
\newcommand{\IsAuxiliary}[2]{\mathsf{Auxil}(#1;\,#2)}
\newcommand{\IsFamily}[3]{\mathsf{Fam}(#1;\,#2,#3)}
\newcommand{\evenClasses}[4]{\mathrm{Ev}(#1,#2;\,#3,#4)}
\newcommand{\IsNewPlaces}[4]{\mathsf{New}_{#1}(#2,#3,#4)}
\newcommand{\withConjInv}[2]{#1^{[#2]}}
\newcommand{\IsPrimeTriangle}[9]{\mathsf{Tri}_{#1}(#2,#3,#4;\,#5;\,#6,#7;\,#8,#9)}
\newcommand{\IsAuxTwist}[4]{\mathsf{AuxTw}(#1,#2;\,#3,#4)}
\newcommand{\unionWithConj}[2]{#1\cup #2\cup\overline{#2}}
\newcommand{\twistAtPoint}[5]{\operatorname{tp}_{#1,#2}(#3,#4)(#5)}
\newcommand{\reduceModPow}[3]{#3\bmod #1^{#2}}
\newcommand{\IsBalanced}[7]{\mathsf{Bal}_{#1}(#2;\,#3,#4;\,#5,#6,#7)}
\newcommand{\IsPencil}[8]{\mathsf{Pencil}_{#1}(#2;\,#3,#4;\,#5,#6,#7;\,#8)}
\newcommand{\tieRelations}[3]{\operatorname{Ties}_{#3}(#1,#2)}
\newcommand{\matrixLattice}[1]{\mathrm{L}(#1)}
\newcommand{\coeffEmb}[2]{\operatorname{emb}_{#1}(#2)}
\newcommand{\ZfAt}[2]{Z_{\mathrm{cris}}(#1;\,#2)}
\newcommand{\crysRels}[1]{\mathcal{R}_{\mathrm{cris}}(#1)}
\newcommand{\blockLoss}[3]{\operatorname{bl}(#1;\,#2,#3)}
\newcommand{\cmSign}[1]{\operatorname{sgn}_{\mathrm{CM}}(#1)}
\newcommand{\complexEmbOf}[2]{\operatorname{ce}(#1,#2)}
\newcommand{\IsRowType}[4]{\mathsf{RowType}_{#1}(#2,#3;\,#4)}
\newcommand{\lambdaLoss}[3]{\operatorname{lam}(#1,#2;\,#3)}
\newcommand{\loss}[2]{\operatorname{loss}_{#1}(#2)}
\newcommand{\AbsGal}[1]{G_{E(#1)}}
\newcommand{\placesAboveG}[1]{\Sigma_{g}(#1)}
\newcommand{\badPlaces}[2]{\Sigma_{\mathrm{bad}}(#1,#2)}
\newcommand{\invAvatar}[2]{\operatorname{av}^{-}_{#1}(#2)}
\newcommand{\UnitMembers}{205}
\newcommand{\UnitInNote}{116}
\newcommand{\UnitProofs}{96}
\newcommand{\UnitSteps}{975}
\newcommand{\UnitAllSteps}{1091}


\linespread{1.08}                       % 11 pt on 14.7 pt
\setlength{\parindent}{0pt}
\setlength{\parskip}{5.5pt plus 1.5pt minus 0.5pt}
\setlength{\abovedisplayskip}{7.5pt plus 2pt minus 3pt}
\setlength{\belowdisplayskip}{7.5pt plus 2pt minus 3pt}
\setlength{\abovedisplayshortskip}{3pt plus 2pt}
\setlength{\belowdisplayshortskip}{6pt plus 2pt minus 2pt}
\emergencystretch=1.5em
\clubpenalty=10000 \widowpenalty=10000 \displaywidowpenalty=10000
\relpenalty=7000 \binoppenalty=9000
\raggedbottom
\pagestyle{fancy}\fancyhf{}\renewcommand{\headrulewidth}{0pt}
\fancyfoot[C]{\sffamily\fontsize{7.6}{9}\selectfont\color{inkfaint}\thepage\,/\,\pageref*{LastPage}}
\setlist{topsep=2pt,partopsep=0pt,parsep=0pt,itemsep=4pt,leftmargin=2.1em,labelsep=0.55em}

\newcommand{\Q}{\mathbb{Q}}
\newcommand{\Z}{\mathbb{Z}}
\newcommand{\F}{\mathbb{F}}
\newcommand{\len}{\operatorname{len}}
\newcommand{\cok}{\operatorname{cok}}
\newcommand{\End}{\operatorname{End}}
\newcommand{\Hom}{\operatorname{Hom}}
\newcommand{\Gal}{\operatorname{Gal}}
\newcommand{\diag}{\operatorname{diag}}
\newcommand{\kw}[1]{\emph{#1}}
\newcommand{\m}[1]{\mbox{$#1$}}        % a short formula that must not be broken across lines
\newcommand{\head}[1]{\par\addvspace{17pt plus 4pt minus 2pt}\Needspace{6\baselineskip}%
  \noindent{\fontsize{13.2}{16}\selectfont\bfseries\color{ink}#1}\par\nopagebreak\vspace{-1pt}}
\newcommand{\sect}[2]{\refstepcounter{section}\label{#2}\par\addvspace{19pt plus 4pt minus 2pt}\Needspace{7\baselineskip}%
  \noindent\llap{{\fontsize{13.2}{16}\selectfont\bfseries\color{inkfaint}\thesection}\hspace{9pt}}%
  {\fontsize{13.2}{16}\selectfont\bfseries\color{ink}#1}\par\nopagebreak\vspace{-1pt}}
\newcommand{\subhead}[1]{\par\addvspace{9pt plus 2pt}\Needspace{4\baselineskip}\noindent\textbf{#1}\enspace\ignorespaces}
% where a statement lives in the Lean archive
\newcommand{\leantag}{{\sffamily\bfseries\fontsize{6.3}{8}\selectfont\color{leancol}\addfontfeatures{LetterSpace=7}LEAN}}
\newcommand{\sep}{\unskip\nobreak\hspace{5pt}{\color{inkfaint}\textperiodcentered}\hspace{5pt plus 2pt}\ignorespaces}
\newcommand{\leanline}[1]{\par\nopagebreak\vspace{-2.5pt}{\fontsize{8.9}{12.4}\selectfont\color{inksoft}\raggedright
  \noindent\llap{\leantag\hspace{9pt}}#1\par}\vspace{1pt}}
% a displayed statement in a tinted box: #1 colour, #2 text
\newcommand{\statement}[2]{\par\addvspace{7pt}\begingroup\setlength{\fboxsep}{8.5pt}\setlength{\fboxrule}{0.6pt}%
  \noindent\fcolorbox{#1!30!white}{#1!3!white}{\parbox{\dimexpr\linewidth-2\fboxsep-2\fboxrule\relax}{#2}}%
  \endgroup\par\addvspace{7pt}}
\newcommand{\nt}[2]{\par\noindent\hangindent=1.1em\hangafter=1 #1\enspace{\color{inksoft}#2}\par\vspace{1pt}}
\newcommand{\legdot}[1]{\tikz[baseline=-0.62ex]\fill[#1] (0,0) circle (2.5pt);\hspace{3pt}}
\newcommand{\crow}[4]{\csname CS@#1\endcsname & #2 & \textcolor{inksoft}{#3} & \textcolor{inksoft}{#4}\\}
\newcommand{\leanquoteline}[2]{\par\noindent\hangindent=\dimexpr#1\fontdimen2\font+2.2em\relax\hangafter=1
  \hspace*{#1\fontdimen2\font}#2}
\newcommand{\quoted}[1]{\par\vspace{3pt}{\ttfamily\fontsize{7.5}{10.4}\selectfont\color{leancol}\raggedright
  \setlength{\parskip}{0pt}\leanquote{#1}\par}\vspace{4pt}}
\newcommand{\bx}[3]{{\sffamily\bfseries\fontsize{7.3}{9.2}\selectfont\color{#1!82!black}#2}\\[1.5pt]#3}
\newcommand{\cs}[2]{\expandafter\newcommand\csname CS@#1\endcsname{#2}}

% theorems
\newtheoremstyle{plainx}{7pt plus 2pt}{7pt plus 2pt}{\itshape}{}{\bfseries}{.}{0.5em}{}
\newtheoremstyle{uprightx}{7pt plus 2pt}{7pt plus 2pt}{}{}{\bfseries}{.}{0.5em}{}
\theoremstyle{plainx}
\newtheorem{theorem}{Theorem}[section]
\newtheorem{lemma}[theorem]{Lemma}
\newtheorem{proposition}[theorem]{Proposition}
\newtheorem{corollary}[theorem]{Corollary}
\theoremstyle{uprightx}
\newtheorem{definition}[theorem]{Definition}
\newtheorem{setting}[theorem]{Setting}          % hypotheses shared by several statements; each of them cites it by number
\newtheorem{remark}[theorem]{Remark}
\newtheorem{example}[theorem]{Example}
\renewcommand{\qedsymbol}{{\setlength{\fboxsep}{0pt}\setlength{\fboxrule}{0.5pt}\fbox{\rule{0pt}{0.46em}\hspace{0.46em}}}}
% what holds throughout the note: the ambient objects only (implicit and instance arguments), never a hypothesis
\newenvironment{conventions}{\par\addvspace{5pt}\noindent\textbf{Conventions.}\enspace\ignorespaces}{\par\addvspace{3pt}}
\makeatletter
\renewenvironment{proof}[1][Proof]{\par\pushQED{\qed}\normalfont\topsep0pt\partopsep0pt
  \trivlist\item[\hskip\labelsep\itshape #1.]\ignorespaces}{\popQED\endtrivlist\@endpefalse}
\makeatother

% ---- macros of the three notes of this unit (background notions; the notions that the unit defines are in data.tex)
\newcommand{\CC}{\mathbb{C}}
\newcommand{\NN}{\mathbb{N}}
\newcommand{\Qb}{\bar{\mathbb{Q}}}
\newcommand{\Ql}{\mathbb{Q}_\ell}
\newcommand{\Zl}{\mathbb{Z}_\ell}
\newcommand{\Qlb}{\bar{\mathbb{Q}}_\ell}
\newcommand{\Zlb}{\bar{\mathbb{Z}}_\ell}
\newcommand{\Cl}{\mathbb{C}_\ell}
\newcommand{\Gl}{G_{\mathbb{Q}_\ell}}
\newcommand{\GE}[1]{G_{#1}}
\newcommand{\cyc}{\operatorname{cyc}}
\newcommand{\vl}{v_\ell}
\newcommand{\Sel}[4]{\operatorname{Sel}_{#2}^{#3}(#1;\,#4)}
\newcommand{\Ob}[2]{\operatorname{Ob}(#1;\,#2)}
\newcommand{\Coc}[1]{\operatorname{Coc}(#1)}
\newcommand{\Rel}[1]{\operatorname{Rel}(#1)}
\newcommand{\pd}[2]{\Sigma_{#2}(#1)}
\newcommand{\inab}[1]{I_{#1}}
\newcommand{\cj}[1]{\bar{#1}}
\newcommand{\Nm}{\operatorname{N}}
\newcommand{\nlen}{\operatorname{nlen}}
\newcommand{\cO}{\mathcal{O}}
\newcommand{\fP}{\mathfrak{P}}
\newcommand{\fq}{\mathfrak{q}}
\newcommand{\HT}{\mathcal{H}}
\newcommand{\GL}{\mathrm{GL}}
\newcommand{\tr}{\operatorname{tr}}
\newcommand{\rk}{\operatorname{rk}}
\newcommand{\Art}{\operatorname{Art}}
\newcommand{\av}{\operatorname{av}}
\newcommand{\Zf}{Z_{\mathrm{f}}}
% the data (as in the note on admissible data)
\newcommand{\Dm}[1]{m(#1)}
\newcommand{\Dr}[1]{r(#1)}
\newcommand{\kD}[1]{k(#1)}
\newcommand{\tD}[1]{t(#1)}
\newcommand{\fD}[1]{f(#1)}
\newcommand{\ED}[1]{E(#1)}
\newcommand{\Dg}[2]{g_{#1}(#2)}
\newcommand{\gD}[1]{g(#1)}
\newcommand{\lvl}[1]{N(#1)}
\newcommand{\hhD}[1]{h_2(#1)}
\newcommand{\hD}[1]{h(#1)}
\newcommand{\bone}[1]{b_1(#1)}
\newcommand{\btwo}[1]{b_2(#1)}
\newcommand{\wt}[3]{w_{#1,#2}(#3)}
\newcommand{\ordch}[1]{\operatorname{ord}_{\mathrm{ch}}(#1)}
\newcommand{\Dmu}[1]{\mu(#1)}
\newcommand{\ND}[1]{N_1(#1)}
\newcommand{\thick}[1]{\operatorname{thick}(#1)}
\newcommand{\inp}[2]{\statement{ana}{\textbf{#1.}\enspace #2}}

% where each declaration of the concordance is treated
\cs{IsConjSelfDual}{Definition~\ref{d:csd}}
\cs{crudeUnit}{Definition~\ref{d:cr}}
\cs{SelmerLengthLe}{Definition~\ref{d:short}}
\cs{IsColumn}{Definition~\ref{d:column}}
\cs{IsRow}{Definition~\ref{d:row}}
\cs{chiOneInv}{Definition~\ref{d:chione}}
\cs{chiOneInvspec}{Lemma~\ref{l:chione}}
\cs{normalisedRep}{Definition~\ref{d:chione}}
\cs{twistedLabel}{Definition~\ref{d:twl}}
\cs{IsApt}{Definition~\ref{d:apt}}
\cs{tupleProj}{Definition~\ref{d:tproj}}
\cs{HasRow}{Definition~\ref{d:hasrow}}
\cs{IsLabelledColumn}{Definition~\ref{d:labcol}}
\cs{gpowsubonedvdchosenGroupOrder}{Lemma~\ref{l:gpow}}
\cs{ellVallecrudeUnit}{Lemma~\ref{l:small3}}
\cs{cardleofsubsetplacesAboveG}{Lemma~\ref{l:gpow}}
\cs{existschosennotMemplacesDividing}{Lemma~\ref{l:gpow}}
\cs{existsalgHom}{Lemma~\ref{l:alghom}}
\cs{continuous}{Lemma~\ref{l:small3}}
\cs{eqone}{Lemma~\ref{l:small3}}
\cs{traceeq}{Lemma~\ref{l:small3}}
\cs{cardle}{Lemma~\ref{l:small3}}
\cs{isSymmetrictupleProj}{Lemma~\ref{l:small3}}
\cs{tupleProjeqselfofmemrange}{Lemma~\ref{l:small3}}
\cs{rangetupleProjmeminvtSubmodule}{Lemma~\ref{l:tprojinv}}
\cs{frequentlyisSplitnotMemplaces}{Lemma~\ref{l:small3}}
\cs{existsbaseLevel}{Theorem C}
\cs{obeyingCocycles}{Definition~\ref{d:obey}}
\cs{existsgraphsubrepresentation}{Lemma~\ref{l:graph}}
\cs{RowBlockIndex}{Definition~\ref{d:blocks}}
\cs{blockRep}{Definition~\ref{d:blocks}}
\cs{AuxPrimeData}{Definition~\ref{d:auxdata}}
\cs{IsAnticyclotomicAt}{Definition~\ref{d:acyc}}
\cs{IsAuxiliary}{Definition~\ref{d:auxil}}
\cs{IsFamily}{Definition~\ref{d:family}}
\cs{existsisAnticyclotomicAt}{Theorem A}
\cs{evenClasses}{Definition~\ref{d:even}}
\cs{IsNewPlaces}{Definition~\ref{d:new}}
\cs{withConjInv}{Definition~\ref{d:new}}
\cs{IsPrimeTriangle}{Definition~\ref{d:tri}}
\cs{IsAuxTwist}{Definition~\ref{d:auxtw}}
\cs{fourdvdcycCharsubone}{Lemma~\ref{l:four}}
\cs{cycInvwithConjInv}{Lemma~\ref{l:newcyc}}
\cs{memselmerModuleoftameeqzero}{Lemma~\ref{l:newsel}}
\cs{unionWithConj}{Definition~\ref{d:union}}
\cs{pointChar}{Definition~\ref{d:point}}
\cs{twistAtPoint}{Definition~\ref{d:point}}
\cs{reduceModPow}{Definition~\ref{d:point}}
\cs{IsBalanced}{Definition~\ref{d:bal}}
\cs{IsPencil}{Definition~\ref{d:pencil}}
\cs{tieRelations}{Definition~\ref{d:ties}}
\cs{idOrConjInvmeminertiaAbove}{Lemma~\ref{l:small3}}
\cs{existscorepresentingcoords}{Lemma~\ref{l:corep}}
\cs{crystallinedefectle}{Lemma~\ref{l:crys}}
\cs{matrixLattice}{Definition~\ref{d:matlat}}
\cs{finrankcrysCocycleseqaddnegjumps}{Lemma~\ref{l:rank}}
\cs{normLengthdefectle}{Lemma~\ref{l:defect}}
\cs{localChart}{Definition~\ref{d:chart}}
\cs{coeffEmb}{Definition~\ref{d:coeffemb}}
\cs{ZfAt}{Definition~\ref{d:zf}}
\cs{crysRels}{Definition~\ref{d:zf}}
\cs{blockLoss}{Definition~\ref{d:blockloss}}
\cs{obeyingcrysRels}{Lemma~\ref{l:small1}}
\cs{cmSign}{Definition~\ref{d:cmsign}}
\cs{complexEmbOf}{Definition~\ref{d:coeffemb}}
\cs{IsRowType}{Definition~\ref{d:rowtype}}
\cs{lambdaLoss}{Definition~\ref{d:loss}}
\cs{loss}{Definition~\ref{d:loss}}
\cs{coeffEmbsurjective}{Lemma~\ref{l:coeffsurj}}
\cs{tateTwistsmul}{Lemma~\ref{l:small1}}
\cs{complexEmbOfeqiff}{Lemma~\ref{l:small1}}
\cs{indexeqcardrange}{Lemma~\ref{l:small1}}
\cs{sumcomplexEmbOfle}{Lemma~\ref{l:sumemb}}
\cs{normgramdetleofoperators}{Lemma~\ref{l:gram}}
\cs{isAptauxPair}{Lemma~\ref{l:aptaux}}
\cs{isAptrankone}{Lemma~\ref{l:aptaux}}
\cs{trivialLiftmemrangetupleProjrankone}{Lemma~\ref{l:aptaux}}
\cs{isConjSelfDualofattached}{Lemma~\ref{l:csdatt}}
\cs{localclausesofconjPlace}{Lemma~\ref{l:conjplace}}
\cs{orthogonalmeminvtSubmodulehecke}{Lemma~\ref{l:small3}}
\cs{leofisotypicmem}{Lemma~\ref{l:isotypic}}
\cs{existsintegralmultipleevalone}{Lemma~\ref{l:intmult}}
\cs{existseventuallynoeigenvalueone}{Lemma~\ref{l:noeig}}
\cs{isRootheckestring}{Lemma~\ref{l:string}}
\cs{eventuallyheckePolyeqstringmul}{Lemma~\ref{l:heckestring}}
\cs{normpowsubpoweqone}{Lemma~\ref{l:rootsunit}}
\cs{existsinertiaeigenvalues}{Lemma~\ref{l:inertia}}
\cs{isRow}{Theorem B}
\cs{normalisedRepquadraticoninertia}{Lemma~\ref{l:quad}}
\cs{twistavatar}{Lemma~\ref{l:twistav}}
\cs{isLabelledColumnofhasProductLabel}{Lemma~\ref{l:prodlabel}}
\newcommand{\Mom}[2]{M_\omega(#1,#2)}           % as in the note on admissible data
\newcommand{\tupleIdem}[2]{e^{\otimes}_{#1,#2}}   % as in the note on primitive labels
\cs{algebraMapmemcycloField}{Lemma~\ref{l:small2}}
\cs{AbsGal}{Definition~\ref{d:places}}
\cs{placesAboveG}{Definition~\ref{d:places}}
\cs{badPlaces}{Definition~\ref{d:places}}
\cs{placesDividingunionsubsetbadPlaces}{Lemma~\ref{l:small2}}
\cs{starmemplacesAboveG}{Lemma~\ref{l:small2}}
\cs{normnatCastlenormsubofpoweqone}{Lemma~\ref{l:small2}}
\cs{continuouscycInv}{Lemma~\ref{l:small2}}
\cs{cycInveqofisFrob}{Lemma~\ref{l:small2}}
\cs{invAvatar}{Definition~\ref{d:invav}}
\cs{avatarconjGalmul}{Lemma~\ref{l:avconj}}
\cs{isLabelledAtavatar}{Lemma~\ref{l:avlabel}}
\cs{isChosenchosenPlace}{Lemma~\ref{l:small2}}
\cs{blockdefectleatplace}{Theorem E}
\cs{fixesemboffixesrootsOfUnity}{Lemma~\ref{l:fixemb}}
\cs{cmSignspec}{Lemma~\ref{l:small3}}
\cs{placeClassratio}{Lemma~\ref{l:ratio}}
\cs{notthickofshortarc}{Lemma~\ref{l:spacing}}
\cs{cmSigneqofthick}{Lemma~\ref{l:cmthick}}
\cs{existsisIntegralMeasure}{Lemma~\ref{l:measure}}
\cs{coprimeqofisAuxiliary}{Lemma~\ref{l:coprime}}
\cs{tameCharOfAuxiliary}{Definition~\ref{d:tameaux}}
\cs{invAvatareqofartinMapeq}{Lemma~\ref{l:small3}}
\cs{existstameCharinvAvatareq}{Lemma~\ref{l:tamechar}}
\cs{existsisFamily}{Theorem D}

\begin{document}
\color{ink}

{\raggedright
{\sffamily\bfseries\fontsize{7.2}{9}\selectfont\color{leancol}\addfontfeatures{LetterSpace=8}READBACK NOTE 31\par}
\vspace{5pt}
{\fontsize{21.5}{25.5}\selectfont\bfseries Columns and rows as Galois representations: labels at the chosen places, anticyclotomic characters at new primes, families of L-values and the crystalline loss at one place\par}
\vspace{5pt}
{\fontsize{12.4}{16}\selectfont\itshape\color{inksoft} One unit of the Lean artifact \ArtifactName, read back into
mathematics\par}}
\vspace{8pt}
{\color{rulegrey}\hrule height 0.6pt}
\vspace{7pt}
\thispagestyle{empty}

\textbf{Summary.} The twists of Chapter~VIII are chosen on the Galois side. For admissible data $D$ and a prime~$\ell$
this note sets up that side. It defines the $\ell$-adic representations that play the part of columns and rows
(conjugate self-dual, with prescribed Hodge--Tate jumps, labelled at the chosen primes by the roots of polynomials
over the residue fields), the bound on Selmer modules that a good twist must achieve, the systems of auxiliary
primes (anticyclotomic data, pencils, ties, prime triangles), and the loss that the crystalline condition at~$\ell$
costs. Five results are proved for the two notes that follow: an anticyclotomic character of order $\ell^N$
ramified at one new prime, trivial at a Frobenius element there that lies in a prescribed open set (Theorem~A); labelled columns of row type are rows
(Theorem~B); the forms cut out by an apt idempotent lie in a base level whose eigensystems are labelled columns
with a quadratic relation on the inertia at the twist primes (Theorem~C); the inverse ratios of twisted $L$-values
are controlled by families of interpolated values (Theorem~D); and at one place above~$\ell$ the crystalline
cocycles of a block lie, up to an explicit loss, in the cup-orthogonal of a class of lattices (Theorem~E).

\statement{leancol}{\fontsize{9.6}{13.2}\selectfont\textbf{What this note is.} The artifact \ArtifactName{} is a
body of Lean code: the compact edition of a formal proof, due to L.~Alpöge with the AI system Claude, of the
Riemann hypothesis and of its generalisation to Dirichlet $L$-functions. That headline is a claim. The authors
report that every file compiles and that Lean's kernel accepted the proofs; an independent replay is recorded in
the handout's record of checks when it reports. No mathematician has read the argument through, and nothing here
says that the Riemann hypothesis is proved. This note is the first of three on the unit \emph{The choice of the
twists} (the folder of the formal text on Selmer groups): it treats \UnitInNote{} of the unit's \UnitMembers{}
declarations, those of the first fifteen files. The other declarations are the subject of the notes \emph{Short
Selmer groups after one anticyclotomic twist} and \emph{A fit family of twists exists}. All definitions are
written out; the lemmas are proved from the inputs named in the proofs, which are results of the background
library or of other units and are quoted with their owners. Where a proof below is an outline of a long formal
computation, it says so. Every statement was compared with its Lean text, argument by argument. The Lean text
read is \PinText; Appendix~A is the concordance of \ConcDecls{} declarations.}

\sect{The statements}{s:intro}
The notions are defined in §\ref{s:back} (background) and §\ref{s:defs} (the definitions of the unit); the
conventions are at the head of §\ref{s:back}. Throughout, $\ell$ is a prime number and
\m{\iota\colon\CC\to\Qlb} a fixed isomorphism of fields.

\statement{arith}{\phantomsection\label{t:A}\textbf{Theorem A} (an anticyclotomic character at one new prime).\enspace
\emph{Let $E$ be a CM field as in the conventions, $N$ a natural number, $S$ a finite set of finite places of~$E$
and $R$ a finite set of finite places of~$E$. Let \m{z\in\GE{E}} lie in the normal subgroup of \m{\GE{E}}
generated by the inertia \m{\inab{R\setminus(S\cup\cj{S}\cup\pd{E}{\ell})}}, and let $W$ be a neighbourhood of~$z$
in~\m{\GE{E}}. Then there are a finite place \m{\fq\notin R\cup\cj{R}} of~$E$ whose norm \m{\Nm(\fq)} is a
prime number, and a datum \m{l=(\fP,\gamma,\phi,h)\in\AuxPrimeData{E}{\ell^N}}, such that}
\[
\phi\in W,\qquad h(\phi)=0,\qquad\text{\emph{and}}\quad \IsAnticyclotomicAt{T}{\fq}{l}\ \ \text{\emph{for every set }}T\subseteq S\cup\pd{E}{\ell}.
\]}
\leanline{\lean{existsisAnticyclotomicAt}}

\statement{arith}{\phantomsection\label{t:B}\textbf{Theorem B} (labelled columns of row type are rows).\enspace
\emph{Let $P$ be parameters, $D$ admissible data for~$P$, and $r$, $n$ finite sets. Let $O$ be an auxiliary pair of
type \m{(r,n)} over~\m{\ED{D}}, $\vartheta$ a twist of~$D$ and \m{c=(c_i)_i} a family of polynomials over the
residue fields of the chosen places of~$O$ with \m{\IsApt{O}{c}{\vartheta}}. Let \m{R\colon\AbsGal{D}\to\GL_r(\Qlb)}
be a homomorphism, $k$ a natural number and $\HT$ a type. If
\m{\IsLabelledColumn{O}{k}{\HT}{\twistedLabel{O}{c}{\cdot}}{\emptyset}{R}} and
\m{\IsRowType{D}{\#r}{k}{\HT}}, then \m{\IsRow{D}{k}{\HT}{R}}.}}
\leanline{\lean{isRow}}

\statement{arith}{\phantomsection\label{t:C}\textbf{Theorem C} (a base level for an apt idempotent).\enspace
\emph{Let $P$, $D$ be as in Theorem~B, $r$ a finite set and $n$ a non-empty finite set, $O$ an auxiliary pair of
type \m{(r,n)} over~\m{\ED{D}}, $\vartheta$ a twist of~$D$ with set of places \m{Q(\vartheta)}, and $c$ a family of
polynomials with \m{\IsApt{O}{c}{\vartheta}}. Suppose \m{\IsAuxiliary{\badPlaces{D}{\ell}}{Q(\vartheta)}}. Then
there are a type $\HT$ and a base level $B$ of forms of the weight of~$O$ on the standard definite hermitian
space of rank~\m{\#n}, with bound \m{\ordch{D}}, such that:}
\begin{enumerate}[label=\textup{(\arabic*)},leftmargin=2.3em,itemsep=2pt]
\item \emph{every form in the image of \m{\tupleProj{O}{c}} is a form of~$B$ of level~$\emptyset$;}
\item \emph{for every form $y$ of~$O$ that is $\ell$-integral, \m{\ordch{D}\cdot\tupleProj{O}{c}(y)} is $\ell$-integral;}
\item \emph{for every finite set $K$, every family \m{u=(u_k)_{k\in K}} of finite places of~\m{\ED{D}} none of
which lies in \m{\badPlaces{D}{\ell}}, and every eigensystem $x$ of~$B$ of level~$u$, the representation
\m{R=\normalisedRep{O}{\varrho_x}} satisfies, with the set of places
\m{S_u=Q(\vartheta)\cup\cj{Q(\vartheta)}\cup u(K)\cup\cj{u(K)}} and the polynomials
\m{c'_i=\twistedLabel{O}{c}{i}\cdot(X-1)^{\#n-\#r}},}
\[
\IsLabelledColumn{O}{\#n-\#r}{\HT}{(c'_i)_i}{S_u}{R}
\]
\emph{and \m{\bigl(R(g)-\invAvatar{\vartheta}{g}\bigr)\bigl(R(g)-1\bigr)=0} for every
\m{g\in\inab{Q(\vartheta)\cup\cj{Q(\vartheta)}}}.}
\end{enumerate}}
\leanline{\lean{existsbaseLevel}}

\statement{arith}{\phantomsection\label{t:D}\textbf{Theorem D} (families of values control ratios of $L$-values).\enspace
\emph{Let $P$, $D$ be as in Theorem~B, $w$ a finite place of~\m{\ED{D}} that is chosen and does not lie
above~$\ell$, $p$ a pair of~$D$ that is ramified, and \m{s\in\{1,-1\}} a sign. Then there are a family
\m{\nu=(\nu_i)_{i\in\{0,1\}\times\{1,-1\}}} of homomorphisms \m{\AbsGal{D}\to\Qlb^{\times}} and a function
\m{\mathrm{val}} with \m{\IsFamily{\badPlaces{D}{\ell}}{\nu}{\mathrm{val}}} such that}
\begin{enumerate}[label=\textup{(\arabic*)},leftmargin=2.3em,itemsep=2pt]
\item \emph{\m{\vl\bigl(\nu_i(\tau_{w})-1\bigr)\le\vl\bigl(\Nm(w)-1\bigr)} for every~$i$, where $\tau_w$ is the chosen tame
generator at the chosen place of~$\Qb$ over~$w$;}
\item \emph{for every twist $\tau$ of~$D$ with \m{\IsAuxiliary{\badPlaces{D}{\ell}}{Q(\tau)}} and every real number
\m{a\ge0}: if \m{\vl\bigl(\mathrm{val}(Q(\tau),\invAvatar{\tau}{\cdot})\bigr)\le a+\vl\bigl(\mathrm{val}(\emptyset,1)\bigr)}, then
\m{\vl\bigl(\iota(\varrho_p(\tau,s))\bigr)\le a}, where \m{\varrho_p(\tau,s)} is the ratio of $L$-values of~$p$ at the
twist~$\tau^{s}$.}
\end{enumerate}}
\leanline{\lean{existsisFamily}}

\statement{arith}{\phantomsection\label{t:E}\textbf{Theorem E} (the loss at one place above~$\ell$).\enspace
\emph{Let $P$, $D$ be as in Theorem~B. Let $\Gamma$ be a topological group with a homomorphism
\m{j\colon\Gamma\to\Gl} of open image, and \m{e\colon\ED{D}\to\Qlb} a ring homomorphism such that $j$ is an open
embedding whose image is the group of the elements of \m{\Gl} that fix \m{e(\ED{D})} pointwise. Let $r$, $t$ be natural
numbers and $\HT$ a type with \m{\IsRowType{D}{r}{t}{\HT}}. Then for all large coefficient rings~$\cO$ the following
holds. Let \m{(L_i)_i} be a family of Galois lattices over~$\cO$ for~$\Gamma$ indexed by
\m{i\in\{\ast\}\sqcup\{0,\dots,t-1\}}, let $F$ be a filter on~$\Gamma$, $K$ a relation between the elements of~$\Gamma$
and the polynomials over~$\cO$, and $U$ a function from the \m{\Zl}-algebra homomorphisms \m{\cO\to\Cl} to the
finite sets of integers. Suppose that}
\begin{enumerate}[label=\textup{(\roman*)},leftmargin=2.3em,itemsep=2pt]
\item \emph{either \m{\ell=\Dg{\nu}{D}} for an index~$\nu$ and, for every \m{s\in\Gamma} with \m{j(s)} in the inertia
group of~\m{\Gl}, the property \m{K(s,\cdot)} is of tame type \m{(\ell,\Dr{D},(\Dr{D}+\tD{D})\Dr{D})}; or \m{F=\top} and
$K$ holds for no pair;}
\item \emph{for every homomorphism \m{g'\colon\cO\to\Cl} and every \m{a\in U(g')}, some $L_i$ has the jump $a$
at~$g'$;}
\item \emph{for every \m{g\in\Gl}, the lattice $L_\ast$ has the jumps \m{\HT(\complexEmbOf{e}{g})} at
\m{\coeffEmb{\cO}{g}}, and for \m{c<t} the lattice $L_c$ has the single jump~$c$ at \m{\coeffEmb{\cO}{g}};}
\item \emph{\m{\rk_{\cO}L_\ast=r}.}
\end{enumerate}
\emph{Then for every index $c$, with \m{x=1} if \m{c=\ast} and \m{x=\tD{D}} otherwise: if $L_c$ belongs to the class
\m{\mathcal{Q}(F,K,U)}, then}
\[
\nlen_{\cO}\Bigl(\Zf\bigl(L_c^{\vee}(1)\bigr)\big/\Zf\bigl(L_c^{\vee}(1)\bigr)\cap\bigl(L_c^{\vee}(1)\bigr)^{\perp\mathcal{Q}(F,K,U)}\Bigr)
\ \le\ [\Gl:j(\Gamma)]\cdot\rk_{\cO}L_c\cdot\frac{\loss{D}{\ell}}{x}.
\]}
\leanline{\lean{blockdefectleatplace}}

Theorem~A is the source of the tied primes of the pencils in the second note; the Taylor--Wiles primes of the
third note and the two further primes of its lemma on reciprocity laws come directly from Chebotarev's theorem.
The proof of Theorem~A uses Chebotarev's theorem (existence form) and the existence theorem of class field theory. Theorems B and~C turn the automorphic input (the
representations attached to eigenforms on definite unitary groups, with their local behaviour) into the two
predicates on which Ribet's method works in the following notes. Theorem~D is used to bound the heights of
ratios of $L$-values in the last clause of a fit family. Theorem~E is the local half of the statement that
``crystalline losses above~$\ell$ are affordable''; the global half is in the third note.

\sect{Conventions and background}{s:back}
\begin{conventions}\phantomsection\label{c:conventions}
The natural numbers are \m{0,1,2,\dots}; differences of natural numbers written in exponents and indices are
truncated at~$0$ where the formal text truncates them, and this is said at each place. Throughout, $\ell$ is a
prime number; \m{\Ql}, \m{\Zl} are the $\ell$-adic numbers and integers, \m{\Qlb} is an algebraic closure
of~\m{\Ql} with the absolute value \m{|\cdot|} normalised by \m{|\ell|=1/\ell}, \m{\Zlb} its valuation ring,
\m{\Cl} its completion, and \m{\Gl=\Gal(\Qlb/\Ql)}. Where a statement involves complex numbers and $\ell$-adic
numbers together, \m{\iota\colon\CC\to\Qlb} is a fixed isomorphism of fields. \m{\Qb} is a fixed algebraic closure
of~$\mathbb{Q}$, and $\mathsf{c}$ a fixed complex conjugation in \m{\Gal(\Qb/\mathbb{Q})}. A \emph{CM field as in the
conventions} is a CM number field $E$ (also written $E'$, $E_1$) with a fixed embedding in~\m{\Qb} that
contains a primitive root of unity of order~$5040$; \m{\GE{E}=\Gal(\Qb/E)} with its profinite topology. In
some statements $E$ is only a number field with a fixed embedding in~\m{\Qb}; this is then said. $P$ denotes parameters and $D$ admissible data for~$P$
(Chapter~III); its field \m{\ED{D}} is a CM field as in the conventions (Lemma~\ref{l:small2}(a)). The letters
$n$, $r$, $m$, $d$, $\sigma$, $J$, $X$, $T$, $K$, when they index matrices or families, denote finite sets (the
letter $d$ is the prime of the parameters where it stands with $N_0$, \m{\ND{D}} or under a square root; as an
index set and as an exponent, \m{1\le d\le\dots}, it is a bound variable of its statement, as in the formal text; in
Definition~\ref{d:obey} and in Definitions \ref{d:new}, \ref{d:pencil}, \ref{d:ties} the sets $X$, $J$ and the index
set of the slopes are arbitrary sets); \m{\#n} is the number of elements, \m{M_n(A)} the ring of matrices over a
commutative ring~$A$ with rows and columns indexed by~$n$, and \m{\GL_n(A)} its group of units. A
\emph{coefficient ring} is the ring of integers $\cO$ (also $\cO'$) of a finite Galois extension of~\m{\Ql}
inside~\m{\Qlb}; ``for all large~$\cO$'' means: there is a coefficient ring $\cO_0$ such that the assertion holds
for all coefficient rings \m{\cO\supseteq\cO_0}. $A$, $\Lambda$ denote commutative rings, $\Gamma$ a group
(a topological group where continuity is mentioned), and $B$ a module.
\end{conventions}

The rest of this section recalls, in words, the notions of the background library and of other units in which
the declarations of the unit are stated. They are not declarations of this unit; their formal definitions are in
the files named in the concordances of the notes of their units.

\subhead{2.1. Places, inertia, Frobenius elements.} Let $E$ be a number field embedded in~\m{\Qb}. A
\emph{finite place} $w$ (also $v$, $\fq$) of~$E$ is a non-zero prime ideal of its ring of integers; \m{\Nm(w)}
is its norm, $\F_w$ its residue field. For a natural number $M$, \m{\pd{E}{M}} is the set of finite places
of~$E$ that contain~$M$. When $E$ is a CM field, \m{\cj{w}} is the complex conjugate place, \m{\cj{S}=\{\cj{w}:w\in S\}}
for a set $S$ of places, and $w$ is \emph{split} if \m{\cj{w}\ne w}; a place of prime norm is split. A \emph{place of~\m{\Qb}} is a valuation subring $\fP$ of~\m{\Qb}; it \emph{lies over}
$w$ if an integer of~$E$ belongs to the maximal ideal of~$\fP$ exactly when it belongs to~$w$. For each $w$ the
library fixes one place \m{\fP_w} of~\m{\Qb} over~$w$. The decomposition group \m{G_{\fP}\subseteq\GE{E}}
consists of the $g$ with \m{g(\fP)=\fP}, the inertia group \m{I_{\fP}} of those which moreover act trivially on
the residue field of~$\fP$; for a set $S$ of finite places, \m{\inab{S}} is the union of the \m{I_{\fP}} over
the places $\fP$ of~\m{\Qb} lying over a place of~$S$. For a natural number $q$, an element \m{\phi\in G_{\fP}} is a
\emph{$q$-Frobenius at~$\fP$} if \m{\phi(x)\equiv x^{q}} modulo the maximal ideal of~$\fP$ for all \m{x\in\fP};
for \m{q=\Nm(w)} this is an arithmetic Frobenius element. For $\fP$ over~$w$ the library fixes a
\emph{tame generator} \m{\tau_{\fP}\in I_{\fP}}: the local Artin symbol of a unit of the completion $E_w$ whose
residue generates \m{\F_w^{\times}}; we write \m{\tau_w=\tau_{\fP_w}}. \m{\fP^{\mathsf c}=\mathsf{c}^{-1}(\fP)} is the conjugate place. For
a CM field, \m{g^{\mathsf c}=\mathsf{c}g\mathsf{c}} for \m{g\in\GE{E}}, and for a sign $b$ we put
\m{g^{[+]}=g} and \m{g^{[-]}=(g^{\mathsf c})^{-1}}. \m{\cyc\colon\GE{E}\to\Zl^{\times}} is the $\ell$-adic
cyclotomic character, read in any \m{\Zl}-algebra; \m{\cyc_M\colon\GE{E}\to(\mathbb{Z}/M)^{\times}} is the action
on the $M$-th roots of unity. The \emph{Artin map} \m{\Art} sends the finite ideles of~$E$ to the abelianisation of~\m{\GE{E}}; we write
\m{\Art(x)=g} when $g$ has the class~\m{\Art(x)}.

\subhead{2.2. Characters from $\mathrm{U}(1)$.} Let $E$ be a CM field as in the conventions, $S$ and $I$ sets of
finite places of~$E$, and \m{\chi\colon\GE{E}\to C} a homomorphism to a commutative monoid. Then $\chi$ is
\emph{from $\mathrm{U}(1)$ relative to~$S$ with modulus~$I$} if $\chi$ is locally constant and there is a
homomorphism $\eta$ from the finite ideles of~$E$ to~$C$, trivial on the principal ideles and on every idele
that is a unit at all places outside \m{S\cup\cj{S}} and is congruent to~$1$ at the places of
\m{I\setminus(S\cup\cj{S})}, such that \m{\chi(g)=\eta(x/\cj{x})} whenever \m{\Art(x)=g}. So such a character is
anticyclotomic, trivial on the decomposition groups above \m{S\cup\cj{S}}, at most tamely ramified at~$I$ and
unramified elsewhere.

\subhead{2.3. Relations and Selmer modules.} Let $\Gamma$ be a group (or a set), $n$ a finite set and $A$ a
commutative ring. A \emph{relation} is a finitely supported function \m{\theta\colon\Gamma\to M_n(A)}; a function
\m{c\colon\Gamma\to B^n} to the $n$-th power of an $A$-module \emph{obeys}~$\theta$ if
\m{\sum_{g}\theta(g)\,c(g)=0}. For a set $L$ of relations, \m{\Ob{L}{B}} is the module of the functions that obey
all of them; for a module $Z$ of functions \m{\Gamma\to A^n}, \m{\Rel{Z}} is the set of the relations that all
elements of $Z$ obey. For a map \m{\varrho\colon\Gamma\to M_n(A)}, \m{\Coc{\varrho}} is the set of the relations
\m{\delta_{gh}-\delta_g-\varrho(g)\delta_h}, where $\delta_g$ is the function with value the identity matrix
at~$g$ and $0$ elsewhere; a function obeys them exactly when it is a cocycle, \m{c(gh)=c(g)+\varrho(g)c(h)}. For
\m{\Gamma=\GE{E}}, a set $S$ of finite places and a set $L$ of relations, the \emph{Selmer module}
\m{\Sel{\varrho}{S}{L}{B}} is the module of the locally constant cocycles \m{\GE{E}\to B^n} for~$\varrho$ that
vanish on \m{\inab{S^{\complement}}}, the inertia above the places outside~$S$, and obey~$L$. For a place $\fP$
of~\m{\Qb}, the module of \emph{local cocycles} at~$\fP$ consists of the restrictions to \m{G_{\fP}} of the
continuous cocycles of~\m{\GE{E}} for~$\varrho$ with values in~$B^n$ that obey~$L$. For \m{a\in A} we write \m{A/a}
for the quotient by the ideal generated by~$a$. \m{\len_A} is the length of an $A$-module; for a coefficient
ring~$\cO$ with uniformiser $\varpi$, \m{\nlen_{\cO}=\vl(\varpi)\cdot\len_{\cO}} is the normalised length, so
that \m{\nlen_{\cO}(\cO/\ell)=1}. \m{\vl\colon\Qlb\to[0,\infty]} is the valuation \m{-\log_\ell|x|}, truncated: it is
$\infty$ at~$0$ and $0$ when \m{|x|\ge1}. For a coefficient ring $\cO$ and \m{e\in\NN}, \m{\cO(-e)} is the
representation \m{g\mapsto\cyc(g)^{-e}} of~\m{\GE{E}} on~$\cO$, a homomorphism to \m{M_1(\cO)}.

\subhead{2.4. Representations and their local conditions.} Let $E$ be a CM field as in the conventions and $n$
a finite set. We consider homomorphisms \m{R\colon\GE{E}\to\GL_n(\Qlb)} (continuity is said when it is assumed).
A homomorphism \m{u\colon\GE{E}\to M_n(\cO)} is a \emph{lattice of~$R$} if \m{u(g)=Y^{-1}R(g)Y} for some
\m{Y\in\GL_n(\Qlb)} and all~$g$. The \emph{trace form of~$R$ is non-degenerate} if an element $x$ of the linear
span of \m{R(\GE{E})} with \m{\tr(x\,R(g))=0} for all $g$ is zero. Let $\fP$ be a place of~\m{\Qb}. An
\emph{$\ell$-adic chart at~$\fP$} is an embedding \m{\Qb\to\Qlb} that induces~$\fP$, together with the open
embedding \m{G_{\fP}\to\Gl} that it defines; it exists when $\fP$ lies above~$\ell$. $R$ is \emph{of finite type
at~$\fP$} if it is trivial on the elements of \m{I_{\fP}} near~$1$ when $\ell$ is a unit at~$\fP$, and potentially
crystalline along every $\ell$-adic chart at~$\fP$. $R$ is \emph{good at~$\fP$} if it is trivial on \m{I_{\fP}} when
$\ell$ is a unit at~$\fP$, and crystalline along every chart at~$\fP$. A \emph{type} is a function $\HT$ from the
complex embeddings of~$E$ to the finite multisets of integers. $R$ \emph{has the jumps~$\HT$ at~$\fP$} if along
every chart at~$\fP$, for every embedding \m{E\to\Qlb} compatible with the chart, read as a complex embedding
$\sigma$ through~$\iota$, the Hodge--Tate jumps of~$R$, with multiplicities, form the multiset \m{\HT(\sigma)}.
For \m{s\in\GE{E}} and a polynomial $Q$ over~\m{\Qlb}, \emph{$Q$ kills $s$ on~$R$ at~$\fP$} if \m{Q(R(s))=0} when
$\ell$ is a unit at~$\fP$, and if a corresponding condition on the crystalline Frobenius holds along every chart
at~$\fP$ when $\fP$ lies above~$\ell$. Let $\fP$ lie over a finite place~$w$, let $T$ be an injective character
\m{\bar{\F}_w^{\times}\to\CC^{\times}} (the library fixes one, the \emph{faithful character} of~$\F_w$) and
\m{f\ge1}. An element \m{s\in I_{\fP}} has \emph{tame exponent~$a$} (for $T$ and~$f$) if for some $z$ with
\m{z^{\Nm(w)^f-1}} a uniformiser of~$E$ at~$w$, $T$ sends the residue of \m{s(z)/z} to
\m{\exp(2\pi i\,a/(\Nm(w)^f-1))}. For a polynomial $c$ over~$\F_w$, $R$ is \emph{labelled by~$c$ at~$\fP$} (for $T$
and~$f$) if for every $s$ of tame exponent~$a$ the polynomial \m{\prod_{x}(X-\iota(T(x)^{a}))}, the product over
the roots $x$ of~$c$ in~$\bar{\F}_w$ with multiplicities, kills $s$ on~$R$ at~$\fP$. For a multiplicative
character $\chi$ of~$\F_w$, \m{x_\chi\in\F_w^{\times}} is the element that names it: \m{T^{a}(x_\chi)=\chi(u)} when
$s$ has tame exponent~$a$ and is the local Artin symbol of a unit with residue~$u$; \m{\chi\mapsto x_\chi} is an
isomorphism of groups.

\subhead{2.5. Crystalline cocycles and classes of lattices.} A \emph{group over~\m{\Gl}} is a group $\Gamma$ with
a homomorphism \m{j\colon\Gamma\to\Gl} of open image. A \emph{Galois lattice over~$\cO$ for~$\Gamma$} is a finite free
$\cO$-module $L$ with an $\cO$-linear action of~$\Gamma$; \m{\rk_{\cO}L} is its rank. \m{\Zf(L)} is the
$\cO$-module of the \emph{crystalline cocycles}: the maps \m{c\colon\Gamma\to L} for which some $y$ in
\m{B_{\mathrm{cris}}\otimes L} has \m{1\otimes c(g)=g\,y-y} for all~$g$ (Colmez's ring $B_{\max}$ serves as period
ring). $L$ is \emph{potentially crystalline} if the periods fixed by the $g$ with \m{j(g)} near~$1$ span
\m{B_{\mathrm{cris}}\otimes L}. For a \m{\Zl}-algebra homomorphism \m{g'\colon\cO\to\Cl} and an integer~$u$, the
\emph{multiplicity of the jump $u$ of~$L$ at~$g'$} is the dimension of the space of $g'$-semilinear maps
\m{L\to\Cl} on which $\Gamma$ acts, near~$1$, through \m{\cyc^{-u}}; $L$ \emph{has the jumps~$\mathfrak{M}$ at~$g'$}, for a
multiset $\mathfrak{M}$ of integers, if these multiplicities are the multiplicities in~$\mathfrak{M}$. \m{L^{\vee}(1)=\Hom_{\cO}(L,\cO(1))}
is the twisted dual. For a filter $F$ on~$\Gamma$, a relation $K$ between elements of~$\Gamma$ and polynomials
over~$\cO$, and a function $U$ from the homomorphisms \m{\cO\to\Cl} to the finite sets of integers, the class
\m{\mathcal{Q}(F,K,U)} consists of the Galois lattices $Y$ whose periods fixed near~$1$ and by $F$-almost all
elements span, whose jumps at each $g'$ lie in \m{U(g')}, and on which $Q$ kills~$s$ whenever \m{K(s,Q)} holds
and \m{j(s)} lies in the inertia group. For \m{F=\top} and $K$ empty this is the class of the crystalline lattices
with jumps in~$U$. A property of polynomials is \emph{of tame type \m{(\ell,\varrho,M)}} if it holds for
\m{X^{M'}-1}, with $M'$ the least common multiple of the \m{\ell^{i}-1}, \m{1\le i\le\varrho}, and for some monic polynomial
over~\m{\Zl} of degree at most~$M$. For a class $\mathcal{Q}$ and a lattice~$L$, the \emph{cup-orthogonal}
\m{L^{\perp\mathcal{Q}}} is the $\cO$-module of the maps \m{c\colon\Gamma\to L} whose cup product with every
crystalline cocycle of every \m{Y\in\mathcal{Q}}, under every pairing of $Y$ and~$L$ modulo a power of~$\ell$ with
values in~\m{\cO(1)}, is a coboundary. These notions belong to the units on $\ell$-adic Hodge theory and on
lifting (Chapter~VIII, the note on crystalline lifts); only their names enter the proofs below.

\subhead{2.6. Characters of $\mathrm{U}(1)$, forms and labels.} Let $E$ be a CM field with its Weil data
(Chapter~VII), which include a Hecke character $\upsilon$ of~$E$ with \m{\upsilon(\varpi_w)\upsilon(\varpi_{\cj w})=1}
at the chosen uniformisers. A \emph{character of~$\mathrm{U}(1)$} is a continuous character $\mu$ of the norm-one
adeles of~$E$, trivial on the norm-one elements of~$E$; it has \emph{weight} \m{e=(e_v)_v} (integers indexed by the
infinite places) and \emph{tame part} \m{\psi=(\psi_w)_w} (characters of the groups \m{\F_w^{\times}}, almost all
trivial) if \m{\mu(u)=\prod_v u_v^{e_v}\prod_w\psi_w(u\bmod w)} for $u$ with integral finite part. For integers $k$,
$j$ the \emph{algebraic Hecke character} \m{\mu[k,j]} is \m{x\mapsto\mu(x/\cj{x})\,\upsilon(x)^{-k}\,|x|^{j-k/2}}, and its
\emph{avatar} \m{\av(\mu[k,j])\colon\GE{E}\to\Qlb^{\times}} is the continuous character with
\m{\av(\mu[k,j])(\Art(x))=\iota(\mu[k,j](x))} times the usual factor at the places above~$\ell$. For admissible
data~$D$ (Chapters III and~V): a \emph{tame character} $\vartheta$ of~$D$ is a character of~$\mathrm{U}(1)$ of
weight~$0$ whose tame part \m{\psi_\vartheta} is trivial off a finite set \m{Q(\vartheta)} of finite places, with
\m{Q(\vartheta)} disjoint from its conjugate and of norms prime to \m{\gD{D}\lvl{D}}; a \emph{twist} is a tame character
whose tame part is non-trivial at every place of~\m{Q(\vartheta)}. A \emph{pair} \m{p=(\nu_p,z_p)} of~$D$ consists of an
integer \m{1\le z_p\le\tD{D}} and a character $\nu_p$ of~$\mathrm{U}(1)$ with tame part supported on the chosen places;
it is \emph{ramified} if its tame part is non-trivial exactly at the chosen places, and \m{\varrho_p(\vartheta,s)} is
its ratio of twisted to untwisted $L$-values (unit of the exchange identity, Chapter~VII). A finite place of~\m{\ED{D}} is
\emph{chosen} if it contains the conjugate of the integer \m{\pi_g} of norm \m{\gD{D}}; \m{\ordch{D}} is the product
of the orders of the groups \m{\GL_{\Dr{D}\tD{D}}(\F_w)} over the chosen places; \m{\omega_D} is the family of
residue characters of the data, of order \m{\Mom{P}{\Dmu{D}}=(N_0d)^{\Dmu{D}+6}} at the chosen places (Chapter~III). An \emph{auxiliary pair $O$ of type
\m{(r,n)} over~$E$} (Chapter~VII) carries a system of places (chosen places \m{w_i}, \m{i\in I_O}, and twist places),
a hermitian space of rank~\m{\#r} with signs, an embedding \m{r\hookrightarrow n}, a character $\mu_O$ of~$\mathrm{U}(1)$
with tame part~$\psi_O$, the character \m{a_O=\upsilon_{\mathrm{U}(1)}^{r_-}} (with $r_-$ the number of signs~$-1$),
and a space \m{\mathcal{F}(O)} of algebraic forms on the standard definite space of rank~\m{\#n}, on which the
group algebra of \m{\prod_i\GL_r(\F_{w_i})} and the Hecke operators at the split places off the places of~$O$ act;
\m{\#n-\#r} (truncated) is its \emph{rank difference}. For an eigenform, its \emph{Hecke polynomial} at a split
place and the representations \emph{attached} to its system of Hecke polynomials are those of Chapter~VII:
attachment asks continuity, a non-degenerate trace form, and that characteristic polynomials of Frobenius elements
are $\ell$-adically close to the Hecke polynomials at most Frobenius elements. Precisely, a \emph{Frobenius
datum} of a number field \m{E\subseteq\Qb} is a triple \m{(w,\fP,\phi)} of a finite place of prime norm, a place
of~\m{\Qb} over it and an element \m{\phi\in\GE{E}} whose inverse is an \m{\Nm(w)}-Frobenius at~$\fP$. A set $W$ of
Frobenius data \emph{contains most Frobenius data} if every non-empty open subset of~\m{\GE{E}} contains a
non-empty open subset $U$ such that, for all but finitely many finite places~$w$, every Frobenius datum
\m{(w,\fP,\phi)} with \m{\phi\in U} belongs to~$W$; it \emph{contains most Frobenius data near~$g$} if it contains
the intersection of such a set with the set of the data whose element $\phi$ lies in some neighbourhood of~$g$;
and a statement holds \emph{for most Frobenius data} (near~$g$) if the set of the data for which it holds has
this property. A \emph{label of degree~$\delta$}
over a finite field is a monic squarefree polynomial of degree~$\delta$ with non-zero constant term, a
\emph{primitive label} an irreducible one with a root of order \m{q^{\delta}-1}; \m{\tupleIdem{r}{c}} is the
idempotent of a tuple $c$ of labels for the index set~$r$, in the group algebra of the product of the groups
\m{\GL_r} of the residue fields (note \emph{Primitive labels}). A \emph{base level} with bound~$M$ (unit of the adelic cusp
forms) is a finite set of places with a space of forms of open level there, stable under the operators that
commute with the translations at these places, all stabilisers having order dividing~$M$; its \emph{eigensystems
of level~$u$} are the isotypic spaces that contain a non-zero form of the base level fixed by the principal
level at the places~$u_k$, each with its chosen attached representation~\m{\varrho_x}.

\sect{The definitions of the unit}{s:defs}
The definitions are given in the order of the formal text. Each names all its arguments; the implicit ones are
the ambient objects of the conventions.

\subhead{3.1. Cocycles, blocks and the crystalline relations.}

\begin{definition}[cocycles that obey the relations among the $u_x(g)-1$]\label{d:obey}
Let $A$ be a commutative ring, $\Gamma$ a group, $X$ a set, $\sigma$ and $d$ finite sets, \m{u=(u_x)_{x\in X}} a
family of homomorphisms \m{u_x\colon\Gamma\to M_\sigma(A)}, \m{\varrho\colon\Gamma\to M_d(A)} a homomorphism and
\m{a\in A}. Then \m{\obeyingCocycles{u}{\varrho}{a}} is the $A$-module of the maps \m{c\colon\Gamma\to(A/a)^d},
continuous or not, that are cocycles for~$\varrho$ and satisfy \m{\sum_g\theta(g)\,c(g)=0} for every finitely
supported function \m{\theta\colon\Gamma\to A} with \m{\sum_g\theta(g)\,(u_x(g)-1)=0} in \m{M_\sigma(A)} for all
\m{x\in X}.
\end{definition}
\leanline{\lean{obeyingCocycles}}

\begin{definition}[blocks of a row with its string]\label{d:blocks}
Let $r$ be a set, $t$ a natural number, and $i$ an element of \m{\{\ast\}\sqcup\{0,\dots,t-1\}}. Then
\m{\RowBlockIndex{r}{i}} is the set $r$ if \m{i=\ast}, and a set with one element otherwise; it is finite when
$r$ is. Let further $E$ be a number field embedded in~\m{\Qb}, $\cO$ a coefficient ring, $d$ a finite set and
\m{T\colon\GE{E}\to M_d(\cO)} a homomorphism. For \m{g\in\GE{E}}, the matrix \m{\blockRep{T}{i}{g}}, indexed by
\m{\RowBlockIndex{d}{i}}, is \m{T(g)} if \m{i=\ast} and the scalar \m{\cyc(g)^{-e}} if \m{i=e<t}. So
\m{\blockRep{T}{i}{\cdot}} is $T$ for \m{i=\ast} and \m{\cO(-e)} for \m{i=e}.
\end{definition}
\leanline{\lean{RowBlockIndex}\sep \lean{blockRep}\sep the finiteness of \m{\RowBlockIndex{d}{i}} is an unnamed instance}

\begin{definition}[the chart at a place above $\ell$]\label{d:chart}
Let $E$ be a CM field as in the conventions and \m{v\in\pd{E}{\ell}}. The formal text fixes an $\ell$-adic chart at the
place \m{\fP_v} of~\m{\Qb}: an embedding \m{e_v\colon\Qb\to\Qlb} under which \m{\Zlb} pulls back to~\m{\fP_v},
with an open embedding \m{j_v\colon G_{\fP_v}\to\Gl}, compatible with~$e_v$, whose image is the group of the
elements that fix \m{e_v(E)} pointwise. It makes \m{G_{\fP_v}} a group over~\m{\Gl}.
\end{definition}
\leanline{\lean{localChart}\sep an instance, chosen from two existence statements of the background library}

\begin{definition}[embeddings read through $\Gl$]\label{d:coeffemb}
Let $\cO'$ be a coefficient ring and \m{g\in\Gl}. Then \m{\coeffEmb{\cO'}{g}\colon\cO'\to\Cl} is the
\m{\Zl}-algebra homomorphism \m{x\mapsto g(x)}. Let $E'$ be a field and \m{e\colon E'\to\Qlb} a ring homomorphism. Then
\m{\complexEmbOf{e}{g}=\iota^{-1}\circ g^{-1}\circ e} is a ring homomorphism \m{E'\to\CC}.
\end{definition}
\leanline{\lean{coeffEmb}\sep \lean{complexEmbOf}}

\begin{definition}[the crystalline relations]\label{d:zf}
Let $E$ be a CM field as in the conventions, $\cO$ a coefficient ring, $n$ a finite set and
\m{u\colon\GE{E}\to M_n(\cO)} a homomorphism. For \m{v\in\pd{E}{\ell}}, \m{\ZfAt{v}{u}} is the $\cO$-module
\m{\Zf(L)} of the crystalline cocycles of the lattice \m{L=\cO^n} on which \m{G_{\fP_v}} acts through~$u$, the group
\m{G_{\fP_v}} being a group over~\m{\Gl} by Definition~\ref{d:chart}. The set \m{\crysRels{u}} of the \emph{crystalline
relations of~$u$} consists of the relations on~\m{\GE{E}} obtained, for some \m{v\in\pd{E}{\ell}}, from a relation
in \m{\Rel{\ZfAt{v}{u}}} by extension by zero from \m{G_{\fP_v}} to~\m{\GE{E}}.
\end{definition}
\leanline{\lean{ZfAt}\sep \lean{crysRels}}

A function on~\m{\GE{E}} obeys \m{\crysRels{u}} exactly when, for each $v$ above~$\ell$, its restriction to
\m{G_{\fP_v}} obeys every relation of the crystalline cocycles there (Lemma~\ref{l:small1}(a)). This is the
Bloch--Kato finite condition at the places above~$\ell$, written as linear relations on global cocycles.

\begin{definition}[the loss of a block]\label{d:blockloss}
Let $E$, $\cO$, $n$ and \m{u\colon\GE{E}\to M_n(\cO)} be as in Definition~\ref{d:zf}, let $X$ be a set, $\sigma$ a
finite set, \m{w=(w_x)_{x\in X}} a family of homomorphisms \m{w_x\colon\GE{E}\to M_\sigma(\cO)}, and
\m{M\in\NN}. Then
\[
\blockLoss{u}{w}{M}=\len_{\cO}\Bigl(\obeyingCocycles{w}{u}{\ell^M}\big/\obeyingCocycles{w}{u}{\ell^M}\cap\Ob{\crysRels{u}}{\cO/\ell^M}\Bigr).
\]
\end{definition}
\leanline{\lean{blockLoss}}

\begin{definition}[the sign of an embedding]\label{d:cmsign}
Let $E$ be a field with a choice, for each infinite place, of a complex embedding that induces it
(for a CM field: a CM type). For a complex embedding $\sigma$ of~$E$, \m{\cmSign{\sigma}} is $1$ if $\sigma$ is the
chosen embedding of its place, and $-1$ otherwise.
\end{definition}
\leanline{\lean{cmSign}}

\begin{definition}[row types]\label{d:rowtype}
Let $D$ be admissible data, $r$ and $t$ natural numbers and $\HT$ a type for \m{\ED{D}}. Then
\m{\IsRowType{D}{r}{t}{\HT}} holds if
\begin{enumerate}[label=\textup{(R\arabic*)},leftmargin=3em,itemsep=1pt]
\item \m{1\le r\le\Dr{D}} and \m{t\le\tD{D}};
\item there is a natural number \m{j\le\Dm{D}} such that for every complex embedding $\sigma$ of~\m{\ED{D}} and
every \m{u\in\HT(\sigma)}: either \m{\hhD{D}\le2|u|} and \m{|u|\le\hhD{D}^2} (a \emph{low} jump), or
\m{\bigl|\cmSign{\sigma}\,u+\wt{D}{j}{\sigma}\bigr|\le\tD{D}}, where \m{\wt{D}{j}{\sigma}} is the weight of the
$j$-th lift at the infinite place of~$\sigma$;
\item for every $\sigma$, \m{\#\{u\in\HT(\sigma):|u|\le2\hhD{D}^2\}\cdot\kD{D}\le r\cdot\Dm{D}};
\item for every $\sigma$, \m{\#\{u\in\HT(\sigma):u<0\}+\#\{u\in\HT(\bar\sigma):u<0\}\le r}, with $\bar\sigma$ the
complex conjugate of~$\sigma$.
\end{enumerate}
The counts are with multiplicity.
\end{definition}
\leanline{\lean{IsRowType}}

\begin{definition}[the local loss]\label{d:loss}
For natural numbers $h$, $b$, $l$, \m{\lambdaLoss{h}{b}{l}} is the exponent of~$l$ in
\m{(2b)!\cdot\bigl(\prod_{b-30h^3<k\le b+30h^3}k\bigr)^{h^3}}, the product over the natural numbers \m{k\ge1} in that
range. For admissible data $D$ and a natural number~$l$,
\[
\loss{D}{l}=10^{6}\cdot\frac{\Dm{D}}{\kD{D}}\Bigl(\bone{D}\cdot v_l(\gD{D})+\lambdaLoss{\hhD{D}}{\bone{D}}{l}\Bigr),
\]
a non-negative real number; \m{v_l} is the exponent of~$l$ (zero when $l$ is not a prime).
\end{definition}
\leanline{\lean{lambdaLoss}\sep \lean{loss}}

\subhead{3.2. The field of the data and the avatar of a tame character.}

\begin{definition}[places of the data]\label{d:places}
Let $D$ be admissible data. Then \m{\AbsGal{D}} is the group \m{\Gal(\Qb/\ED{D})}, and \m{\placesAboveG{D}} is the set
of the finite places of~\m{\ED{D}} that contain \m{\Dg{0}{D}} or \m{\Dg{1}{D}}. For a natural number~$\ell$ (a prime, in all uses),
\m{\badPlaces{D}{\ell}} is the finite set of the finite places of~\m{\ED{D}} that contain
\m{\ell\cdot\gD{D}\cdot\lvl{D}}.
\end{definition}
\leanline{\lean{AbsGal}\sep \lean{placesAboveG}\sep \lean{badPlaces}\sep the structure of a CM field as in the
conventions on \m{\ED{D}} is an unnamed instance}

\begin{definition}[the avatar of the inverse of a tame character]\label{d:invav}
Let $D$ be admissible data and $\vartheta$ a tame character of~$D$. For \m{g\in\AbsGal{D}},
\m{\invAvatar{\vartheta}{g}=\av(\vartheta^{-1}[0,0])(g)\in\Qlb}: the value at~$g$ of the $\ell$-adic avatar, through~$\iota$, of
the Hecke character \m{x\mapsto\vartheta^{-1}(x/\cj{x})}.
\end{definition}
\leanline{\lean{invAvatar}}

\subhead{3.3. Columns, rows and labels.}

\begin{definition}[conjugate self-dual]\label{d:csd}
Let $D$ be admissible data, $n$ a finite set, $t$ a natural number and \m{R\colon\AbsGal{D}\to\GL_n(\Qlb)} a
homomorphism. Then \m{\IsConjSelfDual{t}{R}} holds if $R$ is continuous, the trace form of~$R$ is non-degenerate,
and \m{\tr R(g^{\mathsf c})=\cyc(g)^{-(t-1)}\,\tr R(g^{-1})} for all \m{g\in\AbsGal{D}}, where \m{t-1} is read as $0$
for \m{t=0}.
\end{definition}
\leanline{\lean{IsConjSelfDual}}

\begin{definition}[the crude unit]\label{d:cr}
For a natural number $\ell$ and admissible data~$D$, \m{\crudeUnit{\ell}{D}=v_\ell(\ordch{D})+2^{-\ell}}, a positive
real number.
\end{definition}
\leanline{\lean{crudeUnit}}

\begin{definition}[the Selmer length bound]\label{d:short}
Let $D$ be admissible data, $n$ a finite set, \m{\psi\colon\AbsGal{D}\to\Qlb} a multiplicative homomorphism,
\m{R\colon\AbsGal{D}\to\GL_n(\Qlb)} a homomorphism, $Q$ a set of finite places of~\m{\ED{D}}, \m{m\in\NN},
\m{\gamma\in\AbsGal{D}} and \m{z\in\Zlb}. Then \m{\SelmerLengthLe{\psi}{R}{Q}{m}{\gamma}{z}} holds if for all
large~$\cO$ there are a homomorphism \m{u\colon\AbsGal{D}\to M_n(\cO)}, a multiplicative homomorphism
\m{\psi'\colon\AbsGal{D}\to\cO} and \m{y\in\cO} such that \m{y=z}, $u$ is a lattice of~$R$, \m{\psi'(g)=\psi(g)}
for all~$g$, and, with \m{\varrho=\psi'\cdot u},
\[
\varrho(\gamma)=y\cdot1,\qquad y-1\in\cO^{\times},\qquad\cyc(\gamma)=1,
\]
\begin{align*}
&\nlen_{\cO}\bigl\{c\in\Sel{\varrho}{\pd{\ED{D}}{\ell}\cup\placesAboveG{D}\cup Q\cup\cj{Q}}{\crysRels{\varrho}}{\cO/\ell^{m}}:c(\gamma)=0\bigr\}\\
&\qquad\le\ 10\,\fD{D}\cdot\#n\cdot\crudeUnit{\ell}{D}+\crudeUnit{\ell}{D}.
\end{align*}
\end{definition}
\leanline{\lean{SelmerLengthLe}}

\begin{definition}[columns, as Galois representations]\label{d:column}
Let $D$ be admissible data, $n$ a finite set, $t$ a natural number, $\HT$ a type for \m{\ED{D}}, $S$ a set of
finite places of~\m{\ED{D}} and \m{R\colon\AbsGal{D}\to\GL_n(\Qlb)} a homomorphism. Then
\m{\IsColumn{D}{t}{\HT}{S}{R}} (\emph{$R$ is a column of level~$S$}) holds if \m{\IsConjSelfDual{t}{R}} and, for every
finite place $v$ of~\m{\ED{D}} that does not lie in \m{S\setminus(\pd{\ED{D}}{\ell}\cup\placesAboveG{D})} and every place
$\fP$ of~\m{\Qb} over~$v$: $R$ is of finite type at~$\fP$, has the jumps $\HT$ at~$\fP$, and, if
\m{v\notin\placesAboveG{D}}, is good at~$\fP$.
\end{definition}
\leanline{\lean{IsColumn}}

In the chapters on forms a column is an algebraic modular form; here, as in the formal text of this unit, a
column is a Galois representation of the kind attached to such forms, and a row is one of the kind that occurs
as the irreducible block of the representation of a theta lift.

\begin{definition}[rows, as Galois representations]\label{d:row}
Let $D$, $n$, $t$, $\HT$ and $R$ be as in Definition~\ref{d:column}. Then \m{\IsRow{D}{t}{\HT}{R}} holds if
\begin{enumerate}[label=\textup{(W\arabic*)},leftmargin=3em,itemsep=1pt]
\item \m{\IsColumn{D}{t}{\HT}{\emptyset}{R}} and \m{\IsRowType{D}{\#n}{t}{\HT}};
\item for every \m{v\in\placesAboveG{D}\setminus\pd{\ED{D}}{\ell}} and every place $\fP$ of~\m{\Qb} over~$v$
there are \m{\gamma\in\AbsGal{D}} and \m{N\in\NN} such that \m{\gamma=\tau_{\fP}} if $v$ is chosen, $N$ is a unit
at~$\fP$, \m{R(i^{N})=1} for all \m{i\in I_{\fP}}, and
\m{\vl\bigl(\det(R(\gamma^{[b]})-1)\bigr)\le\#n\cdot\crudeUnit{\ell}{D}} for both signs~$b$;
\item for all large $\cO$ and every lattice \m{u\colon\AbsGal{D}\to M_n(\cO)} of~$R$: the $\cO$-span of
\m{u(\AbsGal{D})} is \m{M_n(\cO)}, and there are \m{g\in\AbsGal{D}}, an index \m{i\in\{0,1\}} and \m{\alpha\in\cO}
with \m{\cyc(g)=\Dg{i}{D}^{\#n}} and \m{u(g)=\alpha\cdot1};
\item there are \m{\gamma\in\inab{\placesAboveG{D}\setminus\pd{\ED{D}}{\ell}}} and \m{y\in\Qlb} with
\m{R(\gamma)=y\cdot1} and \m{\bigl|(y^{\ell^{a}})^{j}-1\bigr|=1} for all \m{a\in\NN} and all \m{1\le j\le4\,\#n+2}.
\end{enumerate}
\end{definition}
\leanline{\lean{IsRow}}

\begin{definition}[the normalising character]\label{d:chione}
Let $D$ be admissible data, $r$, $n$ finite sets and $O$ an auxiliary pair of type \m{(r,n)} over~\m{\ED{D}}. Then
\m{\chiOneInv{O}} is the character \m{\av\bigl(a_O^{-1}[-\#r,0]\bigr)} of~\m{\AbsGal{D}} with values
in~\m{\Qlb^{\times}}, the avatar of the Hecke character
\m{x\mapsto a_O^{-1}(x/\cj{x})\,\upsilon(x)^{\#r}\,|x|^{\#r/2}}. For a homomorphism $R$ from \m{\AbsGal{D}} to the
invertible matrices over~\m{\Qlb} indexed by a finite set, \m{\normalisedRep{O}{R}} is the homomorphism
\m{g\mapsto\chiOneInv{O}(g)\cdot R(g)}.
\end{definition}
\leanline{\lean{chiOneInv}\sep \lean{normalisedRep}}

\begin{definition}[twisted labels]\label{d:twl}
Let $D$, $r$, $n$, $O$ be as in Definition~\ref{d:chione}, let \m{c=(c_i)_{i\in I_O}} be a family of polynomials,
$c_i$ over the residue field of the $i$-th chosen place $w_i$ of~$O$, and let \m{i\in I_O}. Then
\m{\twistedLabel{O}{c}{i}} is the polynomial of the degree of~$c_i$ whose roots are the roots of~$c_i$ multiplied
by \m{x_{\psi_{O,w_i}^{-1}}\in\F_{w_i}^{\times}}, the element that names the inverse of the tame part of~$\mu_O$
at~$w_i$.
\end{definition}
\leanline{\lean{twistedLabel}}

\begin{definition}[apt labels]\label{d:apt}
Let $D$, $r$, $n$, $O$, $c$ be as in Definition~\ref{d:twl} and $\vartheta$ a twist of~$D$. Then
\m{\IsApt{O}{c}{\vartheta}} holds if
\begin{enumerate}[label=\textup{(A\arabic*)},leftmargin=3em,itemsep=1pt]
\item the system of places of~$O$ is that of~$\vartheta$: its chosen places are the chosen places of~$D$ and its
twist places are the places of~\m{Q(\vartheta)};
\item \m{\psi_O=\omega_D^{\,j}\cdot\psi_\vartheta} for some \m{1\le j\le\Dm{D}};
\item \m{\#r} and the rank difference \m{\#n-\#r} are odd, \m{\#r\le\Dr{D}} and \m{\#n-\#r\le\tD{D}};
\item every $c_i$ is a label of degree~\m{\#r};
\item for each \m{k\in\{0,1\}} some chosen place $w_i$ lies above~\m{\Dg{k}{D}} with: \m{\#r=\Dr{D}} and $c_i$ a
primitive label of degree~\m{\#r}, or \m{\#r=1} and \m{c_i=X-1};
\item for every~$i$, $1$ is not a root of \m{\twistedLabel{O}{c}{i}}.
\end{enumerate}
\end{definition}
\leanline{\lean{IsApt}}

\begin{definition}[the projector of a tuple of labels]\label{d:tproj}
Let $D$, $r$, $n$, $O$, $c$ be as in Definition~\ref{d:twl}. Then \m{\tupleProj{O}{c}} is the endomorphism of the
space \m{\mathcal{F}(O)} of forms by which the element \m{\tupleIdem{r}{c}} of the group algebra of \m{\prod_i\GL_r(\F_{w_i})}
acts.
\end{definition}
\leanline{\lean{tupleProj}}

\begin{definition}[a pair that has a row]\label{d:hasrow}
Let $D$, $r$, $n$, $O$ be as in Definition~\ref{d:chione}, let $z$ be Hecke data in rank~\m{\#r} (a polynomial for
each finite place) and \m{R\colon\AbsGal{D}\to\GL_r(\Qlb)} a homomorphism. Then \m{\HasRow{O}{z}{R}} holds if
\begin{enumerate}[label=\textup{(\alph*)},leftmargin=2.3em,itemsep=1pt]
\item \m{\IsRow{D}{\#n-\#r}{\HT}{R}} for some type~$\HT$; and
\item for every tame character $\vartheta$ of~$D$, every auxiliary pair $V$ of type \m{(r,n)} over~\m{\ED{D}} with
\m{\mu_V=\mu_O\cdot\vartheta}, every form $y$ of~$V$ that lifts~$z$, every system $H$ of Hecke polynomials for which
$y$ is an eigenform and every homomorphism \m{S\colon\AbsGal{D}\to\GL_n(\Qlb)} attached to~$H$:
\[
\tr\normalisedRep{V}{S}(g)=\invAvatar{\vartheta}{g}\cdot\tr R(g)+\sum_{e<\#n-\#r}\cyc(g)^{-e}\qquad\text{for all }g\in\AbsGal{D}.
\]
\end{enumerate}
\end{definition}
\leanline{\lean{HasRow}}

Clause (b) says that the normalised representation of an eigenform lifting~$z$ is, on traces, the row twisted
by the avatar of~$\vartheta^{-1}$ plus the string \m{1,\cyc^{-1},\dots,\cyc^{-(\#n-\#r-1)}}.

\begin{definition}[labelled columns]\label{d:labcol}
Let $D$, $r$, $n$, $O$, $c$ be as in Definition~\ref{d:twl}, $m$ a finite set, $k$ a natural number, $\HT$ a type,
$S$ a set of finite places of~\m{\ED{D}} and \m{R\colon\AbsGal{D}\to\GL_m(\Qlb)} a homomorphism. Then
\m{\IsLabelledColumn{O}{k}{\HT}{c}{S}{R}} holds if \m{\IsColumn{D}{k}{\HT}{S}{R}} and, for every \m{i\in I_O} and every
place $\fP$ of~\m{\Qb}, $R$ is labelled by~$c_i$ at~$\fP$ for the faithful character of~$\F_{w_i}$ and
\m{f=(\#r)!}.
\end{definition}
\leanline{\lean{IsLabelledColumn}}

\subhead{3.4. Auxiliary primes.}

\begin{definition}[data at an auxiliary prime]\label{d:auxdata}
Let $E$ be a CM field as in the conventions and $q$ a natural number. \m{\AuxPrimeData{E}{q}} is the set of the
quadruples \m{l=(\fP,\gamma,\phi,h)} of a place $\fP$ of~\m{\Qb}, elements \m{\gamma,\phi\in\GE{E}} and a
homomorphism \m{h\colon\GE{E}\to\mathbb{Z}/q}. (The formal text writes \m{\mathbb{Z}/q} multiplicatively.)
\end{definition}
\leanline{\lean{AuxPrimeData}}

\begin{definition}[anticyclotomic data]\label{d:acyc}
Let $E$ be a CM field as in the conventions, $T$ a set of finite places of~$E$, $\fq$ a finite place, $q$ a natural
number and \m{l=(\fP,\gamma,\phi,h)\in\AuxPrimeData{E}{q}}. Then \m{\IsAnticyclotomicAt{T}{\fq}{l}} holds if $h$ is from
$\mathrm{U}(1)$ relative to~$T$ with modulus~\m{\{\fq\}}, \m{\Nm(\fq)\equiv1\bmod q}, $\fP$ lies over~$\fq$, $\phi$ is an
\m{\Nm(\fq)}-Frobenius at~$\fP$, \m{\fq\notin T}, and \m{\gamma\in I_{\fP}} with \m{h(\gamma)=1} and
\m{\gamma=\tau_{\fP}}.
\end{definition}
\leanline{\lean{IsAnticyclotomicAt}}

\begin{definition}[auxiliary sets]\label{d:auxil}
Let $E'$ be a CM field as in the conventions and $S$, $I$ finite sets of finite places of~$E'$. Then
\m{\IsAuxiliary{S}{I}} holds if every \m{v\in I} satisfies: \m{v\notin S}, \m{\Nm(v)} is a prime number,
\m{\cj{v}\notin I} and \m{\cj{v}\notin S}.
\end{definition}
\leanline{\lean{IsAuxiliary}}

\begin{definition}[families of values]\label{d:family}
Let $E$ be a CM field as in the conventions, $S$ a finite set of finite places of~$E$, $T$ a set,
\m{\nu=(\nu_i)_{i\in T}} a family of homomorphisms \m{\nu_i\colon\GE{E}\to\Qlb^{\times}}, and \m{\mathrm{val}} a function
that assigns an element \m{\mathrm{val}(I,\chi)\in\Qlb} to a finite set $I$ of finite places of~$E$ and a multiplicative
homomorphism \m{\chi\colon\GE{E}\to\Qlb}. Then \m{\IsFamily{S}{\nu}{\mathrm{val}}} holds if
\begin{enumerate}[label=\textup{(F\arabic*)},leftmargin=3em,itemsep=1pt]
\item every $\nu_i$ is continuous;
\item for every $I$ with \m{\IsAuxiliary{S}{I}} there are a finite set~$\alpha$, elements \m{g_a\in\GE{E}} and
\m{c_a\in\Zlb} (\m{a\in\alpha}) with \m{\mathrm{val}(I,\chi)=\sum_{a}c_a\,\chi(g_a)} for all~$\chi$;
\item for every $I$, every finite place \m{\fq\notin I} with \m{\IsAuxiliary{S}{I\cup\{\fq\}}}, every $\chi$ from
$\mathrm{U}(1)$ relative to \m{\pd{E}{\ell}} with modulus~$I$, every place $\fP$ of~\m{\Qb} over~$\fq$, every
\m{\Nm(\fq)}-Frobenius $\phi$ at~$\fP$ and every real \m{a\ge0}: if
\m{\vl\bigl(1-\nu_i(\phi)\chi(\phi)^{\pm1}\bigr)\le a} for all $i$ and both signs, then
\m{\vl(\mathrm{val}(I\cup\{\fq\},\chi))\le\vl(\mathrm{val}(I,\chi))+8a};
\item \m{\mathrm{val}(\emptyset,1)\ne0}.
\end{enumerate}
\end{definition}
\leanline{\lean{IsFamily}}

\begin{definition}[even classes]\label{d:even}
Let $E$ be a CM field as in the conventions, $\cO$ a coefficient ring, $S$ a set of finite places of~$E$ and
$\nu$, $e$ natural numbers. Then \m{\evenClasses{\cO}{S}{\nu}{e}} is the set of the cocycles
\m{c\in\Sel{\cO(-e)}{\pd{E}{\ell}\cup S\cup\cj{S}}{\emptyset}{\cO/\ell^{\nu}}} with \m{c(g^{\mathsf c})=c(g)} for all
\m{g\in\GE{E}} that vanish at every $g$ which lies in the decomposition group of a place of~\m{\Qb}
above~$\ell$ and has \m{\cyc(g)=1}.
\end{definition}
\leanline{\lean{evenClasses}}

\begin{definition}[new places]\label{d:new}
Let $E$ be a CM field as in the conventions, $\nu$ a natural number, $J$ a set, \m{w=(w_j)_{j\in J}} a family of
finite places of~$E$ and \m{\tau=(\tau_j)_j}, \m{\phi=(\phi_j)_j} families in~\m{\GE{E}}. Then
\m{\IsNewPlaces{\nu}{w}{\tau}{\phi}} holds if for every~$j$: \m{w_j\notin\pd{E}{\ell}}, \m{\cyc_{\ell^\nu}(\phi_j)=1}, and
\m{\tau_j=\tau_{\fP}} for some place $\fP$ of~\m{\Qb} over~$w_j$. For \m{k=(j,b)}, a pair of an index and a sign,
\m{\withConjInv{\tau}{k}} denotes \m{\tau_j^{[b]}}, that is $\tau_j$ or \m{(\tau_j^{\mathsf c})^{-1}}.
\end{definition}
\leanline{\lean{IsNewPlaces}\sep \lean{withConjInv}}

\begin{definition}[prime triangles]\label{d:tri}
Let $E$ be a CM field as in the conventions, $\nu$ a natural number, $J$ a finite totally ordered set,
\m{w=(w_j)_{j\in J}} finite places of~$E$, \m{\tau=(\tau_j)_j} and \m{\phi=(\phi_j)_j} families in~\m{\GE{E}},
\m{a\in\GE{E}}, $t$ a natural number, \m{\mathrm{ty}\colon J\to\{e:0<e<t\}} a monotone map, and $h$, $c$ natural
numbers. Then \m{\IsPrimeTriangle{\nu}{w}{\tau}{\phi}{a}{t}{\mathrm{ty}}{h}{c}} holds if
\begin{enumerate}[label=\textup{(T\arabic*)},leftmargin=3em,itemsep=1pt]
\item \m{\IsNewPlaces{\nu}{w}{\tau}{\phi}};
\item there is a finite set \m{s\subseteq\GE{E}} with \m{(\#J+\#s)\cdot h\le c} such that for every coefficient
ring~$\cO$ and every \m{e\in\NN}, an element of \m{\Sel{\cO(-e)}{\pd{E}{\ell}}{\emptyset}{\cO/\ell^{\nu}}} that vanishes
on~$s$ is zero;
\item for every \m{e<t}, the product \m{\prod_{k<t,\,k\ne e}\bigl(\cyc(a^{-1})^{e}-\cyc(a^{-1})^{k}\bigr)} divides
\m{\ell^{h}} in~\m{\Zl};
\item for every $\cO$ and every \m{j\in J}: if \m{z\in\evenClasses{\cO}{w(J)}{\nu}{\mathrm{ty}(j)}} satisfies, for all
\m{i<j}, \m{z(\tau_i)=0}, and \m{z(\phi_i)=0} when \m{\mathrm{ty}(i)=\mathrm{ty}(j)}, then
\m{\ell^{h}\cdot z(\tau_j)=0};
\item for every $\cO$ and every \m{0<e<t}, with \m{\delta(g)=\cyc(g)^{-e}-1\in\cO}: \m{4\,\delta(a)^2} divides
\m{\ell^{h}} in~$\cO$, and every \m{z\in\evenClasses{\cO}{\emptyset}{\nu}{e}} with \m{z(\phi_j)=0} for all $j$ with
\m{\mathrm{ty}(j)=e} satisfies \m{\delta(a)\cdot2z(g)=\delta(g)\cdot2z(a)} for all \m{g\in\GE{E}}.
\end{enumerate}
\end{definition}
\leanline{\lean{IsPrimeTriangle}}

In a prime triangle the ordered places $w_j$ serve the twists \m{0<\mathrm{ty}(j)<t} and cut down the even classes
modulo~$\ell^\nu$ one after the other, up to~$\ell^h$; the element $a$ separates the characters \m{\cyc^{-e}} up
to~$\ell^h$; and $c$ bounds the total cost.

\begin{definition}[auxiliary twists]\label{d:auxtw}
Let $E$ be a CM field as in the conventions, $S$ and $Q$ finite sets of finite places of~$E$,
\m{\psi\colon\GE{E}\to\Qlb} a multiplicative homomorphism and \m{c\in[0,\infty]}. Then
\m{\IsAuxTwist{S}{Q}{\psi}{c}} holds if \m{\IsAuxiliary{S}{Q}}, $\psi$ is from $\mathrm{U}(1)$ relative to
\m{\pd{E}{\ell}} with modulus~$Q$, and there is a family \m{\sigma=(\sigma_{\fq})_{\fq\in Q}} in~\m{\GE{E}} with
\m{\IsNewPlaces{0}{(\fq)_{\fq\in Q}}{\sigma}{1}} (each \m{\fq\in Q} is not above~$\ell$ and \m{\sigma_{\fq}} is the tame
generator at a place over~$\fq$) and \m{\sum_{\fq\in Q}\vl(\psi(\sigma_{\fq})-1)\le c}.
\end{definition}
\leanline{\lean{IsAuxTwist}}

\begin{definition}[a union with conjugates]\label{d:union}
For finite sets $S$, $J$ of finite places of a CM field as in the conventions, the set
\m{\unionWithConj{S}{J}} is the union of $S$, of $J$ and of the set of the conjugates of the places of~$J$.
\end{definition}
\leanline{\lean{unionWithConj}}

\begin{definition}[twists at a point]\label{d:point}
Let $\Gamma$ be a monoid, $A$ and $A'$ commutative rings, $n$ and $\sigma$ finite sets, \m{q\ge1} an integer,
\m{h=(h_j)_{j\in\sigma}} a family of homomorphisms \m{h_j\colon\Gamma\to\mathbb{Z}/q} and \m{y=(y_j)_{j\in\sigma}} a
family in~$A$. The \emph{point character} of $h$ and $y$ is defined under the hypothesis that \m{y_j^{q}=1} for
all~$j$: it is the homomorphism \m{\Gamma\to A} whose value at \m{g\in\Gamma} is \m{\prod_jy_j^{h_j(g)}}, the exponents
being the representatives in \m{[0,q)}. For a ring homomorphism \m{e\colon A'\to A}, a function
\m{\varrho\colon\Gamma\to M_n(A')} and \m{g\in\Gamma},
\[
\twistAtPoint{h}{y}{e}{\varrho}{g}=\Bigl(\prod_jy_j^{h_j(g)}\Bigr)\cdot e(\varrho(g))\in M_n(A),
\]
with $e$ applied to the entries. For natural numbers $l$, $M$ and a homomorphism \m{\varrho\colon\Gamma\to M_n(A)},
\m{\reduceModPow{l}{M}{\varrho}} is $\varrho$ followed by the reduction \m{M_n(A)\to M_n(A/l^{M})}.
\end{definition}
\leanline{\lean{pointChar}\sep \lean{twistAtPoint}\sep \lean{reduceModPow}}

\begin{definition}[balanced representations]\label{d:bal}
Let $E$ be a CM field as in the conventions, $A$ the ring of integers of a finite extension of~\m{\Ql}, with
valuation \m{v_A} normalised by \m{v_A(\varpi)=1} and residue field~$\kappa_A$, $n$ a finite set, $S$ a finite set of finite
places of~$E$, \m{\varrho\colon\GE{E}\to M_n(A)} a homomorphism, $L$ a set of relations on~\m{\GE{E}} with values in
\m{M_n(A)}, \m{\gamma\in\GE{E}}, $\beta$ a non-negative extended real number (a bound for lengths) and \m{M\in\NN}. Then
\m{\IsBalanced{\ell}{S}{\varrho}{L}{\gamma}{\beta}{M}} holds if
\begin{enumerate}[label=\textup{(B\arabic*)},leftmargin=3em,itemsep=1pt]
\item \m{\reduceModPow{\ell}{M}{\varrho}} is locally constant, and \m{\varrho(g)=1} for all
\m{g\in\inab{S^{\complement}}};
\item the $A$-span of \m{\varrho(\GE{E})} is \m{M_n(A)}, and \m{4\,\#n+2<\#\kappa_A};
\item \m{\gamma\in\inab{S\setminus\pd{E}{\ell}}}, and there is \m{s\in A} with
\m{\varrho(\gamma^{[-]})=s\cdot1}, \m{\varrho(\gamma)=s\cdot1} and \m{s+1,\ s-1\in A^{\times}};
\item for every \m{a\in A} and \m{b\in(A/a)^n}, the coboundary \m{g\mapsto\varrho(g)b-b} obeys~$L$;
\item there are \m{g\in\GE{E}} and \m{\alpha,\alpha'\in A} with \m{\varrho(g)=\alpha\cdot1} and
\m{\varrho(g^{[-]})=\alpha'\cdot1} such that for every $a$ dividing~\m{\ell^{M}} in~$A$:
\[
\sum_{v\in S}\min_{\fP\text{ over }v}\lambda_{\fP}(a)
+\#n\cdot v_A\bigl(\alpha-\cyc(g^{\mathsf c})\,\alpha'\bigr)\ \le\ \bigl(\#S+f_E\bigr)\cdot\len_A\bigl((A/a)^n\bigr)+\beta ,
\]
where \m{\lambda_{\fP}(a)} is the $A$-length of the module of the local cocycles of $\varrho$ and~$L$ at~$\fP$ with
values in \m{(A/a)^n}, and $f_E$ is the number of infinite places of~$E$.
\end{enumerate}
\end{definition}
\leanline{\lean{IsBalanced}}

\begin{definition}[pencils of tied primes]\label{d:pencil}
Let $E$ be a CM field as in the conventions, $A$ a commutative ring, $\mathcal{I}$ a set (of rows), $S_0$ a finite set
of finite places of~$E$, \m{U=(U_i)_{i\in\mathcal{I}}} a family of properties of elements of~\m{\GE{E}}, $W$ a subset
of~\m{\GE{E}}, \m{N\in\NN}, $\fq_0$ a finite place, $J$ a finite set of finite places, and $x$ a function that
assigns to each finite place~$v$ a pair \m{x(v)=(l_v,s_v)} of a datum
\m{l_v=(\fP_v,\gamma_v,\phi_v,h_v)\in\AuxPrimeData{E}{\ell^{N}}} and a family \m{s_v=(s_{v,i})_{i\in\mathcal{I}}} in~$A$
(the \emph{slopes}). Then \m{\IsPencil{\ell}{S_0}{U}{W}{N}{\fq_0}{J}{x}} holds if
\begin{enumerate}[label=\textup{(P\arabic*)},leftmargin=3em,itemsep=1pt]
\item \m{\IsAuxiliary{S_0}{J}} and \m{\fq_0\in J};
\item \m{\IsAnticyclotomicAt{\pd{E}{\ell}}{v}{l_v}} for every \m{v\in J};
\item for all \m{v,v'\in J} with \m{v\ne\fq_0}: \m{h_{v'}(\phi_v)=1} if \m{v'=\fq_0}, and \m{h_{v'}(\phi_v)=0}
otherwise;
\item \m{\phi_{\fq_0}\in W};
\item for every \m{v\in J} and $i$ with \m{s_{v,i}\ne0}, the property $U_i$ holds at \m{\phi_v} and at
\m{\phi_v^{[-]}}.
\end{enumerate}
\end{definition}
\leanline{\lean{IsPencil}}

\begin{definition}[ties]\label{d:ties}
Let $E$ be a CM field as in the conventions, $A$ a commutative ring, $n$ a finite set, $\mathcal{I}$ a set,
\m{q\ge1}, $J$ a finite set of finite places of~$E$, $x$ a function assigning to each finite place $v$ a pair
\m{(l_v,s_v)} with \m{l_v=(\fP_v,\gamma_v,\phi_v,h_v)\in\AuxPrimeData{E}{q}} and \m{s_v=(s_{v,i})_{i\in\mathcal{I}}}
in~$A$, and \m{i\in\mathcal{I}}. Then \m{\tieRelations{J}{x}{i}} is the set of the relations
\m{\delta_{\gamma_v^{[b]}}-s_{v,i}\,\delta_{\phi_v^{[b]}}} with values in \m{M_n(A)}, for \m{v\in J} and both signs~$b$.
A function $c$ obeys them when \m{c(\gamma_v)=s_{v,i}\,c(\phi_v)} and \m{c(\gamma_v^{[-]})=s_{v,i}\,c(\phi_v^{[-]})}
for all \m{v\in J}.
\end{definition}
\leanline{\lean{tieRelations}}

\subhead{3.5. Two further definitions.}

\begin{definition}[the tame character of an auxiliary set]\label{d:tameaux}
Let $D$ be admissible data, $I$ a finite set of finite places of~\m{\ED{D}} such that
\m{\IsAuxiliary{\badPlaces{D}{\ell}}{I}}, and $\chi$ a character of~$\mathrm{U}(1)$ of weight~$0$ whose tame part is
trivial off~$I$. The \emph{tame character of~$D$ defined by $I$ and~$\chi$} is $\chi$ with \m{Q=I}; it is a tame
character because $I$ is disjoint from its conjugate and the norms of its places are prime to \m{\gD{D}\lvl{D}}
(Lemma~\ref{l:coprime}).
\end{definition}
\leanline{\lean{tameCharOfAuxiliary}}

\begin{definition}[the lattice of a matrix representation]\label{d:matlat}
Let $\Gamma$ be a group, $n$ a finite set, $A$ a commutative \m{\Zl}-algebra and
\m{\varrho\colon\Gamma\to M_n(A)} a homomorphism. Then \m{\matrixLattice{\varrho}} is the module $A^n$ with $\Gamma$
acting through~$\varrho$, as a Galois lattice over~$A$ for~$\Gamma$.
\end{definition}
\leanline{\lean{matrixLattice}}

\sect{Cocycles, embeddings and a determinant bound}{s:first}

\begin{lemma}[the graph of a cocycle that obeys the relations]\label{l:graph}
Let $A$ be a noetherian commutative ring, $\Gamma$ a group, $X$ a finite set, $\sigma$ and $d$ finite sets,
\m{u=(u_x)_{x\in X}} a family of homomorphisms \m{u_x\colon\Gamma\to M_\sigma(A)}, \m{\varrho\colon\Gamma\to M_d(A)}
a homomorphism and \m{a\in A}, and let \m{c\in\obeyingCocycles{u}{\varrho}{a}}. Let $\Gamma$ act on
\m{\mathbb{M}=M_\sigma(A)^{X}} by \m{g\cdot(m_x)_x=(u_x(g)m_x)_x} and on \m{(A/a)^d} through~$\varrho$. Then there are
\m{k\in\NN} and a $\Gamma$-stable $A$-submodule $Y$ of \m{\mathbb{M}\times(A/a)^d} such that the first projection
\m{Y\to\mathbb{M}} is injective and
\[
\bigl((u_x(g)-1)_x+a^{k}v,\ c(g)\bigr)\in Y\qquad\text{for all }g\in\Gamma,\ v\in\mathbb{M}.
\]
\end{lemma}
\begin{proof}
Let $\Lambda_0$ be the module of the finitely supported functions \m{\theta\colon\Gamma\to A} and
\m{\mathcal{L}\colon\Lambda_0\to\mathbb{M}}, \m{\theta\mapsto\sum_g\theta(g)(u_x(g)-1)_x}. By the lemma of Artin and Rees
for the ideal \m{(a)}, the finitely generated module $\mathbb{M}$ and its submodule \m{\mathcal{L}(\Lambda_0)}, there
is $k_0$ with
\[a^{k_0+1}\mathbb{M}\cap\mathcal{L}(\Lambda_0)=a\,\bigl(a^{k_0}\mathbb{M}\cap\mathcal{L}(\Lambda_0)\bigr).\]
Put
\m{k=k_0+1} and let $Y$ be the image of
\[
\mathcal{N}\colon\Lambda_0\times\mathbb{M}\to\mathbb{M}\times(A/a)^d,\qquad
\mathcal{N}(\theta,v)=\Bigl(\mathcal{L}\theta+a^{k}v,\ \sum_g\theta(g)\,c(g)\Bigr).
\]
\emph{Stability.} For \m{g\in\Gamma} put \m{\theta^{g}=\sum_h\theta(h)(\delta_{gh}-\delta_g)}. Then
\begin{align*}
\mathcal{L}(\theta^{g})&=\sum_h\theta(h)\bigl(u(gh)-u(g)\bigr)=u(g)\,\mathcal{L}\theta,\\
\sum_h\theta(h)\bigl(c(gh)-c(g)\bigr)&=\varrho(g)\sum_h\theta(h)c(h),
\end{align*}
the second by the cocycle identity; hence \m{g\cdot\mathcal{N}(\theta,v)=\mathcal{N}(\theta^{g},u(g)v)}.
\emph{Injectivity.} If \m{\mathcal{L}\theta+a^{k}v=0}, then \m{\mathcal{L}\theta\in a^{k}\mathbb{M}\cap\mathcal{L}(\Lambda_0)},
so \m{\mathcal{L}\theta=-a\,\mathcal{L}\eta} for some $\eta$, that is \m{\mathcal{L}(a\eta+\theta)=0}. Since $c$ obeys
the relations among the \m{u_x(g)-1}, \m{\sum_g(a\eta+\theta)(g)c(g)=0}, and $a$ kills \m{(A/a)^d}; so
\m{\sum_g\theta(g)c(g)=0}. \emph{The points.} \m{\mathcal{N}(\delta_g,v)=((u(g)-1)+a^{k}v,c(g))}, with $\delta_g$ here
the scalar indicator of~$g$.
\end{proof}
\leanline{\lean{existsgraphsubrepresentation}}

So $c$ is, on the elements \m{(u(g)-1)+a^kv}, the value of a $\Gamma$-equivariant map defined on the projection
of~$Y$, and $Y$ is its graph. The third note uses this to turn a cocycle that obeys the relations of the lifts
into an extension.

\begin{lemma}[four small facts]\label{l:small1}
\textup{(a)} Let $E$, $\cO$, $n$, $u$ be as in Definition~\ref{d:zf} and \m{a\in\cO}. A function
\m{c\colon\GE{E}\to(\cO/a)^n} obeys \m{\crysRels{u}} if and only if, for every \m{v\in\pd{E}{\ell}}, its restriction to
\m{G_{\fP_v}} obeys \m{\Rel{\ZfAt{v}{u}}}.
\textup{(b)} Let $E$ be a CM field as in the conventions, $\fP$ a place of~\m{\Qb} with an $\ell$-adic chart, of
open embedding \m{j_{\fP}\colon G_{\fP}\to\Gl}, and \m{c\in\NN}. Then \m{g\in G_{\fP}} acts on \m{\cO(-c)} by the scalar
\m{\cyc(j_{\fP}(g))^{-c}}, with \m{\cyc} the cyclotomic character of~\m{\Gl}.
\textup{(c)} Let $E$ be a field, $\Gamma$ a topological group, \m{e\colon E\to\Qlb} a ring homomorphism and
\m{j\colon\Gamma\to\Gl} an open embedding whose image is the group of the elements fixing \m{e(E)} pointwise. Then
for \m{g,g'\in\Gl}: \m{\complexEmbOf{e}{g}=\complexEmbOf{e}{g'}} if and only if \m{g\,g'^{-1}\in j(\Gamma)}.
\textup{(d)} In the situation of \textup{(c)}, the index \m{[\Gl:j(\Gamma)]} is the number of the homomorphisms
\m{\complexEmbOf{e}{g}}, \m{g\in\Gl}.
\end{lemma}
\begin{proof}
(a) unfolds Definition~\ref{d:zf}: a relation extended by zero is evaluated on the restriction. (b) is the
compatibility of the chart with the cyclotomic characters. (c): \m{\iota^{-1}g^{-1}e=\iota^{-1}g'^{-1}e} means that
\m{g'g^{-1}}, hence its inverse, fixes \m{e(E)}. (d) follows from (c): the fibres of \m{g\mapsto\complexEmbOf{e}{g}} are
the cosets of~\m{j(\Gamma)}.
\end{proof}
\leanline{\lean{obeyingcrysRels}\sep \lean{tateTwistsmul}\sep \lean{complexEmbOfeqiff}\sep \lean{indexeqcardrange}}

\begin{lemma}\label{l:coeffsurj}
Let $\cO$ be a coefficient ring. Every \m{\Zl}-algebra homomorphism \m{\cO\to\Cl} is \m{\coeffEmb{\cO}{g}} for some
\m{g\in\Gl}.
\end{lemma}
\begin{proof}
$\cO$ is integral over~\m{\Zl}, so the homomorphism extends to a \m{\Zl}-algebra homomorphism \m{\Qlb\to\Cl}
(the field \m{\Cl} is algebraically closed and torsion-free over~$\cO$). Its image is the algebraic closure
of~\m{\Ql} in~\m{\Cl}, which is~\m{\Qlb}; so it is an element of~\m{\Gl} (an input of the unit on Tate--Sen theory).
\end{proof}
\leanline{\lean{coeffEmbsurjective}}

\begin{lemma}\label{l:sumemb}
Let $E$ be a CM field as in the conventions and $a$ a function from the complex embeddings of~$E$ to~$\NN$. For
\m{v\in\pd{E}{\ell}} let \m{\mathfrak{S}_v} be the set of the complex embeddings \m{\complexEmbOf{e_v|_E}{g}},
\m{g\in\Gl}, with $e_v$ the chart of Definition~\ref{d:chart}. Then
\m{\sum_{v\in\pd{E}{\ell}}\sum_{\sigma\in\mathfrak{S}_v}a(\sigma)\le\sum_{\sigma}a(\sigma)}.
\end{lemma}
\begin{proof}
The sets \m{\mathfrak{S}_v} are pairwise disjoint: if \m{\sigma=\iota^{-1}g^{-1}e_v|_E}, then an integer $x$ of~$E$
lies in~$v$ exactly when \m{|\iota\sigma(x)|<1}, because $e_v$ induces the place \m{\fP_v} over~$v$ and $g$ is an
isometry; so $\sigma$ determines~$v$. The left side is thus a sum over a subset of the embeddings.
\end{proof}
\leanline{\lean{sumcomplexEmbOfle}}

\begin{lemma}[operators that clear denominators bound a Gram determinant]\label{l:gram}
Let $V$ and $W$ be complex vector spaces, \m{\mathcal{A}\subseteq V} a set whose linear span $Z$ has finite
dimension~$N$, \m{j\colon V\to W} a linear map, \m{\eta_1,\dots,\eta_N} linear forms on~$V$, \m{x_1,\dots,x_N} vectors
of~$Z$, $\gamma$ a real number and \m{m\in\Qlb}. Suppose that for every non-zero \m{p\in\mathcal{A}} there is a linear
map \m{G\colon V\to W} such that
\textup{(i)} for every \m{v\in\mathcal{A}} there is \m{l\in\CC} with \m{G(v)=l\cdot j(v)};
\textup{(ii)} \m{e^{-\gamma}\le|\iota(l)\,m|} for every \m{l\in\CC} with \m{G(p)=l\cdot j(p)};
\textup{(iii)} \m{|m\cdot\iota(\eta_a(w))|\le1} for all $a$, $b$ and every \m{w\in Z} with \m{j(w)=G(x_b)}.
Then \m{\bigl|\iota\bigl(\det(\eta_a(x_b))_{a,b}\bigr)\bigr|\le e^{N\gamma}}.
\end{lemma}
\begin{proof}
Choose a basis \m{\omega_1,\dots,\omega_N} of~$Z$ inside~$\mathcal{A}$, for each $p$ a map $G_p$ as in the hypothesis
for~\m{\omega_p}, and scalars \m{l_{p,k}} with \m{G_p(\omega_k)=l_{p,k}\,j(\omega_k)}. For finitely many polynomials over
an algebraically closed valued field there is a point with integral coordinates at which each of them has
absolute value at least that of each of its coefficients (an input of the unit on local heights). Applied to
the linear forms \m{\sum_p\iota(l_{p,k})\,m\,X_p} it gives \m{d_1,\dots,d_N\in\Zlb} with
\m{|\lambda_k|\ge|\iota(l_{k,k})\,m|\ge e^{-\gamma}} for \m{\lambda_k=\sum_p\iota(l_{p,k})\,m\,d_p}, by~(ii). Write
\m{x_b=\sum_kP_{b,k}\omega_k}. The matrix
\[
\bigl(\iota\eta_a(\omega_k)\bigr)_{a,k}\cdot\diag(\lambda_k)\cdot\bigl(\iota P_{b,k}\bigr)_{k,b}
\]
has the entries \m{\sum_pd_p\cdot m\,\iota\eta_a(w_{p,b})}, where \m{w_{p,b}} is the element
\m{\sum_kP_{b,k}l_{p,k}\omega_k} of~$Z$. Moreover \m{j(w_{p,b})=G_p(x_b)}, so by (iii) the entries lie in~\m{\Zlb}. Hence its determinant
\m{\iota(\det(\eta_a(x_b)))\cdot\prod_k\lambda_k} has absolute value at most~$1$.
\end{proof}
\leanline{\lean{normgramdetleofoperators}}

Lemma~\ref{l:gram} is the form in which a congruence bound is used in the third note: the maps $G$ are Hecke
operators that act on lifts by scalars, $m$ is a power of~$\ell$, and the determinant is a Gram determinant.

\sect{The field of the data, avatars and lattices}{s:avatars}

\begin{lemma}[small facts on the field of the data]\label{l:small2}
Let $D$ be admissible data.
\textup{(a)} \m{\ED{D}} is a CM field as in the conventions: it is embedded in~\m{\Qb} through the cyclotomic
field of conductor \m{\ND{D}\cdot d}, and the image of every element lies in that cyclotomic field.
\textup{(b)} A finite place of~\m{\ED{D}} in \m{\pd{\ED{D}}{\ell}\cup\placesAboveG{D}} lies in \m{\badPlaces{D}{\ell}}.
\textup{(c)} \m{\cj{v}\in\placesAboveG{D}} if and only if \m{v\in\placesAboveG{D}}.
\textup{(d)} If \m{x\ne y} in~\m{\Qlb} and \m{x^N=y^N=1} for a natural number~$N$, then \m{|N|\le|x-y|}.
\textup{(e)} \m{\cyc^{-1}\colon\AbsGal{D}\to\Qlb} is continuous.
\textup{(f)} If \m{(w,\fP,s)} is a Frobenius datum of~\m{\ED{D}} (a finite place of prime norm, a place of~\m{\Qb}
over it and $s$ with $s^{-1}$ an \m{\Nm(w)}-Frobenius at~$\fP$) and \m{w\notin\pd{\ED{D}}{\ell}}, then
\m{\cyc(s)^{-1}=\Nm(w)}.
\textup{(g)} If the system of places of an auxiliary pair is that of a tame character~$\vartheta$, then each of its
chosen places $w_i$ is chosen, lies in \m{\placesAboveG{D}}, and \m{\psi_{\vartheta,w_i}=1}.
\end{lemma}
\begin{proof}
(a) is how the formal text equips \m{\ED{D}} with its embedding. (b), (c), (g) unfold the definitions; for (g),
the set \m{Q(\vartheta)} has norms prime to \m{\gD{D}}, so it contains no chosen place. (d): \m{x-y} divides
\m{\prod_{\zeta^N=1,\zeta\ne1}(1-\zeta)=N} up to a unit, since \m{x^{-1}y} is an $N$-th root of unity different
from~$1$. (e) and (f) are the continuity of the cyclotomic character and its value at a Frobenius element.
\end{proof}
\leanline{\lean{algebraMapmemcycloField}\sep \lean{placesDividingunionsubsetbadPlaces}\sep \lean{starmemplacesAboveG}\sep
\lean{normnatCastlenormsubofpoweqone}\sep \lean{continuouscycInv}\sep \lean{cycInveqofisFrob}\sep \lean{isChosenchosenPlace}}

\begin{lemma}[conjugate duality of an avatar]\label{l:avconj}
Let $D$ be admissible data, $a$ a character of~$\mathrm{U}(1)$ of~\m{\ED{D}} with a weight~$e$ and a tame
part~$\psi$, and $k$, $j$ integers. Then, for every \m{g\in\AbsGal{D}},
\[
\av(a[k,j])(g^{\mathsf c})\cdot\av(a[k,j])(g)=\cyc(g)^{2j-k}.
\]
\end{lemma}
\begin{proof}
Both sides are continuous in~$g$, and the Frobenius elements of most Frobenius data are dense (an input of the
unit on Chebotarev's theorem); so it suffices to prove the identity for \m{g=s} with \m{(w,\fP,s)} a Frobenius
datum, $w$ and \m{\cj{w}} outside~$\ell$ and outside the support of~$\psi$. Then $s^{\mathsf c}$ is the
Frobenius element of a datum at~\m{\cj{w}}, and the avatar at such an element is $\iota$ of the value of the
Hecke character at the uniformiser. The factor \m{a(x/\cj{x})} takes inverse values at \m{\varpi_w} and
\m{\varpi_{\cj w}}; \m{\upsilon(\varpi_w)\upsilon(\varpi_{\cj w})=1}; and \m{|\varpi_w|^{2(j-k/2)}=\Nm(w)^{-(2j-k)}=\cyc(s)^{2j-k}}
by Lemma~\ref{l:small2}(f).
\end{proof}
\leanline{\lean{avatarconjGalmul}}

\begin{lemma}[the label of an avatar]\label{l:avlabel}
Let $D$ be admissible data, $a$ a character of~$\mathrm{U}(1)$ of~\m{\ED{D}} with a weight~$e$ and a tame
part~$\psi$, and $k$, $j$ integers. Let $w$ be a finite place of~\m{\ED{D}}, $\fP$ a place of~\m{\Qb} and
\m{f\in\NN}. Suppose that \m{\psi_{\cj w}=1} and that $w$ is split. Then the
one-dimensional representation \m{\av(a[k,j])} is labelled at~$\fP$, for the faithful character of~$\F_w$ and~$f$,
by the polynomial \m{X-x_{\psi_w}}.
\end{lemma}
\begin{proof}
Let \m{s\in I_{\fP}} have tame exponent~$a'$; in particular $\fP$ lies over~$w$. By the local compatibility of
avatars (unit of the Galois representations), $s$ is the local Artin symbol of a unit $u$ of~$E_w$ and
\m{X-\iota(a[k,j]_w(u))} kills~$s$. The characters $\upsilon$ and \m{|\cdot|} are trivial on units, and
\m{a(u/\cj{u})=\psi_w(u\bmod w)} because \m{\psi_{\cj w}=1} and \m{\cj{w}\ne w}. By the property that defines
\m{x_{\psi_w}}, \m{\psi_w(u\bmod w)=T(x_{\psi_w})^{a'}}.
\end{proof}
\leanline{\lean{isLabelledAtavatar}}

\begin{lemma}[the normalising character at Frobenius elements]\label{l:chione}
Let $D$ be admissible data, $r$, $n$ finite sets, $O$ an auxiliary pair of type \m{(r,n)} over~\m{\ED{D}}, $w$ a
finite place of~\m{\ED{D}} and $\fP$ a place of~\m{\Qb}. Suppose that $\fP$ lies over~$w$ and that $w$ does not
contain~$\ell$. Let \m{k\in\NN} and let \m{s\in\AbsGal{D}} be such that $s^{-1}$ is an \m{\Nm(w)^{k}}-Frobenius
at~$\fP$. Then \m{\chiOneInv{O}(s)=\iota(\nu_w^{-1})^{k}}, where
\m{\nu_w=\chi(\varpi_w)\,\Nm(w)^{\#r/2}}, with $\chi$ the Hecke character \m{x\mapsto a_O(x/\cj{x})\,\upsilon(x)^{-\#r}}, is the
\emph{Frobenius value} of~$a_O$ in rank~\m{\#r} at~$w$.
\end{lemma}
\begin{proof}
$s$ is the local Artin symbol of an element of the completion at~$w$ of order~$k$, and the avatar at~$s$ is
$\iota$ of the value there of the Hecke character \m{a_O^{-1}[-\#r,0]}. That character is unramified at~$w$, so the
value is the $k$-th power of its value at~$\varpi_w$, which is \m{\chi(\varpi_w)^{-1}\Nm(w)^{-\#r/2}=\nu_w^{-1}}.
\end{proof}
\leanline{\lean{chiOneInvspec}}

\begin{lemma}[divisibility and counting at the level]\label{l:gpow}
Let $D$ be admissible data.
\textup{(a)} For \m{i\in\{0,1\}} and an integer $d$ with \m{1\le d\le\Dr{D}\tD{D}}, the number \m{\Dg{i}{D}^{d}-1}
divides \m{\ordch{D}}.
\textup{(b)} If $T$ is a finite set of finite places of~\m{\ED{D}} with \m{T\subseteq\placesAboveG{D}}, then
\m{\#T\le4\,\fD{D}}.
\textup{(c)} There are a finite place $w$ of~\m{\ED{D}} and an index~$i$ such that $w$ is chosen,
\m{w\in\placesAboveG{D}\setminus\pd{\ED{D}}{\ell}} and \m{\Nm(w)=\Dg{i}{D}}.
\end{lemma}
\begin{proof}
(a) \m{\ordch{D}} is divisible by the order of \m{\GL_d(\F_w)} for a chosen place $w$ above~\m{\Dg{i}{D}}
(Chapter~III), and that order \m{\prod_{k<d}(q^d-q^k)}, \m{q=\Nm(w)=\Dg{i}{D}}, has the factor \m{q^d-1}.
(b) A place above \m{\Dg{0}{D}} or \m{\Dg{1}{D}} is a chosen place or the conjugate of one, and there are
\m{\fD{D}} chosen places above each of the two primes. (c) Take $i$ with \m{\Dg{i}{D}\ne\ell} and a chosen place
above~\m{\Dg{i}{D}}.
\end{proof}
\leanline{\lean{gpowsubonedvdchosenGroupOrder}\sep \lean{cardleofsubsetplacesAboveG}\sep \lean{existschosennotMemplacesDividing}}

\begin{lemma}[a lattice as a conjugate]\label{l:alghom}
Let $D$ be admissible data, $n$ a finite set, $\cO$ a coefficient ring, \m{R\colon\AbsGal{D}\to\GL_n(\Qlb)} a
homomorphism and \m{u\colon\AbsGal{D}\to M_n(\cO)} a lattice of~$R$. Then there is an injective homomorphism of
$\cO$-algebras \m{F\colon M_n(\cO)\to M_n(\Qlb)} with \m{F(u(g))=R(g)} for all~$g$ and
\m{\det F(Y)=\det Y}, \m{\tr F(Y)=\tr Y} for all \m{Y\in M_n(\cO)}.
\end{lemma}
\begin{proof}
With \m{u(g)=Y_0^{-1}R(g)Y_0}, take \m{F(Y)=Y_0\,Y\,Y_0^{-1}}; conjugation preserves determinant and trace.
\end{proof}
\leanline{\lean{existsalgHom}}

\begin{lemma}[small facts used later]\label{l:small3}
Let $D$ be admissible data.
\textup{(a)} For \m{i\in\{0,1\}} and \m{1\le d\le\Dr{D}\tD{D}}: \m{\vl(\Dg{i}{D}^{d}-1)\le\crudeUnit{\ell}{D}}.
\textup{(b)} Let $u$ be a lattice of~$R$ as in Lemma~\ref{l:alghom}. If $R$ is continuous, so is~$u$; if
\m{R(g)=1} then \m{u(g)=1}; and \m{\tr u(g)=\tr R(g)} for all~$g$.
\textup{(c)} Let $O$ be an auxiliary pair of type \m{(r,n)} over~\m{\ED{D}} and $c$ a family of polynomials as in
Definition~\ref{d:twl}. If \m{\IsApt{O}{c}{\vartheta}} for a twist~$\vartheta$, then \m{\#n\le2\,(\kD{D}+\Dm{D}(\tD{D}+1))}.
The endomorphism \m{\tupleProj{O}{c}} is symmetric for the inner product of \m{\mathcal{F}(O)}, and it fixes its
image. Infinitely many finite places of~\m{\ED{D}} are split and lie outside the places of~$O$ and their
conjugates. If a subspace $Z$ of \m{\mathcal{F}(O)} is stable under the Hecke operators \m{H_w(X)} of all split
places $w$ off the places of~$O$ and their conjugates and all \m{X\in\CC}, so is its orthogonal complement.
\textup{(d)} Let $\fP$ be a place of~\m{\Qb} over a finite place $\fq$ of a CM field $E$ as in the conventions,
\m{\tau\in I_{\fP}}, and $S$ a set of finite places with \m{\fq,\cj{\fq}\in S}. Then \m{\tau^{[b]}\in\inab{S}} for both
signs~$b$.
\textup{(e)} For a tame character $\tau$ of~$D$, a finite idele $x$ of~\m{\ED{D}} and \m{g\in\AbsGal{D}} with
\m{\Art(x)=g}: \m{\invAvatar{\tau}{g}=\iota\bigl(\tau(x/\cj{x})\bigr)^{-1}}.
\textup{(f)} For a complex embedding $\sigma$ of~\m{\ED{D}}: \m{\sigma(\sqrt{-d})=\cmSign{\sigma}\sqrt{-d}}, and the
chosen embedding of the place of~$\sigma$ sends the root of unity $\varsigma$ of~\m{\ED{D}} to
\m{\sigma(\varsigma)^{\cmSign{\sigma}}}.
\end{lemma}
\begin{proof}
(a) By Lemma~\ref{l:gpow}(a), \m{\vl(\Dg{i}{D}^d-1)\le\vl(\ordch{D})}. (b) follows from Lemma~\ref{l:alghom} and
from the fact that conjugation is a homeomorphism onto its image. (c) The bound on \m{\#n} is (A3) with the
inequalities of the data. The idempotent \m{\tupleIdem{r}{c}} is central and self-adjoint (note \emph{Primitive labels},
Input~1), so its action is a symmetric projector. A place of prime norm is split, and all but finitely many
places avoid a finite set. The Hecke operator \m{H_w(X)} has an adjoint of the same kind, at~\m{\cj{w}}, so the
orthogonal complement of a stable space is stable. (d) \m{\tau^{[-]}=(\tau^{\mathsf c})^{-1}} lies in the inertia
group of the conjugate place, which is over~\m{\cj{\fq}}. (e) is the defining property of the avatar; the weight is~$0$,
so there is no factor at~$\ell$. (f) is read off from the two possible values of the chosen embedding.
\end{proof}
\leanline{\lean{ellVallecrudeUnit}\sep \lean{continuous}\sep \lean{eqone}\sep \lean{traceeq}\sep \lean{cardle}\sep
\lean{isSymmetrictupleProj}\sep \lean{tupleProjeqselfofmemrange}\sep \lean{frequentlyisSplitnotMemplaces}\sep
\lean{orthogonalmeminvtSubmodulehecke}\sep \lean{idOrConjInvmeminertiaAbove}\sep \lean{invAvatareqofartinMapeq}\sep \lean{cmSignspec}}

\begin{lemma}[the image of the projector is Hecke-stable]\label{l:tprojinv}
Let $D$ be admissible data, $r$, $n$ finite sets, $O$ an auxiliary pair of type \m{(r,n)} over~\m{\ED{D}} and $c$ a
family of polynomials as in Definition~\ref{d:twl}. Let $w$ be a finite place of~\m{\ED{D}} with
\m{\cj{w}\ne w} that lies outside the places of~$O$ and their conjugates, and \m{X\in\CC}. Then the image of
\m{\tupleProj{O}{c}} is stable under the Hecke operator \m{H_w(X)=\sum_i(-X)^i\,\Nm(w)^{i(i-1)/2}\,T_{w,i}} on
\m{\mathcal{F}(O)}.
\end{lemma}
\begin{proof}
The image of the action of an element of a group algebra is stable under every operator that commutes with the
action of the group. The Hecke operators at~$w$ are translations by elements supported at~$w$, which commute
with the translations by the lifts of elements of \m{\prod_i\GL_r(\F_{w_i})}, supported at the chosen places and
trivial at~$w$.
\end{proof}
\leanline{\lean{rangetupleProjmeminvtSubmodule}}

\sect{Apt labels occur; attached representations are conjugate self-dual}{s:apt}

\begin{lemma}[the auxiliary pairs have apt labels]\label{l:aptaux}
\textup{(a)} Let $D$ be admissible data, $\vartheta$ a twist of~$D$, $j$ a lift index \m{(1\le j\le\Dm{D})} and $c$ a
family of labels that belongs to the good tuple families of~$D$. Then \m{\IsApt{O_D(\vartheta,j)}{c}{\vartheta}}, where
\m{O_D(\vartheta,j)} is the $j$-th auxiliary pair of~$\vartheta$.
\textup{(b)} Let $D$ be admissible data, $\vartheta$ a twist of~$D$, $p$ a pair of~$D$ and $j$ such that $p$ is of rank
one for~$j$ (third note). Let \m{O_p(\vartheta)} be the auxiliary pair of the $L$-value data of $p$ and~$\vartheta$.
Then \m{\IsApt{O_p(\vartheta)}{(X-1)_i}{\vartheta}}.
\textup{(c)} In the situation of \textup{(b)}, the lift \m{\Xi_p(\vartheta)} of the trivial character, a form
of~\m{O_p(\vartheta)}, lies in the image of \m{\tupleProj{O_p(\vartheta)}{(X-1)_i}}.
\end{lemma}
\begin{proof}
(a) (A1) and (A2) hold by the construction of the auxiliary pair (its tame part is \m{\omega_D^{j}\psi_\vartheta}).
(A3): its type is \m{(r,n)} with \m{\#r=\Dr{D}} and rank difference \m{\tD{D}}, both odd. Each component of a good
family is a good tuple: every $c_i$ is a label of degree \m{\Dr{D}} none of whose roots is a root of unity of
order dividing \m{\Mom{P}{\Dmu{D}}}, and over each \m{\Dg{k}{D}} some $c_i$ is primitive; this gives (A4) and (A5). (A6): the roots of the
twisted label are the roots of~$c_i$ times an element whose \m{\Mom{P}{\Dmu{D}}}-th power is~$1$ (the name of a power
of~$\omega_D$, the twist being trivial at the chosen places by Lemma~\ref{l:small2}(g)); if $1$ were among them,
a root of~$c_i$ would be a root of unity of order dividing \m{\Mom{P}{\Dmu{D}}}.

(b) Unfold the hypothesis: the hermitian space is a line, \m{\#r=1}, the rank difference is \m{2z_p-1}, odd and at
most \m{\tD{D}}, and the tame part of~$\nu_p$ is \m{\omega_D^{j}}. \m{X-1} is a label of degree~$1$, a chosen place
above each \m{\Dg{k}{D}} carries it, and the twisted label \m{X-x} has \m{x\ne1} because \m{\omega_D^{j}} is
non-trivial at the chosen places for \m{1\le j\le\Dm{D}}.

(c) Write \m{\Xi=\Xi_p(\vartheta)} and \m{G=\prod_i\GL_1(\F_{w_i})}. For \m{s\in G} and a form~$y$,
\m{\langle\Xi,s\cdot y\rangle=\langle\Xi,y\rangle}: both sides are class sums of the theta descent of~$y$
(Chapter~VII), and the descent of \m{s\cdot y} is the descent of~$y$ translated by a lift of~$s$, which the class
sum absorbs. The idempotent of the label \m{X-1} in rank one is the average over the group twisted by the
trivial character, so \m{\tupleIdem{r}{X-1}=\#G^{-1}\sum_{s\in G}s}, $r$ being the one-element index set of the rows.
Write $\mathsf{e}$ for this element. As \m{\tupleProj{O_p(\vartheta)}{(X-1)_i}} is symmetric,
\m{\langle\mathsf{e}\,\Xi,y\rangle=\langle\Xi,\mathsf{e}\,y\rangle=\langle\Xi,y\rangle} for all~$y$, so
\m{\mathsf{e}\,\Xi=\Xi}.
\end{proof}
\leanline{\lean{isAptauxPair}\sep \lean{isAptrankone}\sep \lean{trivialLiftmemrangetupleProjrankone}}

\begin{lemma}[attached representations are conjugate self-dual]\label{l:csdatt}
Let $D$ be admissible data, $n$ a finite set, $f$ a form in rank~\m{\#n} over~\m{\ED{D}} (a function on the finite
adelic points with values in the functions on a set~$M$), $H$ a system of Hecke polynomials and
\m{\varrho\colon\AbsGal{D}\to\GL_n(\Qlb)} a homomorphism. If $f$ is an eigenform for~$H$ and $\varrho$ is attached
to~$H$, then \m{\IsConjSelfDual{\#n}{\varrho}}.
\end{lemma}
\begin{proof}
Continuity and the non-degeneracy of the trace form are part of attachment. For the trace identity let
\m{\#n=k+1} (for \m{\#n=0} there is nothing to prove) and write \m{c_i(Y)} for the coefficient of~$X^i$ in the
characteristic polynomial of a matrix~$Y$. At a split place $w$ of prime norm outside~$\ell$ at which $f$ is
spherical at $w$ and~\m{\cj{w}}, the translate by the scalar~$\varpi_w$ of the $i$-th Hecke operator at~\m{\cj{w}}
is the \m{(\#n-i)}-th Hecke operator at~$w$ (duality of the Hecke operators at a split place, an input of the
unit on the Weil representation). Applied to~$f$ and compared coefficient by coefficient this gives
\m{H_w(0)\,H_{\cj w}(k)=H_w(1)\,\Nm(w)^{k}}, where \m{H_w(i)} denotes the coefficient of~$X^i$. Now let the
Frobenius data approach~$g$. By attachment, \m{\iota H_w} tends coefficientwise to the characteristic polynomial
of \m{\varrho(g)}; the system \m{w\mapsto H_{\cj w}} has \m{\varrho\circ(\cdot)^{\mathsf c}} attached, because conjugation
is a homeomorphism of \m{\AbsGal{D}} that carries Frobenius data at~$w$ to Frobenius data at~\m{\cj{w}}; and
\m{\Nm(w)=\cyc(s)^{-1}} tends to \m{\cyc(g)^{-1}} (Lemma~\ref{l:small2}(e),(f)). By the uniqueness of limits,
\[
c_0(\varrho(g))\cdot c_k(\varrho(g^{\mathsf c}))=c_1(\varrho(g))\cdot\cyc(g)^{-k}.
\]
Since \m{\tr\varrho(g^{\mathsf c})=-c_k(\varrho(g^{\mathsf c}))} and
\m{\tr\varrho(g^{-1})=-c_1(\varrho(g))/c_0(\varrho(g))}, this is the identity of Definition~\ref{d:csd} with
\m{t=\#n}.
\end{proof}
\leanline{\lean{isConjSelfDualofattached}}

For a weight $\kappa$ and a natural number~$N$, \m{\HT_{\kappa,N}} denotes the type of the representations
attached to forms of weight~$\kappa$ in rank~$N$ (Chapter~VII): at the embedding of the CM type the multiset of the
\m{i-1-\lambda_i}, at its conjugate that of the \m{\lambda_i+N-i}, \m{1\le i\le N}, with $\lambda$ the partition
of~$\kappa$.

\begin{lemma}[local conditions pass to the conjugate place]\label{l:conjplace}
Let $D$ be admissible data, $r$ a finite set, $j$ a natural number, \m{R\colon\AbsGal{D}\to\GL_r(\Qlb)} a
homomorphism and $\fP$ a place of~\m{\Qb}. Suppose \m{\IsConjSelfDual{j}{R}} and that $R$ is of finite type at the
conjugate place~\m{\fP^{\mathsf c}}. Then
\textup{(1)} $R$ is of finite type at~$\fP$;
\textup{(2)} for every weight~$\kappa$: if $R$ has the jumps \m{\HT_{\kappa,j}} at \m{\fP^{\mathsf c}}, it has them
at~$\fP$;
\textup{(3)} for every \m{s\in I_{\fP^{\mathsf c}}} and every polynomial $Q$: if $Q$ kills $s$ on~$R$
at~\m{\fP^{\mathsf c}}, then $Q$ kills \m{(s^{-1})^{\mathsf c}} on~$R$ at~$\fP$.
\end{lemma}
\begin{proof}
Let \m{R'(g)=\cyc(g)^{-(j-1)}\,{}^{t}R(g^{-1})}, the dual of~$R$ twisted by \m{\cyc^{1-j}}. Its trace form is
non-degenerate, and by the trace identity of Definition~\ref{d:csd} it has the traces of
\m{g\mapsto R(g^{\mathsf c})}, whose span is semiprime. By the theorem of Brauer and Nesbitt in the form of the
background library there is \m{Y\in\GL_r(\Qlb)} with \m{R'(g)\,Y=Y\,R(g^{\mathsf c})} for all~$g$. Conjugation by
$\mathsf{c}$ carries the place $\fP$ to \m{\fP^{\mathsf c}}, so the local clauses of \m{R'} at~\m{\fP^{\mathsf c}}
are those of~$R$ at~$\fP$ (transport along a conjugation, an input of the unit on $\ell$-adic characters). It
remains to read the clauses of~$R'$ at \m{\fP^{\mathsf c}} from those of~$R$ there. The pairing
\m{\langle x,y\rangle={}^{t}x\,y} between $R'$ and $R$ is equivariant up to \m{\cyc^{-(j-1)}}; so along a chart
$R'$ is potentially crystalline when $R$ is, and the jump $u$ of~$R'$ has the multiplicity of the jump
\m{(j-1)-u} of~$R$. On inertia at a place prime to~$\ell$ the cyclotomic character is trivial, so $R'$ is
trivial near~$1$ where $R$ is: this is~(1). For (2): the type \m{\HT_{\kappa,j}} is carried to itself by
\m{u\mapsto(j-1)-u} together with the exchange of an embedding with its conjugate, which is what the passage
from \m{\fP^{\mathsf c}} to~$\fP$ does to the embeddings. For (3) at a place prime to~$\ell$: \m{Q(R(s))=0} gives,
after transposition and because \m{\cyc(s)=1}, \m{Q(R'(s^{-1}))=0}, hence \m{Q(R((s^{-1})^{\mathsf c}))=0}; above~$\ell$ an
inertial element in a Weil pair has degree~$0$, and the clause of the pairing gives the same conclusion.
\end{proof}
\leanline{\lean{localclausesofconjPlace}}

\begin{lemma}[a Hecke-stable space is spanned by isotypic vectors]\label{l:isotypic}
Let $D$ be admissible data, $r$, $n$ finite sets, $O$ an auxiliary pair of type \m{(r,n)} over~\m{\ED{D}}, and $Z$,
$Z'$ subspaces of~\m{\mathcal{F}(O)}. Suppose that $Z$ is stable under the Hecke operators \m{H_w(X)} of all split
finite places $w$ outside the places of~$O$ and their conjugates and all \m{X\in\CC}, and that every \m{v\in Z}
which is non-zero and lies in an isotypic space belongs to~$Z'$. Then \m{Z\subseteq Z'}.
\end{lemma}
\begin{proof}
$Z$ is finite-dimensional, fixed by the level that is trivial at the places of~$O$, and stable under all the
operators \m{T_{w,i}} at the split places off these places (they are the coefficients of the polynomial
\m{H_w(X)}). Such a space is spanned by its Hecke eigenvectors off the places of~$O$, and a non-zero eigenvector
of maximal level off these places lies in an isotypic space (two inputs of the units on local--global
compatibility and on adelic cusp forms).
\end{proof}
\leanline{\lean{leofisotypicmem}}

\sect{Auxiliary primes}{s:aux}

\begin{proof}[Proof of Theorem~A]
\emph{Step 1: units.} Choose a finite set $T_0$ of finite places of~$E$ that contains
\m{\cj{S}\cup S\cup\pd{E}{\ell}}, meets $R$ only inside that set, and is such that every finite idele is a principal
idele times an idele that is a unit outside~$T_0$ (finiteness of the class group). By Dirichlet's theorem the
group of $T_0$-units of~$E$ has a finite set $\mathfrak{s}$ of generators. For \m{u\in\mathfrak{s}} choose
\m{y_u\in\Qb} with \m{y_u^{\ell^N}=u}.

\emph{Step 2: the prime.} Let \m{\mathcal{U}} be the intersection of the kernel of \m{\cyc_{\ell^N}}, of the
stabilisers of the~$y_u$, and of the interior of~$W$. It is open, and it contains~$z$: the first two groups
contain every conjugate of an inertia element at a place $v$ of~$R$ outside \m{S\cup\cj{S}\cup\pd{E}{\ell}} (inertia at a
place prime to~$\ell$ fixes the $\ell^N$-th roots of unity and the $\ell^N$-th roots of the units at~$v$, and the
$T_0$-units are units at~$v$ because \m{v\notin T_0}), hence the normal subgroup generated by these inertia
groups. By Chebotarev's theorem (Input~1) there are a finite place \m{\fq\notin T_0\cup R\cup\cj{R}} of prime
norm, a place $\fP$ of~\m{\Qb} over~$\fq$ and an \m{\Nm(\fq)}-Frobenius \m{\sigma_0\in\mathcal{U}} at~$\fP$ such that
every \m{\Nm(\fq)}-Frobenius at~$\fP$ lies in~\m{\mathcal{U}}. Since \m{\sigma_0} fixes the $\ell^N$-th roots of unity,
\m{\Nm(\fq)=1+\ell^Nk} for some~$k$.

\emph{Step 3: a character of the ideles.} Let \m{\gamma=\tau_{\fP}}, the local Artin symbol of a unit \m{u_0} of
the completion \m{E_{\fq}} whose residue generates \m{\F_{\fq}^{\times}}. Let \m{\rho_0\colon E_{\fq}^{\times}\to\F_{\fq}^{\times}}
send $t$ to the residue of \m{t\varpi_{\fq}^{-\operatorname{ord}t}}, and let \m{f_0\colon E_{\fq}^{\times}\to\mathbb{Z}/\ell^N} be
\m{\rho_0} followed by the homomorphism that sends the residue of~$u_0$ to~$1$ (it exists because $\ell^N$ divides
\m{\Nm(\fq)-1}). If \m{\rho_0(t^k)=1}, then \m{f_0(t)=0}. For \m{u\in\mathfrak{s}}: $u$ is a unit at~$\fq$, the root
$y_u$ is a unit at~$\fP$, and since $\sigma_0$ fixes $y_u$ and is a Frobenius element,
\m{y_u\equiv y_u^{\Nm(\fq)}} modulo~$\fP$; so \m{u^{k}=y_u^{\Nm(\fq)-1}\equiv1} modulo~$\fq$ and \m{f_0(u)=0}. Hence
the homomorphism \m{a\mapsto f_0(a_{\fq})} on the ideles that are units outside~$T_0$ is trivial on the $T_0$-units,
and by the choice of~$T_0$ it extends uniquely to a homomorphism $\eta$ on the finite ideles that is trivial on
the principal ideles. It is continuous, with open kernel.

\emph{Step 4: the Galois character.} By the existence theorem of class field theory (Input~2) there is a
continuous homomorphism \m{h\colon\GE{E}\to\mathbb{Z}/\ell^N} with \m{h(g)=\eta(x/\cj{x})} whenever \m{\Art(x)=g}. As
\m{\Nm(\fq)} is prime, \m{\cj{\fq}\ne\fq}, so the conjugate of the local idele~$u_0$ is a unit idele supported
at~\m{\cj{\fq}} and \m{h(\gamma)=f_0(u_0)=1}. Put \m{\phi=\sigma_0\gamma^{a}} with $a$ the representative of
\m{-h(\sigma_0)}; it is again an \m{\Nm(\fq)}-Frobenius at~$\fP$, so \m{\phi\in\mathcal{U}\subseteq W}, and
\m{h(\phi)=0}.

\emph{Step 5: the clauses.} Let \m{T\subseteq S\cup\pd{E}{\ell}} and \m{l=(\fP,\gamma,\phi,h)}. The character $\eta$ is
trivial on the ideles that are units outside \m{T\cup\cj{T}\subseteq T_0} and congruent to~$1$ at~$\fq$ (they have
\m{f_0(a_{\fq})=0}), and $h$ is locally constant: $h$ is from $\mathrm{U}(1)$ relative to~$T$ with
modulus~\m{\{\fq\}}. The other clauses of Definition~\ref{d:acyc} were established: \m{\Nm(\fq)\equiv1\bmod\ell^N},
$\fP$ over~$\fq$, $\phi$ a Frobenius element, \m{\fq\notin T_0\supseteq T}, and \m{\gamma=\tau_{\fP}\in I_{\fP}} with
\m{h(\gamma)=1}. Finally \m{\fq\notin R\cup\cj{R}} and \m{\Nm(\fq)} is prime.
\end{proof}

\inp{Input 1 (Chebotarev's theorem, existence form; background library, unit of Chebotarev's theorem)}{Let
$E$ be a number field embedded in~\m{\Qb}, \m{\mathcal{U}} a non-empty open subset of~\m{\GE{E}} and $T$ a finite set
of finite places. Then there are a finite place \m{\fq\notin T} of prime norm, a place $\fP$ of~\m{\Qb} over it and
an \m{\Nm(\fq)}-Frobenius at~$\fP$ in~\m{\mathcal{U}}, such that every \m{\Nm(\fq)}-Frobenius at~$\fP$ lies
in~\m{\mathcal{U}}. \emph{Reference} (as located by the second reading of this write-up): J.~S.~Milne,
\emph{Class Field Theory}, v4.03, Chapter~VIII, §7, Theorem~7.4. \emph{Relation:} an existence consequence of the
density theorem for finite Galois extensions, with a chosen place above the prime; no density is used.}

\inp{Input 2 (existence theorem of class field theory; background library, unit of the global Artin map)}{Let
$E$ be a totally complex number field embedded in~\m{\Qb} and $\theta$ a continuous homomorphism from the finite
ideles of~$E$ to a group with a $T_1$ topology, trivial on the principal ideles. Then there is a continuous
homomorphism $\mu$ on~\m{\GE{E}} with \m{\mu(g)=\theta(x)} whenever \m{\Art(x)=g}. \emph{Reference} (second reading):
Milne, \emph{Class Field Theory}, v4.03, Chapter~V, §5, Theorems 5.3 and~5.5 and Remark~5.7(a), pp.~178--180.
\emph{Relation:} a special case (totally complex fields, finite ideles): the factorisation follows from the
description of the kernel of the reciprocity map.}

\begin{lemma}\label{l:four}
Let $E$ be a CM field as in the conventions and \m{a\in\GE{E}}. Then $4$ divides \m{\cyc(a)-1} in~\m{\Zl}.
\end{lemma}
\begin{proof}
If \m{\ell=2}: $E$ contains a primitive fourth root of unity, which $a$ fixes, so \m{\cyc(a)\equiv1\bmod4}. If
$\ell$ is odd, $4$ is a unit of~\m{\Zl}.
\end{proof}
\leanline{\lean{fourdvdcycCharsubone}}

\begin{lemma}[the cyclotomic character at new places]\label{l:newcyc}
Let $E$ be a CM field as in the conventions, $J$ a set, \m{\nu\in\NN}, and \m{w=(w_j)_{j\in J}},
\m{\tau=(\tau_j)_j}, \m{\phi=(\phi_j)_j} with \m{\IsNewPlaces{\nu}{w}{\tau}{\phi}}. Let $\cO$ be a coefficient ring and
\m{k=(j,b)} a pair of an index and a sign. Then \m{\cyc(\withConjInv{\tau}{k})=1}, and for every \m{e\in\NN} the
element \m{\ell^{\nu}} divides \m{\cyc(\withConjInv{\phi}{k})^{-e}-1} in~$\cO$.
\end{lemma}
\begin{proof}
$\tau_j$ and \m{(\tau_j^{\mathsf c})^{-1}} lie in inertia groups at places prime to~$\ell$, where \m{\cyc} is trivial.
$\phi_j$ fixes the $\ell^\nu$-th roots of unity, and so do its conjugate by~$\mathsf{c}$, its inverse and its
powers; so \m{\cyc} of these elements is \m{\equiv1\bmod\ell^\nu}.
\end{proof}
\leanline{\lean{cycInvwithConjInv}}

\begin{lemma}[vanishing at the tame generators makes a cocycle unramified]\label{l:newsel}
Let $E$ be a CM field as in the conventions, $\Lambda$ a commutative ring, \m{e,\nu\in\NN}, $J$ a set and
\m{w=(w_j)_{j\in J}}, $\tau$, $\phi$ with \m{\IsNewPlaces{\nu}{w}{\tau}{\phi}}. Let $n$ be a finite set,
\m{\varrho\colon\GE{E}\to M_n(\Lambda)} a homomorphism, $S$ and $T$ sets of finite places of~$E$, $L$ a set of
relations and \m{c\colon\GE{E}\to(\Lambda/\ell^{e})^n} a function. Suppose that
\m{c\in\Sel{\varrho}{T}{L}{\Lambda/\ell^{e}}}, that \m{c(\withConjInv{\tau}{k})=0} for every \m{k=(j,b)}, that
\m{T\subseteq S\cup w(J)\cup\cj{w(J)}}, and that \m{\varrho(g)=1} for every \m{g\in\inab{S^{\complement}}} with
\m{\cyc(g)=1}. Then \m{c\in\Sel{\varrho}{S}{L}{\Lambda/\ell^{e}}}.
\end{lemma}
\begin{proof}
Only the vanishing of $c$ on the inertia above the places of \m{T\setminus S} is to be shown. Such a place is
$w_j$ or \m{\cj{w_j}} for some~$j$, it is prime to~$\ell$, and $\varrho$ is trivial on its inertia. For the tame
generator at a place over $w_j$ and a module killed by a power of~$\ell$, a Selmer cocycle that kills the tame
generator on the appropriate side, $\tau_j$ or \m{(\tau_j^{\mathsf c})^{-1}}, is unramified at the corresponding
place (an input of the unit on tame cocycles: wild inertia is pro-prime-to-$\ell$ and tame inertia is
topologically generated by the tame generator). Both sides vanish by hypothesis.
\end{proof}
\leanline{\lean{memselmerModuleoftameeqzero}}

\begin{lemma}[Selmer modules are finitely presented, uniformly in the coefficients]\label{l:corep}
Let $E$ be a number field embedded in~\m{\Qb}, $n$ a finite set, $\Lambda$ a commutative ring, \m{m\ge1} an
integer, \m{\varrho\colon\GE{E}\to M_n(\Lambda)} a function, $W$ a neighbourhood of~$1$ in~\m{\GE{E}}, $S$ a finite set of
finite places of~$E$ and $L$ a set of relations. Then there are a finite set~$\mathcal{K}$, a family
\m{(g_k,i_k)_{k\in\mathcal{K}}} in \m{\GE{E}\times n} and a $\Lambda$-submodule $\mathcal{R}$ of~$\Lambda^{\mathcal{K}}$
containing \m{m\Lambda^{\mathcal{K}}} such that, for every $\Lambda$-module $B$ with \m{mB=0} on which $W$ acts
trivially (\m{\varrho(g)b=b} for \m{g\in W}, \m{b\in B^n}), the map \m{c\mapsto(c(g_k)_{i_k})_{k}} is a bijection from
\m{\Sel{\varrho}{S}{L}{B}} onto the set of the \m{y\in B^{\mathcal{K}}} with \m{\sum_k\rho_ky_k=0} for all
\m{\rho\in\mathcal{R}}.
\end{lemma}
\begin{proof}
This is an outline of the formal proof.
Choose an open subgroup \m{U\subseteq W}. The continuous cocycles of~\m{\GE{E}} that are unramified outside~$S$,
with values in the finite module of the \m{\mathbb{Z}/m}-valued functions on \m{\GE{E}/U}, are finite in number
(Hermite--Minkowski, in the form of the unit on Selmer groups); so there is an open subgroup $V$ on which all
of them vanish. Take \m{\mathcal{K}=(\GE{E}/V)\times n}, with $g_k$ a representative of the coset, and let $\mathcal{R}$
be generated by \m{m\Lambda^{\mathcal{K}}} and by the vectors obtained by writing the relations of
\m{\Coc{\varrho}}, the relations $\delta_g$ for \m{g\in\inab{S^{\complement}}} and the relations of~$L$ in the
coordinates. A \m{y\in B^{\mathcal{K}}} killed by~$\mathcal{R}$ is a function on \m{\GE{E}/V}, locally constant, that
obeys all these relations: a Selmer cocycle. Conversely a Selmer cocycle $c$ with values in~$B^n$ is constant
on the cosets of~$V$: for a homomorphism \m{\chi\colon B\to\mathbb{Z}/m} and an index~$i$, the function
\m{g\mapsto\bigl(xU\mapsto\chi((c(x)-c(g^{-1}x))_i)\bigr)} is a continuous cocycle of the kind above (because $W$ acts
trivially), so it vanishes on~$V$; and the homomorphisms $\chi$ separate the elements of~$B$.
\end{proof}
\leanline{\lean{existscorepresentingcoords}}

\sect{Families of interpolated values: proof of Theorem D}{s:families}

For a pair $p$ of~$D$ and a tame character~$\vartheta$, \m{y_v(p,\vartheta)} denotes the value at the uniformiser of
a finite place $v$ of the Hecke character \m{x\mapsto(\nu_p\vartheta)(x/\cj{x})}, and \m{\Nm(Q(\vartheta))} the product of
the norms of the places of~\m{Q(\vartheta)}. A function $F$ on the tame characters of~$D$ is an \emph{integral
measure at the modulus of~$\vartheta$} if there are finitely many points $u_a$ of the norm-one adeles and complex
coefficients $c_a$, integral at every prime that is prime to \m{\Nm(Q(\vartheta))}, with
\m{F(\vartheta')=\sum_ac_a\vartheta'(u_a)} for all tame $\vartheta'$ with \m{Q(\vartheta')=Q(\vartheta)} (unit of the exchange
identity).

\begin{lemma}[an integral measure for a ramified pair]\label{l:measure}
Let $D$ be admissible data and $p$ a pair of~$D$ that is ramified. Then there are a complex function $F$ on the
tame characters of~$D$, a set \m{K\subseteq\{0,1\}} and \m{m\in\NN} such that
\begin{enumerate}[label=\textup{(\arabic*)},leftmargin=2.3em,itemsep=1pt]
\item \m{F(1)\ne0}, and $F$ is an integral measure at the modulus of every tame character~$\vartheta$;
\item for tame characters $\vartheta$, $\vartheta'$ and a finite place \m{v\notin Q(\vartheta)} with
\m{Q(\vartheta')=Q(\vartheta)\cup\{v\}}, the same character of~$\mathrm{U}(1)$ and the same tame part:
\[F(\vartheta')=F(\vartheta)\cdot\prod_{k\in K}\ \prod_{\varepsilon=\pm1}\bigl(1-y_v(p,\vartheta)^{\varepsilon}\,\Nm(v)^{k-z_p}\bigr);\]
\item for every twist $\vartheta$: \m{F(\vartheta)=F(1)\cdot\Nm(Q(\vartheta))^{m}\cdot\varrho_p(\vartheta,1)^{m+1}}.
\end{enumerate}
\end{lemma}
\begin{proof}
\emph{The case \m{z_p\ge2}.} Take for $F$ the integral algebraic part of~$p$: the normalised squared norm
\m{\ordch{D}\cdot\mathrm{den}^2\cdot\|\Xi_p(\vartheta)\|^2/\Nm(Q(\vartheta))^{\#n-1}} of the lift of the trivial character,
with \m{K=\{0\}} and \m{m=0}. It is non-zero at~$1$, an integral measure at each modulus, and
\m{\varrho_p(\vartheta,1)=F(\vartheta)/F(1)} (three inputs of the unit of the exchange identity), which is (1) and~(3).
For (2): \m{F(\vartheta)=A\cdot L(z_p,(\nu_p\vartheta)^{-1})} with a constant \m{A\ne0} and the Euler factors removed at
the chosen places, at \m{Q(\vartheta)} and at their conjugates. Passing from $\vartheta$ to $\vartheta'$ removes the Euler
factors at $v$ and~\m{\cj{v}}; at \m{z_p>1} this multiplies the Euler product by
\m{\prod_{w\in\{v,\cj{v}\}}(1-a_w\Nm(w)^{-z_p})}, and the coefficients at $v$ and~\m{\cj{v}} are \m{y_v} and~\m{y_v^{-1}}.

\emph{The case \m{z_p=1}.} Let \m{\Lambda(\vartheta)} be the period, at the level of the chosen places and
of~\m{Q(\vartheta)}, of the determinant kernel of rank~$2$ for the character \m{\nu_p\vartheta\,\upsilon_{\mathrm{U}(1)}}
(unit of the columns), let \m{N'\ne0} be an integer that makes the theta kernels of rank~$2$ integral at the
primes prime to the modulus (unit of the rows; it is not the parameter~$N_0$), and put
\m{F(\vartheta)=|\ordch{D}\,N'\,\Lambda(\vartheta)|^2},
\m{K=\{0,1\}}, \m{m=1}. \m{\Lambda(\vartheta)} is a sum, over the classes and over \m{\GL_2} of the residue ring of
the modulus, of values of the kernel, each divided by the order of a stabiliser, which divides~\m{\ordch{D}}; the
values times~$N'$ are integral at the primes prime to~\m{\Nm(Q(\vartheta))} after the point is moved into the
integral points above the prime. So $F$, the product of such a sum with its conjugate, is an integral measure.
Adding $v$ to the modulus multiplies the period by the Euler factor \m{\prod_{k<2}(1-\Nm(v)^{k}X)} of rank~$2$
at~$v$, hence $F$ by the product in~(2). By the class-sum identity of the determinant kernel,
\m{|\Lambda(\vartheta)|^{2}} is a constant times the product of $L$-values of \m{(\nu_p\vartheta)^{-1}} (here $p$ ramified is
used), and that product is \m{\Nm(Q(\vartheta))^{z_p^2}} times the product for~\m{\nu_p^{-1}} times \m{\varrho_p(\vartheta,1)^2};
with \m{z_p=1} this is (3), and \m{F(1)\ne0}. The formal proof of this case is a computation of some forty
lines; the steps above are its outline.
\end{proof}
\leanline{\lean{existsisIntegralMeasure}}

\begin{lemma}\label{l:coprime}
Let $D$ be admissible data and $I$ a finite set of finite places of~\m{\ED{D}} with
\m{\IsAuxiliary{\badPlaces{D}{\ell}}{I}}. Then $\ell$ is prime to \m{\prod_{w\in I}\Nm(w)}, and so is
\m{\gD{D}\,\lvl{D}}.
\end{lemma}
\begin{proof}
Each \m{\Nm(w)}, \m{w\in I}, is a prime number, and it does not divide \m{\ell\,\gD{D}\,\lvl{D}}: otherwise $w$ would
contain that number and lie in \m{\badPlaces{D}{\ell}}.
\end{proof}
\leanline{\lean{coprimeqofisAuxiliary}}

\begin{lemma}[characters from $\mathrm{U}(1)$ are avatars of tame characters]\label{l:tamechar}
Let $D$ be admissible data, $I$ a finite set of finite places of~\m{\ED{D}} with
\m{\IsAuxiliary{\badPlaces{D}{\ell}}{I}}, and \m{\psi\colon\AbsGal{D}\to\Qlb} a multiplicative homomorphism that is from
$\mathrm{U}(1)$ relative to \m{\pd{\ED{D}}{\ell}} with modulus~$I$. Then there is a tame character $\tau$ of~$D$ with
\m{Q(\tau)=I} and \m{\psi(g)=\invAvatar{\tau}{g}} for all \m{g\in\AbsGal{D}}.
\end{lemma}
\begin{proof}
Let $\eta$ be the idele character of~$\psi$. Then \m{u\mapsto\iota^{-1}\eta(u_{\mathrm f})}, with \m{u_{\mathrm f}} the
finite part of a norm-one adele~$u$ read as a finite idele, is trivial on the norm-one elements of~\m{\ED{D}} and
on the integral norm-one adeles that are congruent to~$1$ at the places of~$I$: a character of the class group
of~$\mathrm{U}(1)$ of tame level~$I$. Such a character is a character $\mu$ of~$\mathrm{U}(1)$ of weight~$0$ with
tame part trivial off~$I$ (an input of the unit on Hecke characters). Let $\tau$ be the tame character defined by $I$ and~$\mu^{-1}$
(Definition~\ref{d:tameaux}). By Lemma~\ref{l:small3}(e), \m{\invAvatar{\tau}{g}=\iota(\mu(x/\cj{x}))=\eta(x/\cj{x})=\psi(g)}
when \m{\Art(x)=g}, and the Artin map is surjective.
\end{proof}
\leanline{\lean{existstameCharinvAvatareq}}

\begin{proof}[Proof of Theorem~D]
Let $F$, $K$, $m$ be as in Lemma~\ref{l:measure}. For each $I$ with \m{\IsAuxiliary{\badPlaces{D}{\ell}}{I}} fix the
presentation \m{F(\vartheta')=\sum_ac_{I,a}\vartheta'(u_{I,a})} at the modulus~$I$, and for each point a finite idele
\m{y_{I,a}} with \m{u_{I,a}} of finite part \m{y_{I,a}/\cj{y_{I,a}}}; for other~$I$ let the sum be empty. Define
\[
\mathrm{val}(I,\chi)=\sum_a\iota(c_{I,a})\cdot\chi\bigl(g_{I,a}\bigr)^{-s},\qquad g_{I,a}\ \text{an element with}\ \Art(y_{I,a})=g_{I,a},
\]
and let \m{\nu_{(k,\varepsilon)}} be the avatar of \m{\nu_p^{\varepsilon}[0,-(z_p-k)]}. By Lemma~\ref{l:small3}(e),
\begin{equation}\label{e:val}
\mathrm{val}(I,\invAvatar{\tau}{\cdot})=\iota\bigl(F(\tau^{s})\bigr)\quad\text{for tame }\tau\text{ with }Q(\tau)=I,\qquad \mathrm{val}(\emptyset,1)=\iota(F(1)).
\end{equation}
(F1) holds because avatars are continuous; (F2) because $\ell$ is prime to the norms of~$I$
(Lemma~\ref{l:coprime}), so the \m{\iota(c_{I,a})} are integral; (F4) by \eqref{e:val} and \m{F(1)\ne0}.

(F3). Let $I$, $\fq$, $\chi$, $\fP$, $\phi$, $a$ be as there. By Lemma~\ref{l:tamechar}, \m{\chi=\invAvatar{\tau}{\cdot}}
with \m{Q(\tau)=I}, and by \eqref{e:val} and Lemma~\ref{l:measure}(2),
\m{\mathrm{val}(I\cup\{\fq\},\chi)=\mathrm{val}(I,\chi)\cdot\prod_{k\in K}\prod_{\varepsilon}\iota\bigl(1-y_{\fq}^{\varepsilon}\Nm(\fq)^{k-z_p}\bigr)}.
Each factor is \m{1-\nu_{(k,-\varepsilon)}(\phi)\chi(\phi)^{s\varepsilon}}, by the value of an avatar at a Frobenius element
of a place where the character is unramified; it has valuation at most~$a$. There are at most four factors, and
the valuation of a product is at most the sum: the increase is at most \m{4a\le8a}.

(1). The tame generator \m{\tau_w} is the Artin symbol of a unit whose residue $\gamma_0$ generates \m{\F_w^{\times}};
as $w$ is chosen and \m{\cj{w}} is not, \m{\nu_i(\tau_w)=\iota(\psi_{p,w}(\gamma_0))^{\pm1}}. This is a root of unity
of order dividing \m{\Nm(w)-1}, different from~$1$ because $p$ is ramified. By Lemma~\ref{l:small2}(d),
\m{|\Nm(w)-1|\le|\nu_i(\tau_w)-1|}.

(2). Let $\tau$ be a twist with \m{Q(\tau)} auxiliary. By \eqref{e:val} and Lemma~\ref{l:measure}(3) (the character
\m{\tau^{s}} is again a twist),
\m{\mathrm{val}(Q(\tau),\invAvatar{\tau}{\cdot})=\iota(F(1))\cdot\Nm(Q(\tau))^{m}\cdot\iota(\varrho_p(\tau,s))^{m+1}}. The norm is prime
to~$\ell$ (Lemma~\ref{l:coprime}) and \m{|\iota F(1)|\le1}, so the hypothesis gives
\m{\ell^{-a}\le|\iota\varrho_p(\tau,s)|^{m+1}}, that is \m{(m+1)\,\vl(\iota\varrho_p(\tau,s))\le a}.
\end{proof}

\sect{Labelled columns, rows and the base level: proofs of Theorems B and C}{s:rows}

\begin{lemma}\label{l:rootsunit}
Let \m{x\in\Qlb} be a root of unity and $a$, $b$ natural numbers such that \m{(x^{a})^{\ell^{u}}\ne(x^{b})^{\ell^{u}}}
for all \m{u\in\NN}. Then \m{|x^{a}-x^{b}|=1}.
\end{lemma}
\begin{proof}
Write the order of~$x$ as \m{\ell^{u}k} with $k$ prime to~$\ell$. Then \m{(x^{a})^{\ell^{u}}} and \m{(x^{b})^{\ell^{u}}}
are distinct $k$-th roots of unity, so their difference has absolute value at least \m{|k|=1}
(Lemma~\ref{l:small2}(d)); and it is a multiple of \m{x^{a}-x^{b}} in~\m{\Zlb}.
\end{proof}
\leanline{\lean{normpowsubpoweqone}}

\begin{lemma}[the eigenvalues of inertia above $g$]\label{l:inertia}
Let $D$ be admissible data, $r$, $n$, $m$ finite sets, $O$ an auxiliary pair of type \m{(r,n)} over~\m{\ED{D}},
$\vartheta$ a twist and $c$ a family of polynomials with \m{\IsApt{O}{c}{\vartheta}}. Let \m{l\in\{0,1\}}, let $v$ be a
finite place of~\m{\ED{D}} containing \m{\Dg{l}{D}}, $\fP$ a place of~\m{\Qb} over~$v$, and \m{t\in\NN}. Then there is
\m{\gamma\in I_{\fP}}, equal to \m{\tau_{\fP}} if $v$ is chosen, such that for every \m{s\in I_{\fP}} there is a multiset
$\mathfrak{M}$ of elements of~\m{\Qlb} with \m{\#r+t} elements such that
\begin{enumerate}[label=\textup{(\arabic*)},leftmargin=2.3em,itemsep=1pt]
\item every \m{y\in\mathfrak{M}} satisfies \m{y^{\Dg{l}{D}^{(\#r)!}-1}=1} and \m{y^{\Dg{l}{D}^{d}-1}=1} for some
\m{1\le d\le\Dr{D}}, and \m{y\ne1} if \m{t=0} and \m{s=\gamma};
\item for every homomorphism \m{R\colon\AbsGal{D}\to\GL_m(\Qlb)}, every \m{k\in\NN}, every type $\HT$ and every set
$S$ of finite places with \m{\IsLabelledColumn{O}{k}{\HT}{(\twistedLabel{O}{c}{i}\cdot(X-1)^{t})_i}{S}{R}}, the
polynomial \m{\prod_{y\in\mathfrak{M}}(X-y)} kills $s$ on~$R$ at~$\fP$ and kills \m{(s^{-1})^{\mathsf c}} on~$R$
at~\m{\fP^{\mathsf c}}.
\end{enumerate}
\end{lemma}
\begin{proof}
If $v$ is not chosen, \m{\cj{v}} is; the statement for \m{\cj{v}} and \m{\fP^{\mathsf c}} gives the statement
for~$v$ on applying \m{s\mapsto(s^{-1})^{\mathsf c}}. So let \m{v=w_i} be a chosen place of~$O$ (A1), with faithful
character~$T$. By (A4) the twisted label is a label of degree~\m{\#r}. Take \m{\gamma=\tau_{\fP}}. Every
\m{s\in I_{\fP}} has a tame exponent~$a$ for \m{f=(\#r)!}, which for \m{s=\gamma} can be taken prime to
\m{\Nm(v)^{f}-1} (an input of the unit on local--global compatibility). Let $\mathfrak{M}$ be the multiset of the
\m{\iota(T(x)^{a})} over the roots $x$ of \m{\twistedLabel{O}{c}{i}\cdot(X-1)^{t}}. Then (2) at~$\fP$ is the definition
of a labelled column, and at \m{\fP^{\mathsf c}} it follows from Lemma~\ref{l:conjplace}(3), a column being
conjugate self-dual and of finite type at~$\fP$. For (1): a root $x$ of a label of degree~\m{\#r} over~\m{\F_v}
satisfies \m{x^{q^{f}-1}=1} and \m{x^{q^{d}-1}=1} for the degree \m{d\le\#r\le\Dr{D}} of its minimal polynomial,
\m{q=\Nm(v)=\Dg{l}{D}}; for \m{x=1} take \m{d=1}. If \m{t=0} and \m{s=\gamma}, then \m{\iota T(x)^{a}=1} would give
\m{x=1} ($a$ prime to the order, $T$ and $\iota$ injective), against~(A6).
\end{proof}
\leanline{\lean{existsinertiaeigenvalues}}

\begin{proof}[Proof of Theorem~B]
(W1) is given. Choose $l$ with \m{\Dg{l}{D}\ne\ell} and, by (A5), a chosen place $w_i$ above \m{\Dg{l}{D}} with $c_i$
primitive and \m{\#r=\Dr{D}}, or \m{c_i=X-1} and \m{\#r=1}. Put \m{q=\Nm(w_i)=\Dg{l}{D}}, let $\fP$ be a place over~$w_i$
(it is prime to~$\ell$), $\phi$ a $q$-Frobenius at~$\fP$, and \m{s\in I_{\fP}} of tame exponent~$1$. By the labelled
clause, \m{\prod_{\xi}(X-\iota T(\xi))}, over the roots $\xi$ of \m{\twistedLabel{O}{c}{i}}, kills \m{R(s)}; so the
eigenvalues of \m{R(s)} are among the \m{\iota T(\xi)}, and every polynomial vanishing at these numbers kills
\m{R(s)} ($R$ is of finite type at~$\fP$).

\emph{The roots.} Recall that \m{\Mom{P}{\Dmu{D}}=(N_0d)^{\Dmu{D}+6}}, that \m{9} and \m{\Mom{P}{\Dmu{D}}^2}
divide \m{q-1} (third note; Chapter~III), and let $j$ be the index of~(A2). In the primitive case the twisted
label is again primitive (its roots are those of~$c_i$ times a root of unity of order dividing
\m{\Mom{P}{\Dmu{D}}}; note \emph{Primitive labels}, the
lemma on the scaling of roots): its roots are \m{\xi^{q^{a}}}, \m{a<\#r}, for a generator $\xi$ of the
multiplicative group of the extension of degree~\m{\#r} of~\m{\F_{w_i}}. They satisfy \m{\xi^{q^{\#r}}=\xi}; the
\m{\xi^{q^{a}\ell^{u}}}, \m{a<\#r}, stay distinct for every~$u$ (\m{\#r=\Dr{D}} is a power of~$3$ and \m{9\mid q-1};
an input of the unit on labels); and all have the same power \m{y=\xi^{e}}, \m{e=(q^{\#r}-1)/(q-1)}, their norm,
which has order \m{q-1}. In the rank-one case the only root is the scaling element \m{y=\xi}, which names the
character \m{\omega_{D,w_i}^{-j}}, and \m{e=1}. In both cases \m{y^{x}=1} implies that \m{\Mom{P}{\Dmu{D}}} divides
\m{j\cdot x}: in the first because \m{\Mom{P}{\Dmu{D}}^2\mid q-1}, in the second because \m{\omega_{D,w_i}} has
order \m{\Mom{P}{\Dmu{D}}} (so that the order of~$y$ is that number divided by its greatest common divisor
with~$j$).

\emph{(W3).} Let $\cO$ contain the eigenvalues of~\m{R(s)} and let $u$ be a lattice of~$R$. Tame inertia gives
\m{R(\phi s\phi^{-1})=R(s)^{q}} for the arithmetic $q$-Frobenius~$\phi$ (only the fact that conjugation by~$\phi$
permutes the eigenvalues by the $q$-th power is used); \m{R(s)^{q^{\#r}}=R(s)}; and \m{R(s)} has an eigenvalue \m{\zeta=\iota T(\xi)} whose
conjugates \m{\zeta^{q^{a}}}, \m{a<\#r}, are distinct modulo the maximal ideal, by Lemma~\ref{l:rootsunit} and the
distinctness above. Under these conditions \m{u(\AbsGal{D})} spans \m{M_r(\cO)} and \m{u(\phi^{-\#r})} is a scalar
(an input of the background library on residual irreducibility); and \m{\cyc(\phi^{\#r})=q^{\#r}=\Dg{l}{D}^{\#r}}.
(Take \m{g=\phi^{\#r}}: its inverse acts by a scalar, hence so does~$g$.)

\emph{(W4).} \m{\gamma=s^{e}} lies in \m{\inab{\placesAboveG{D}\setminus\pd{\ED{D}}{\ell}}} and acts by the scalar
\m{\iota T(y)}, since \m{X^{e}-\iota T(y)} vanishes at all \m{\iota T(\xi^{q^a})}. For \m{1\le j\le4\#r+2} and
\m{a\in\NN}, \m{|(\iota T(y)^{\ell^{a}})^{j}-1|=1} by Lemma~\ref{l:rootsunit}, provided \m{y^{j\ell^{a+u}}\ne1} for
all~$u$. Suppose \m{y^{j'\ell^{a+u}}=1} with \m{1\le j'\le4\#r+2}. Then \m{\Mom{P}{\Dmu{D}}} divides
\m{j\,j'\ell^{a+u}}. Let \m{p_0\ne\ell} be a prime dividing \m{N_0d}; one of $2$ and $5$ will do, as both
divide~$N_0$. Then \m{p_0^{\,\Dmu{D}}} divides \m{j\,j'}, so that
\[2^{\Dmu{D}}\le j\,j'\le\Dm{D}\,(4\Dr{D}+2),\]
which one of the inequalities of the data excludes (Chapter~III).

\emph{(W2).} Let \m{v\in\placesAboveG{D}\setminus\pd{\ED{D}}{\ell}} contain \m{\Dg{i'}{D}} and let $\fP'$ be a place
over~$v$. Take $\gamma$ and the multisets of Lemma~\ref{l:inertia} with \m{t=0}, and
\m{N=\Dg{i'}{D}^{(\#r)!}-1}, a unit at~$\fP'$. Then \m{R(s'^{N})=1} for \m{s'\in I_{\fP'}}, because \m{X^{N}-1} vanishes
at the eigenvalue candidates. The matrix \m{R(\gamma^{[b]})} has \m{\#r} eigenvalues $x$, each \m{\ne1} with
\m{x^{\Dg{i'}{D}^{d}-1}=1} for some \m{1\le d\le\Dr{D}} (at $\fP'$ or at the conjugate place, according to~$b$); so
\m{\vl(x-1)\le\vl(\Dg{i'}{D}^{d}-1)\le\crudeUnit{\ell}{D}} by Lemma~\ref{l:small2}(d) and Lemma~\ref{l:small3}(a), and
\m{\vl(\det(R(\gamma^{[b]})-1))\le\#r\cdot\crudeUnit{\ell}{D}}.
\end{proof}

\begin{lemma}[the quadratic relation on the inertia at the twist primes]\label{l:quad}
Let $D$ be admissible data, $r$ a finite set, $n$ a non-empty finite set, $O$ an auxiliary pair of type
\m{(r,n)} over~\m{\ED{D}}, $\vartheta$ a twist and $c$ a family of polynomials with \m{\IsApt{O}{c}{\vartheta}}; suppose
\m{\IsAuxiliary{\badPlaces{D}{\ell}}{Q(\vartheta)}}. Let $f$ be a form of the weight of~$O$ on the standard definite
space of rank~\m{\#n}, $H$ a system of Hecke polynomials and \m{R\colon\AbsGal{D}\to\GL_n(\Qlb)} a homomorphism.
Suppose that $f$ is an eigenform for~$H$, that $R$ is attached to~$H$, and that at each twist place $w$ of~$O$
the form $f$ has the parahoric type with characters \m{(1,\psi_{O,w}^{-1})}. Then for every
\m{g\in\inab{Q(\vartheta)\cup\cj{Q(\vartheta)}}},
\[
\bigl(\normalisedRep{O}{R}(g)-\invAvatar{\vartheta}{g}\bigr)\bigl(\normalisedRep{O}{R}(g)-1\bigr)=0.
\]
\end{lemma}
\begin{proof}
Let \m{g=s\in I_{\fP}} with $\fP$ over \m{w\in Q(\vartheta)\cup\cj{Q(\vartheta)}}; $w$ is prime to~$\ell$. As
\m{\chiOneInv{O}(s)=1} (Lemma~\ref{l:chione} with \m{k=0}), it suffices that $R$ is of finite type at~$\fP$ and that
\m{(X-\invAvatar{\vartheta}{s})(X-1)} kills $s$ on~$R$ at~$\fP$. If \m{w\notin Q(\vartheta)}, apply the case of \m{\cj{w}}
to \m{(s^{-1})^{\mathsf c}} and transfer it by Lemmas \ref{l:csdatt} and~\ref{l:conjplace}, using
\m{\invAvatar{\vartheta}{s^{\mathsf c}}\cdot\invAvatar{\vartheta}{s}=1} (Lemma~\ref{l:avconj} with \m{k=j=0}). So let $w$ be
a twist place of~$O$ and $a$ a tame exponent of~$s$. A parahoric type gives a product label (Frobenius
reciprocity), here with the characters $1$ and \m{\psi_{\vartheta,w}^{-1}} (the tame part of~$O$ at~$w$ is that
of~$\vartheta$, since $\omega_D$ is trivial off the chosen places); and a representation attached to an eigenform
with a product label is of finite type and labelled accordingly (two inputs of the unit on local--global
compatibility): the polynomial with the roots $1$ and \m{\iota T(x_{\psi_{\vartheta,w}^{-1}})^{a}} kills~$s$. By
Lemma~\ref{l:avlabel} the second root is \m{\invAvatar{\vartheta}{s}}.
\end{proof}
\leanline{\lean{normalisedRepquadraticoninertia}}

\begin{lemma}[twisting a labelled column by an avatar]\label{l:twistav}
Let $D$ be admissible data, $r$, $n$, $m$ finite sets, $O$ an auxiliary pair of type \m{(r,n)} over~\m{\ED{D}},
$c$ and $c'$ families of polynomials as in Definition~\ref{d:twl}, $\vartheta$ a twist, \m{k\in\NN}, $\HT$ a type, $S$ a
set of finite places and \m{R\colon\AbsGal{D}\to\GL_m(\Qlb)} a homomorphism with
\m{\IsLabelledColumn{O}{k}{\HT}{c}{S}{R}}. Let $a$ be a character of~$\mathrm{U}(1)$ with a weight $e$ and a tame
part~$\psi$, let $k'$, $j$ be integers and $b$ a natural number with \m{(b-1)+(2j-k')=k-1} (both differences
\m{b-1} and \m{k-1} truncated at~$0$). Suppose \m{\IsApt{O}{c'}{\vartheta}} and that \m{\psi_v=1} for every finite place $v$
that is not chosen. Then
\[
\IsLabelledColumn{O}{b}{\HT'}{\tilde{c}}{S}{\av(a[k',j])\cdot R},
\]
where \m{\HT'(\sigma)} is \m{\HT(\sigma)} shifted by minus the infinity type of \m{a[k',j]} at~$\sigma$, and
\m{\tilde{c}_i} is $c_i$ with its roots multiplied by \m{x_{\psi_{w_i}}}.
\end{lemma}
\begin{proof}
The twist of a continuous representation with non-degenerate trace form by a continuous character is of the
same kind, and by Lemma~\ref{l:avconj} the trace identity for \m{\av(a[k',j])R} holds with the exponent
\m{(k-1)-(2j-k')=b-1}. At a place $\fP$ over~$v$ as in Definition~\ref{d:column}, the avatar is of finite type, has
its infinity type as jumps, and is good where the character is unramified, that is at the places that are not
chosen (the hypothesis on~$\psi$); these clauses and the jumps add under twisting (an input of the unit on
$\ell$-adic characters). For the labels: if \m{X-\alpha} kills $s$ on the avatar and \m{Q(\alpha X)} kills $s$ on~$R$,
then $Q$ kills $s$ on the twist; by Lemma~\ref{l:avlabel}, \m{\alpha=\iota T(x_{\psi_{w_i}})^{a'}} for $s$ of tame
exponent~$a'$ (the chosen place $w_i$ is split and its conjugate is not chosen), which scales the roots of~$c_i$
by~\m{x_{\psi_{w_i}}}.
\end{proof}
\leanline{\lean{twistavatar}}

\begin{lemma}[product labels give labelled columns]\label{l:prodlabel}
Let $D$ be admissible data, $r$, $n$ finite sets, $m$ a non-empty finite set, $O$ an auxiliary pair of type
\m{(r,n)} over~\m{\ED{D}}, $c$ and $c'$ families of polynomials as in Definition~\ref{d:twl}, $\vartheta$ a twist with
\m{\IsApt{O}{c'}{\vartheta}}, $S$ a set of finite places, $\kappa$ a weight, $f$ a form of weight~$\kappa$ on the standard
definite space of rank~\m{\#m}, $H$ a system of Hecke polynomials and \m{R\colon\AbsGal{D}\to\GL_m(\Qlb)}. Suppose
that $f$ is an eigenform for~$H$ and $R$ is attached to~$H$; that $f$ is spherical at every finite place $v$
which lies neither in \m{S\setminus(\pd{\ED{D}}{\ell}\cup\placesAboveG{D})} nor among the chosen places of~$O$ and
their conjugates; and that for a natural number $N$ and block sizes \m{t_1,\dots,t_{N}}, at each
chosen place $w_i$ the form $f$ has the product label of~$c_i$, of degree~\m{\#r}, with trivial characters on the
blocks of sizes~$t_\bullet$. Then
\[
\IsLabelledColumn{O}{\#m}{\HT_{\kappa,\#m}}{(c_i\cdot(X-1)^{\#m-\#r})_i}{S}{R}.
\]
\end{lemma}
\begin{proof}
$R$ is conjugate self-dual by Lemma~\ref{l:csdatt}. Where $f$ is spherical, $R$ is good with the jumps of the
weight (Input~3 below); at a chosen place it is of finite type with these jumps and labelled by the product
label (the input used in Lemma~\ref{l:quad}), the trivial characters contributing the factor
\m{(X-1)^{\sum t_\bullet}} with \m{\#r+\sum t_\bullet=\#m}; at the conjugates of the chosen places the clauses follow
from Lemma~\ref{l:conjplace}.
\end{proof}
\leanline{\lean{isLabelledColumnofhasProductLabel}}

\inp{Input 3 (Galois representations of eigenforms on definite unitary groups; background library, unit of the Galois representations)}{Every
eigenform $f$ on the standard definite unitary group of rank~$N$ over a CM field as in the conventions has an attached
$\ell$-adic representation. Each attached representation is, at a place where $f$ is spherical, unramified (prime
to~$\ell$) or crystalline (above~$\ell$) with the jumps of the weight of~$f$, and the Hecke polynomial kills the
geometric Frobenius element. \emph{References} (second reading): G.~Chenevier, M.~Harris, \emph{Construction of
automorphic Galois representations~II}, Theorem~3.2.3; J.-P.~Labesse, \emph{Changement de base CM et séries
discrètes}, Corollaire~5.3; M.~Harris, K.-W.~Lan, R.~Taylor, J.~Thorne, \emph{On the rigid cohomology of certain
Shimura varieties}, Theorem~A. \emph{Relation:} a weaker form of the published attachment: closeness of
characteristic polynomials to Hecke polynomials at most Frobenius elements, in place of equality at all
unramified split places.}

\begin{proof}[Proof of Theorem~C]
This is an outline of the formal proof: it unfolds definitions of the units on forms, and the identities
between adelic sums are left to the formal text.
Let $G$ be the group \m{\prod_i\GL_r(\F_{w_i})} and $\pi$ the action on forms on the standard space of rank~\m{\#n}
by right translation by the adelic points supported at the places of~$O$. The group algebra of~$G$ acts on the
forms of the level and character of~$O$ (through lifts, twisted by~$\psi_O\circ\det$), and for a form $y$
of~$O$ the action of~\m{\tupleIdem{r}{c}} on its smooth form is the smooth form of \m{\tupleProj{O}{c}(y)}.

\emph{The base level.} Its places are those of~$O$; its space is the image of~\m{\tupleIdem{r}{c}} in the smooth forms of that
level. It is stable under the operators that commute with the translations at these places, the level is open,
and the stabilisers have order dividing \m{\ordch{D}} because \m{\#n\le2(\kD{D}+\Dm{D}(\tD{D}+1))}
(Lemma~\ref{l:small3}(c) and an input of Chapter~III).

\emph{(1)} For $y$ in the image of the projector, \m{\tupleProj{O}{c}(y)=y}, so its smooth form is fixed by~\m{\tupleIdem{r}{c}}.

\emph{(2)} \m{\#G} divides \m{\ordch{D}} (Chapter~III; \m{\#r\le\Dr{D}}), and \m{\ordch{D}\cdot\tupleIdem{r}{c}} is a sum over
\m{g\in G} of an integer times \m{\#G\cdot\tupleIdem{r}{c}(g)} times~$g$. The number \m{\#G\cdot\tupleIdem{r}{c}(g)} is a value of a character
of~$G$, an algebraic integer; the scalar by which the lift of $g$ is twisted is a root of unity; and a translate
of an $\ell$-integral form by an element of the integral level is $\ell$-integral.

\emph{(3)} Let $x$ be an eigensystem of level~$u$ and $f$ a non-zero form in it; it is the smooth form of a
form of~$O$ in the image of~\m{\tupleIdem{r}{c}}, an eigenform, and \m{\varrho_x} is attached to its Hecke system. It is
spherical outside \m{Q(\vartheta)\cup\cj{Q(\vartheta)}\cup u(K)\cup\cj{u(K)}} and the chosen places and their conjugates. At
a chosen place $w_i$ its type under \m{\GL_r(\F_{w_i})} is cut out by the $i$-th factor of~\m{\tupleIdem{r}{c}}, twisted by
\m{\psi_{O,w_i}}: this is a parahoric type, which by Frobenius reciprocity gives the product label of
\m{\twistedLabel{O}{c}{i}} with one block of size \m{\#n-\#r} and trivial character (here (A6) is used: the label
must be prime to \m{X-1}). By Lemma~\ref{l:prodlabel}, \m{\varrho_x} is a labelled column with jumps
\m{\HT_{\kappa,\#n}}, exponent~\m{\#n} and the labels \m{\twistedLabel{O}{c}{i}(X-1)^{\#n-\#r}}. The character
\m{\chiOneInv{O}} is the avatar of \m{a_O^{-1}[-\#r,0]}, whose tame part is trivial; so Lemma~\ref{l:twistav}, with
\m{k'=-\#r}, \m{j=0} and \m{b=\#n-\#r}, shows that \m{R=\normalisedRep{O}{\varrho_x}} is a labelled column with
exponent \m{\#n-\#r}, the same labels and the shifted type, which is the type~$\HT$ of the theorem; it depends
on~$O$ only. At a twist place $w$ the forms of~$O$ transform under the parahoric subgroup by the character
\m{\psi_{O,w}^{-1}} of the determinant of the block indexed by~$r$; this is the parahoric type of
Lemma~\ref{l:quad}, which gives the quadratic relation.
\end{proof}

\sect{Strings of roots of Hecke polynomials}{s:strings}

\begin{lemma}\label{l:intmult}
Let \m{F_1,F_1',F_2,F_3} be polynomials over~\m{\Qlb}. Suppose that $F_2$ is monic, that all coefficients of~$F_3$
have absolute value at most~$1$, that all coefficients of \m{F_1F_2-F_3} and of \m{F_3-F_1'F_2} have absolute value
less than~$1$, and that \m{|F_1'(1)|\ge1}. Then there is a polynomial $e$ over~\m{\Zlb} such that $F_1$ divides $e$
in~\m{\Qlb[X]} and \m{e(1)=1}.
\end{lemma}
\begin{proof}
Division by a monic polynomial does not increase the maximum of the absolute values of the coefficients. So
\m{F_1-F_1'}, the quotient of \m{(F_1F_2-F_3)+(F_3-F_1'F_2)} by~$F_2$, has coefficients of absolute value less
than~$1$, and \m{|F_1(1)|\ge1}; likewise $F_1$ has coefficients of absolute value at most~$1$. Take
\m{e=F_1/F_1(1)}.
\end{proof}
\leanline{\lean{existsintegralmultipleevalone}}

\begin{lemma}[no eigenvalue $1$ near one of the elements $\delta_i$]\label{l:noeig}
Let $D$ be admissible data, $n$ a finite set, \m{T\colon\AbsGal{D}\to M_n(\Qlb)} a continuous function,
\m{(\delta_i)_{i\in\NN}} a sequence in~\m{\AbsGal{D}}, \m{y\in\Qlb} and \m{Y\in M_n(\Qlb)}. Suppose that
\m{T(\delta_i)=y^{i}\,Y} for all~$i$ and that \m{|y^{j}-1|=1} for \m{1\le j\le4\#n+2}. Then there is an $i$ such
that \m{|\det(1-T(\phi))|\ge1} for all $\phi$ in a neighbourhood of~$\delta_i$.
\end{lemma}
\begin{proof}
Otherwise, by the continuity of $T$ and of the determinant, \m{|\det(1-y^{i}Y)|<1} for \m{i=0,\dots,\#n}. Now
\m{\det(1-y^{i}Y)=\prod_\alpha(1-y^{i}\alpha)} over the eigenvalues $\alpha$ of~$Y$ with multiplicity, so for each such
$i$ some eigenvalue $\alpha_i$ has \m{|1-y^{i}\alpha_i|<1}. There are at most \m{\#n} eigenvalues, so
\m{\alpha_i=\alpha_j} for some \m{i<j\le\#n}; call it~$\alpha$. Then \m{|y^{i}\alpha|=|y^{j}\alpha|=1}, hence \m{|y|=|\alpha|=1}, and
\m{|y^{j-i}-1|=|y^{j}\alpha-y^{i}\alpha|<1}, against the hypothesis for the exponent \m{j-i\le\#n}.
\end{proof}
\leanline{\lean{existseventuallynoeigenvalueone}}

\begin{lemma}[the string of roots]\label{l:string}
Let $D$ be admissible data, $r$ a finite set, $n$ a non-empty finite set, $O$ an auxiliary pair of type
\m{(r,n)} over~\m{\ED{D}}, $x$ an isotypic space of forms of the weight of~$O$ on the standard definite space of
rank~\m{\#n} and \m{f\in x}, \m{f\ne0}. Suppose that for every \m{g\in\AbsGal{D}} the matrix
\m{\normalisedRep{O}{\varrho_x}(g)} has the eigenvalue~$1$. Let $w$ be a finite place of~\m{\ED{D}} that is split and
does not contain~$\ell$, and $S$ a polynomial such that $f$ has the Hecke polynomial $S$ at~$w$. Then, for every
\m{i<\#n-\#r}, the number \m{\nu_w\,\Nm(w)^{i}} is a root of~$S$, with \m{\nu_w} as in Lemma~\ref{l:chione}.
\end{lemma}
\begin{proof}
Let $s$ be a geometric Frobenius element at~\m{\fP_w} (the inverse of an \m{\Nm(w)}-Frobenius). By Input~3
the polynomial \m{\iota S} kills \m{\varrho_x(s)}, and by Lemma~\ref{l:chione} with \m{k=1},
\m{\chiOneInv{O}(s)=\iota(\nu_w)^{-1}}. So the eigenvalue $1$ of \m{\normalisedRep{O}{\varrho_x}(s)} is an eigenvalue
\m{\iota(\nu_w)} of \m{\varrho_x(s)}, and $\nu_w$ is a root of~$S$. The local
component of $f$ at~$w$ generates a unitary unramified representation of \m{\GL_{\#n}}, and for the Satake
parameters of such a representation a root $c$ of the Hecke polynomial with \m{|c|^2\,\Nm(w)^{k}=\Nm(w)^{\#n}} brings
the roots \m{c,c\,\Nm(w),\dots,c\,\Nm(w)^{k-1}} (an input of the unit on unramified representations of~\m{\GL_N}); here
\m{k=\#n-\#r} by the absolute value of~\m{\nu_w}.
\end{proof}
\leanline{\lean{isRootheckestring}}

\begin{lemma}[the Hecke polynomial is the string times an integral factor]\label{l:heckestring}
Let $D$ be admissible data, $r$ a finite set, $n$ a non-empty finite set, $O$ an auxiliary pair of type
\m{(r,n)} over~\m{\ED{D}}, $x$ an isotypic space of forms of the weight of~$O$ on the standard definite space of
rank~\m{\#n}, and $f$ a form with \m{f\in x}, \m{f\ne0}, such that every \m{\normalisedRep{O}{\varrho_x}(g)},
\m{g\in\AbsGal{D}}, has the eigenvalue~$1$. Suppose moreover that all coefficients of the
characteristic polynomials of the matrices \m{\normalisedRep{O}{\varrho_x}(g)} have absolute value at most~$1$. Then
for most Frobenius data \m{(w,\fP,\phi)} the following holds: for every polynomial $S$ over~\m{\Qlb} with
\m{|S(1)|\ge1} such that the characteristic polynomial of \m{\normalisedRep{O}{\varrho_x}(\phi)} is congruent to
\m{S\cdot\prod_{i<\#n-\#r}(X-\cyc(\phi)^{-i})} coefficientwise modulo the maximal ideal, there is a complex polynomial
$e$ with $\ell$-integral coefficients and \m{e(1)=1} such that the Hecke polynomial \m{H_x(w)} of~$x$ at~$w$
satisfies
\[
H_x(w)(d_0X)=\prod_{i<\#n-\#r}\bigl(X-\Nm(w)^{i}\bigr)\cdot(\text{a divisor of }e),\qquad d_0=\iota^{-1}\bigl(\chiOneInv{O}(\phi)\bigr)^{-1}.
\]
\end{lemma}
\begin{proof}
For most Frobenius data, $w$ is split and prime to~$\ell$, $f$ has the Hecke polynomial \m{H_x(w)} at~$w$, and
the characteristic polynomial of \m{\varrho_x(\phi)} is within \m{1/2} of \m{\iota H_x(w)} (attachment). By
Lemma~\ref{l:string} and Lemma~\ref{l:chione}, \m{H_x(w)(d_0X)} is divisible by the string
\m{\prod_{i<\#n-\#r}(X-\Nm(w)^{i})}; let $F_1$ be $\iota$ of the quotient, suitably scaled. By
Lemma~\ref{l:small2}(f) the string is \m{\prod(X-\cyc(\phi)^{-i})}, monic. Lemma~\ref{l:intmult}, with $F_2$ the
string, $F_3$ the characteristic polynomial of the normalised representation and \m{F_1'=S}, gives~$e$.
\end{proof}
\leanline{\lean{eventuallyheckePolyeqstringmul}}

\sect{The loss at one place above $\ell$: proof of Theorem E}{s:loss}

In this section $\varsigma$ is the root of unity of order \m{\ND{D}} that generates \m{\ED{D}} over
\m{\mathbb{Q}(\sqrt{-d})}; the infinite places of~\m{\ED{D}} form a group, identified with
\m{(\mathbb{Z}/\ND{D})^{\times}}, and for a divisor $m$ of \m{\ND{D}}, \m{\operatorname{cl}_m} is the reduction of a place
modulo~$m$ (Chapter~III); \m{[\sigma]} is the infinite place of a complex embedding~$\sigma$; \m{\thick{D}} is the set
of thick places.

\begin{lemma}\label{l:fixemb}
Let $D$ be admissible data, \m{e\colon\ED{D}\to\Qlb} a ring homomorphism and \m{g\in\Gl}. If
\m{g(\zeta)=\zeta} for every \m{\zeta\in\Qlb} with \m{\zeta^{\ND{D}d}=1}, then \m{g(e(x))=e(x)} for every
\m{x\in\ED{D}}.
\end{lemma}
\begin{proof}
\m{\ED{D}} lies in the cyclotomic field of conductor \m{\ND{D}d} (Lemma~\ref{l:small2}(a)), $e$ extends to it, and
$g$ fixes the image of its generator.
\end{proof}
\leanline{\lean{fixesemboffixesrootsOfUnity}}

\begin{lemma}[the class of a ratio of two places]\label{l:ratio}
Let $D$ be admissible data, \m{e\colon\ED{D}\to\Qlb} a ring homomorphism, \m{g,g'\in\Gl} and \m{a,n\in\NN}. Suppose
that \m{(g'g^{-1})\bigl(e(\varsigma)^{\ell^{a}}\bigr)=\bigl(e(\varsigma)^{\ell^{a}}\bigr)^{\ell^{n}}}. Let $y$ be an infinite
place of~\m{\ED{D}} such that
\[[\complexEmbOf{e}{g'}]\cdot y=[\complexEmbOf{e}{g}],\]
and $m$ a natural number that divides \m{\ND{D}} and is prime to~\m{\ell^{a}}. Then
\[
\operatorname{cl}_m(y)=\cmSign{\complexEmbOf{e}{g}}\cdot\cmSign{\complexEmbOf{e}{g'}}\cdot\ell^{n}\quad\text{in }(\mathbb{Z}/m)^{\times}.
\]
\end{lemma}
\begin{proof}
Write \m{\sigma=\complexEmbOf{e}{g}}, \m{\sigma'=\complexEmbOf{e}{g'}}. The chosen embedding of the place
\m{[\sigma']\,y} sends $\varsigma$ to the \m{\operatorname{cl}(y)}-th power of its image under the chosen embedding
of~\m{[\sigma']}, and the chosen embedding of a place \m{[\sigma]} sends $\varsigma$ to \m{\sigma(\varsigma)^{\cmSign{\sigma}}}
(Lemma~\ref{l:small3}(f)). Applying \m{g'\iota}, raising to the power \m{\ell^{a}} and inserting the hypothesis, one
finds that \m{e(\varsigma)} to the power \m{\ell^{a}\bigl(\cmSign{\sigma}\cmSign{\sigma'}\operatorname{cl}(y)-\ell^{n}\bigr)}
is~$1$. As \m{e(\varsigma)} has order \m{\ND{D}}, the number \m{\ND{D}} divides that exponent, and $m$, prime
to~\m{\ell^{a}}, divides the second factor.
\end{proof}
\leanline{\lean{placeClassratio}}

\begin{lemma}[spacing of thick places]\label{l:spacing}
Let $D$ be admissible data, \m{e\colon\ED{D}\to\Qlb} a ring homomorphism, \m{g,g',g_1\in\Gl} and
\m{a,i,n\in\NN}. Suppose that \m{[\complexEmbOf{e}{g}]} and \m{[\complexEmbOf{e}{g'}]} are thick, that
\m{(g'g^{-1})(e(\varsigma)^{\ell^{a}})=(e(\varsigma)^{\ell^{a}})^{\ell^{n}}} and
\m{(g'g_1^{-1})(e(\varsigma)^{\ell^{a}})=(e(\varsigma)^{\ell^{a}})^{\ell^{i}}}, and that \m{\ell\le\kD{D}\tD{D}},
\m{0<i}, \m{i<n} and \m{n\le27\,\Dmu{D}}. Then \m{[\complexEmbOf{e}{g_1}]} is not thick.
\end{lemma}
\begin{proof}
Suppose it is. Squaring the identity of Lemma~\ref{l:ratio} removes the signs: the ratios of the three places
have the classes \m{\ell^{2n}} and \m{\ell^{2i}} modulo every divisor of~\m{\ND{D}} that is prime to~\m{\ell^{a}}. If
\m{\ell=2}, put \m{k_3=\lfloor\log_3\Dmu{D}\rfloor+5}. The divisor \m{3^{k_3}} of
\m{\ND{D}=N_0\cdot3^{k_3}\cdot2^{\Dmu{D}}} is prime to~$2$, and the thick places are \m{\equiv-1} modulo~\m{3^{k_3}}
(Chapter~III: they are \m{\equiv-1} modulo the divisors of \m{N_0\cdot3^{k_3}}, not modulo other powers of~$3$); so
\m{4^{i}\equiv1} modulo~\m{3^{k_3}}. As \m{v_3(4-1)=1}, the order of~$4$ modulo~\m{3^{k_3}} is
\m{3^{k_3-1}=27\cdot3^{\lfloor\log_3\Dmu{D}\rfloor+1}}, which is greater than \m{27\Dmu{D}\ge n>i}: a contradiction. If $\ell$ is
odd, take a power of~$2$ dividing~\m{\ND{D}}: the three-term condition on thick places (Chapter~III) gives
\m{\ell^{2i}\equiv1} modulo \m{2^{\Dmu{D}+1}}, hence \m{2^{\Dmu{D}}\le(\ell+1)\cdot2i}, against
\m{\ell\le\kD{D}\tD{D}}, \m{i<27\Dmu{D}} and the inequalities of the data.
\end{proof}
\leanline{\lean{notthickofshortarc}}

\begin{lemma}[thick places have one sign when $\ell$ divides the conductor]\label{l:cmthick}
Let $D$ be admissible data, \m{e\colon\ED{D}\to\Qlb} a ring homomorphism and \m{g,g'\in\Gl}. Suppose that
\m{[\complexEmbOf{e}{g}]} and \m{[\complexEmbOf{e}{g'}]} are thick and that $\ell$ divides \m{\ND{D}\,d}. Then
\m{\cmSign{\complexEmbOf{e}{g}}=\cmSign{\complexEmbOf{e}{g'}}}.
\end{lemma}
\begin{proof}
Suppose the signs differ. If \m{\ell\mid\ND{D}}: put $c_0=(1+\sqrt{-d})/2$, so that $c_0(1-c_0)=(d+1)/4$, which
$\ell$ divides (Chapter~III). The two embeddings send \m{\sqrt{-d}} to opposite values, so \m{g\,g'^{-1}} exchanges
\m{e(c_0)} and \m{1-e(c_0)}; it is an isometry, so both have the same absolute value, their product has absolute
value less than~$1$, and their sum is~$1$: a contradiction with the ultrametric inequality. If \m{\ell\nmid\ND{D}},
then \m{\ell=d}, and Lemma~\ref{l:ratio} with \m{a=0} and \m{m=N_0} gives the class \m{-d^{n}} for the ratio of two
thick places, whose class modulo $N_0$ is~$1$; but no power of~$d$ is \m{\equiv-1} modulo~$N_0$ (Chapter~III).
\end{proof}
\leanline{\lean{cmSigneqofthick}}

\begin{lemma}[the rank of the crystalline cocycles]\label{l:rank}
Let $\cO$ be a coefficient ring, $\Gamma$ a topological group over~\m{\Gl} with homomorphism~$j$, $L$ a Galois lattice
over~$\cO$ for~$\Gamma$, $E$ a CM field as in the conventions and \m{e\colon E\to\Qlb} a ring homomorphism. Suppose
that $j$ is an open embedding whose image is the group of the elements fixing \m{e(E)} pointwise, that $L$ is
potentially crystalline, and that $\HT$ is a type such that for every \m{g\in\Gl} the lattice $L$ has the jumps
\m{\HT(\complexEmbOf{e}{g})} at \m{\coeffEmb{\cO}{g}}. Then
\[
\rk_{\cO}\Zf(L)=\rk_{\cO}L+\sum_{\sigma}\#\{u\in\HT(\sigma):u<0\},
\]
the sum over the complex embeddings $\sigma$ of the form \m{\complexEmbOf{e}{g}}, \m{g\in\Gl}.
\end{lemma}
\begin{proof}
By the rank formula of Bloch and Kato over~\m{\Zl} (Input~4),
\[\rk_{\Zl}\Zf(L)=\rk_{\Zl}L+[\Gl:j(\Gamma)]\cdot\#\{\text{negative jumps over }\Zl\}.\]
The jumps over~\m{\Zl} are the union
of the jumps at the \m{[\cO:\Zl]} embeddings of~$\cO$, each of which is \m{\coeffEmb{\cO}{g}} for some~$g$
(Lemma~\ref{l:coeffsurj}); the index is the number of the embeddings \m{\complexEmbOf{e}{g}} (Lemma~\ref{l:small1}(d));
and each such complex embedding is read by the same number of embeddings of~$\cO$. Dividing by \m{[\cO:\Zl]}
gives the formula.
\end{proof}
\leanline{\lean{finrankcrysCocycleseqaddnegjumps}}

\inp{Input 4 (the rank formula of Bloch and Kato; background library, unit of the Bloch--Kato condition)}{For a
group $\Gamma$ over~\m{\Gl} with homomorphism~$j$ and a potentially crystalline lattice $T$, finite free
over~\m{\Zl}, with Hodge--Tate jumps $\mathfrak{M}$ over~\m{\Zl}: \m{\Zf(T)} is finitely generated over~\m{\Zl} of rank
\m{\rk T+[\Gl:j(\Gamma)]\cdot\#\{u\in\mathfrak{M}:u<0\}}. \emph{References} (second reading): S.~Bloch, K.~Kato,
\emph{$L$-functions and Tamagawa numbers of motives}, §3.7.2--3 and Corollary~3.8.4, pp.~353--355; P.~Colmez,
\emph{Théorie d'Iwasawa des représentations de de Rham d'un corps local}, §§III.2 and III.4; J.-M.~Fontaine,
\emph{Représentations $p$-adiques semi-stables}, Exposé~III, Théorème~3.8(i). \emph{Relation:} the local dimension
formula applied to the integral module of cocycles (rank of the coboundaries plus \m{\dim H^1_{\mathrm f}}), with
Colmez's ring \m{B_{\max}} as period ring.}

\begin{lemma}[a power of $\ell$ bounds the defect]\label{l:defect}
Let $\cO$ be a coefficient ring, $\Gamma$ a group over~\m{\Gl} with homomorphism~$j$, $L$ a potentially
crystalline Galois lattice over~$\cO$ for~$\Gamma$, $S$ an $\cO$-submodule of the maps \m{\Gamma\to L}, \m{k\in\NN} and
\m{y\ge0} a real number. Suppose \m{\rk_{\cO}L\cdot2k\le y} and \m{\ell^{k}c\in S} for every \m{c\in\Zf(L)}. Then
\m{\nlen_{\cO}\bigl(\Zf(L)/\Zf(L)\cap S\bigr)\le[\Gl:j(\Gamma)]\cdot y}.
\end{lemma}
\begin{proof}
The quotient is a quotient of \m{\Zf(L)/\ell^{k}\Zf(L)}, which is isomorphic to \m{(\cO/\ell^{k})^{\rho}} with
\m{\rho=\rk_{\cO}\Zf(L)} and has normalised length \m{k\rho}. By Input~4 (with a count of jumps at each embedding),
\m{\rho\le(1+[\Gl:j(\Gamma)])\rk_{\cO}L\le2\,[\Gl:j(\Gamma)]\rk_{\cO}L}.
\end{proof}
\leanline{\lean{normLengthdefectle}}

In the next lemma a \emph{Kummer chart} $k$ of~$\Gamma$ (unit on Wach and Kisin modules) has a ramification
index $e_k$ and a residue degree $f_k$, a set of \emph{letters} on which \m{\Gl} acts through the roots of unity of
order \m{\ell^{f_k}-1}, and may be \emph{unramified}; a \emph{root of the letters} is a \m{\zeta\in\cO} with
\m{\zeta^{\ell^{f_k}-1}=1}. A finite set $U$ of integers \emph{lies in the window \m{(\eta,\epsilon,\beta,W)}}, for
natural numbers $\eta$, $\beta$, $W$ and a sign~$\epsilon$, if every \m{u\in U} has \m{|u|\le\eta} or
\m{\beta\le\epsilon u\le\beta+W}.

\begin{lemma}[the crystalline case]\label{l:crys}
Let $\Gamma$ be a topological group over~\m{\Gl} with homomorphism~$j$, $\cO$ a coefficient ring, $L$ a Galois
lattice over~$\cO$ for~$\Gamma$, and \m{Q,m,h,\varkappa,G,b,\delta,\eta,W} natural numbers. Let \m{\mathrm{th}} and
$\epsilon$ be functions on the \m{\Zl}-algebra homomorphisms \m{g'\colon\cO\to\Cl}, with truth values and with
signs as values, and $U$, $U'$ functions from these homomorphisms to the finite sets of integers; put
\m{\beta_{g'}=b+\delta} if \m{\mathrm{th}(g')} holds and \m{\beta_{g'}=G+\delta} otherwise. Let \m{y\ge0} be a real number,
$k$ a Kummer chart of~$\Gamma$, $\zeta$ a root of its letters in~$\cO$, and \m{T_0}, \m{s_0} functions on~\m{\Cl}
with truth values and signs as values. Suppose:
\begin{enumerate}[label=\textup{(\alph*)},leftmargin=2.3em,itemsep=1pt]
\item $0<Q$; $99\,(Q^2+(W+1)Q+Q+W+m)\le h$; $99h^3\le G$; $99\,(Gh^3+G)\le b$; $2\log(4b)<m$;
$\delta<\eta$; $99\,((\eta+W)Q+\eta+h)\le h^3$;
\item for every $g'$, \m{U(g')} lies in the window \m{(\eta,-\epsilon(g'),\beta_{g'},W)} and \m{U'(g')} in the window
\m{(\eta,\epsilon(g'),\beta_{g'},W)};
\item $L$ belongs to the class \m{\mathcal{Q}(\top,\emptyset,U')};
\item either, for every~$g'$, every negative \m{u\in U(g')\cup U'(g')} has \m{2u\le-h} and the jumps of $L$ at~$g'$ in
\m{[-\eta,\eta]} number at most~$\varkappa$ with multiplicity; or, for every~$g'$, \m{U'(g')\subseteq[-\eta,0]};
\item \m{10^{5}\bigl(\rk_{\cO}L\cdot\lfloor2\lambda/Q\rfloor\cdot m+\varkappa\lambda\bigr)\le y}, where
\m{\lambda=\lambdaLoss{h}{b}{\ell}};
\item $j$ is an open embedding, \m{e_kf_k=[\Gl:j(\Gamma)]} and \m{e_k\le h};
\item for every $g'$ with \m{\mathrm{th}(g')} or \m{Q<\ell}: \m{\mathrm{th}(g')=T_0(g'(\zeta))} and
\m{\epsilon(g')=s_0(g'(\zeta))};
\item either $k$ is unramified and \m{7\le\ell}; or \m{\ell\le Q} and $s_0$ is constant;
\item if \m{\ell\le Q}: for every \m{z\in\Cl} and all \m{1\le a,a'\le\lfloor\log_\ell(h^3)\rfloor}, \m{T_0(z)} and
\m{T_0(z^{\ell^{a}})} imply that \m{T_0(z^{\ell^{a+a'}})} does not hold.
\end{enumerate}
Then \m{\nlen_{\cO}\bigl(\Zf(L)/\Zf(L)\cap L^{\perp\mathcal{Q}(\top,\emptyset,U)}\bigr)\le[\Gl:j(\Gamma)]\cdot y}.
\end{lemma}
\begin{proof}
The lemma assembles three theorems of the note on crystalline lifts; its content lies there. Put
\m{K_0=\lfloor2\lambda/Q\rfloor}. By the definition of $\lambda$ and Legendre's formula, \m{2b<\ell(\lambda+1)}, and an
integer that divides \m{(2b)!\prod k^{h^3}} at most to the power~$\lambda$ gives control of the divisibility by~$\ell$
of the integers near~$b$. By (a), \m{\log_\ell H<m} for every \m{H\le4b}. If, for some \m{x_0\le m}, every pairing
of~$L$ with a lattice of the class \m{\mathcal{Q}(\top,\emptyset,U)} satisfies the scaled cup bounds with exponent
\m{10^{4}K_0x_0}, then \m{\ell^{10^{4}K_0x_0}\Zf(L)\subseteq L^{\perp\mathcal{Q}(\top,\emptyset,U)}} (first input), and
Lemma~\ref{l:defect} with (e) gives the claim. \emph{Case \m{\ell\le Q}.} By (g)--(i) the thick letters are sparse
(no three at short gaps) and the signs are constant; the theorem on sparse thick letters gives the scaled cup
bounds. \emph{Case \m{Q<\ell} and \m{h^3\le\lambda}.} The chart is unramified and \m{\ell\ge7} by~(h), and the theorem
for two sizes of windows applies, the inequalities following from \m{2b<\ell(\lambda+1)} and~(a). \emph{Case
\m{Q<\ell} and \m{\lambda<h^3}.} Then \m{\ell\ge9G} (from the inequalities~(a) and the case hypothesis), and
$\ell$ divides no non-zero integer in the windows in play, for two different reasons. In the low window and in
the far windows near~\m{\pm G} the integers are smaller than~$9G$ in absolute value, hence smaller than~$\ell$. In
the far windows near~\m{\pm b} the integers are factors of the product in the definition of~$\lambda$, which enters
it to the power~$h^3$: a multiple of~$\ell$ among them would give \m{\lambda\ge h^3}. The theorem on few low jumps then shows that the defect is zero unless $\lambda$ is at least \m{h/5},
and in that case bounds its length by \m{11f_k(\varkappa(2\lambda+3)+(W+1)\rk_{\cO}L)} times the length of~\m{\cO/\ell},
which is at most \m{[\Gl:j(\Gamma)]\,y} by (e), (f) and~(a); here alternative (d) is used.
\end{proof}
\leanline{\lean{crystallinedefectle}}

\begin{proof}[Proof of Theorem~E]
Choose a Kummer chart $k$ of~$\Gamma$ with \m{e_kf_k=[\Gl:j(\Gamma)]}, unramified when \m{j(\Gamma)} contains the
inertia group (an input of the unit on Wach and Kisin modules), and let $\cO$ be large enough to contain a
primitive root of unity $\zeta$ of order \m{\ell^{f_k}-1}. Let $j_0$ be the index of (R2) for~$\HT$. By
Lemma~\ref{l:coeffsurj} every embedding \m{g'\colon\cO\to\Cl} is \m{\coeffEmb{\cO}{\rho(g')}} for some
\m{\rho(g')\in\Gl}; put \m{\sigma_{g'}=\complexEmbOf{e}{\rho(g')}}. By (ii), (iii) and (R2), for every~$g'$ the set
\m{U(g')} lies in the window
\[
\bigl(\hhD{D}^2+1,\ -\cmSign{\sigma_{g'}},\ \beta+j_0\hD{D}+\hhD{D}-2\tD{D},\ 4\tD{D}\bigr),
\]
with \m{\beta=\bone{D}} if \m{[\sigma_{g'}]\in\thick{D}} and \m{\beta=\btwo{D}} otherwise, because the weight \m{\wt{D}{j_0}{\sigma}} is \m{\bone{D}} or \m{\btwo{D}} plus \m{j_0\hD{D}+\hhD{D}}, and the
string jumps \m{0\le c<t} are low; the set \m{\{-u-1:u\in U(g')\}}, which contains the jumps of the twisted
duals, lies in the window with the opposite sign.

\emph{The tame case: \m{\ell=\Dg{\nu}{D}}.} Then \m{\ell\equiv1} modulo \m{\ND{D}d}, so every element of~\m{\Gl} fixes the
roots of unity of that order, hence \m{e(\ED{D})} (Lemma~\ref{l:fixemb}): $j$ is a homeomorphism onto~\m{\Gl}. The
theorem on tame types of the note on crystalline lifts gives an explicit exponent $k_1$ with
\m{\ell^{k_1}\Zf(L_c^{\vee}(1))} in the cup-orthogonal, and Lemma~\ref{l:defect} applies:
\m{2k_1\rk_{\cO}L_c\le\rk_{\cO}L_c\cdot\loss{D}{\ell}/x} because \m{v_\ell(\gD{D})\ge1}, so that
\m{\loss{D}{\ell}\ge10^{6}\,\Dm{D}\bone{D}/\kD{D}}, and \m{x\le\tD{D}}.

\emph{The crystalline case: \m{F=\top} and $K$ empty.} Apply Lemma~\ref{l:crys} to \m{L_c^{\vee}(1)} with
\[
Q=\kD{D}\tD{D},\quad m=\Dm{D},\quad h=\hhD{D},\quad\varkappa=\Bigl\lfloor\frac{\rk_{\cO}L_c\cdot\Dm{D}}{\kD{D}\,x}\Bigr\rfloor,\quad
G=\btwo{D},
\]
\m{b=\bone{D}}, \m{\delta=j_0\hD{D}+\hhD{D}-2\tD{D}}, \m{\eta=\hhD{D}^2+1}, \m{W=4\tD{D}}, \m{\mathrm{th}(g')} the statement that
\m{[\sigma_{g'}]} is thick, and \m{\epsilon(g')=\cmSign{\sigma_{g'}}}. Condition~(a) consists of inequalities of the data
(Chapter~III). (b) was shown; (c) is the hypothesis on $L_c$, dualised. (d): for \m{c=\ast} the first alternative
holds, the bound on the low jumps being~(R3); for a string index the dual of~$L_c$ has the single jump
\m{-c-1\in[-\eta,0]}. (e) is a comparison with Definition~\ref{d:loss}:
\m{10^{5}(\rk\cdot\tfrac{2\lambda}{\kD{D}\tD{D}}\Dm{D}+\varkappa\lambda)\le3\cdot10^{5}\,\rk\cdot\tfrac{\Dm{D}}{\kD{D}}\cdot\tfrac{\lambda}{x}}.
(f): \m{e_k\le e_kf_k=[\Gl:j(\Gamma)]}, the number of embeddings \m{\complexEmbOf{e}{g}} (Lemma~\ref{l:small1}(d)), at most
\m{[\ED{D}:\mathbb{Q}]\le\hhD{D}}. (g)--(i): thickness and sign of \m{\sigma_{g'}} must be functions of \m{g'(\zeta)}.
If $\ell$ does not divide \m{\ND{D}d}, the chart is unramified (inertia fixes the roots of unity of order
\m{\ND{D}d}, hence \m{e(\ED{D})}), \m{\ell\ge7} because \m{2\cdot3\cdot5} divides~$N_0$, and \m{g'(\zeta)} determines the
action of \m{\rho(g')} on the roots of unity of order prime to~$\ell$, hence \m{\sigma_{g'}} (Lemma~\ref{l:small1}(c)).
If $\ell$ divides \m{\ND{D}d}, then \m{\ell\le\kD{D}\tD{D}=Q}; thickness is by definition the property
``\m{g'(\zeta)=z} for some $g'$ with \m{[\sigma_{g'}]} thick'', the sign is constant on thick places by
Lemma~\ref{l:cmthick}, and (i) is Lemma~\ref{l:spacing}: a Frobenius element of the chart raises the roots of unity
of order prime to~$\ell$ to the power \m{\ell^{f_k}}, the exponents \m{a,a'} are at most
\m{\log_\ell(\hhD{D}^3)\le12\Dmu{D}}, and \m{a+a'\le27\Dmu{D}}. The formal proof of this theorem has seventy
lines, mostly inequalities between the numbers of the data; the paragraph above is its outline.
\end{proof}

\sect{Where the results are used}{s:use}
All results of this note are used inside the unit, by the two notes that follow. Theorem~A is the source of the
pencils of tied primes in the note \emph{Short Selmer groups after one anticyclotomic twist}, which also uses
Lemma~\ref{l:corep} (to read Selmer modules in finitely many coordinates), Lemma~\ref{l:graph} and
Lemma~\ref{l:rank} (for balanced rows). The note \emph{A fit family of twists exists} uses Theorem~B and
Lemmas \ref{l:twistav}, \ref{l:aptaux} to show that the auxiliary pairs have rows; Theorem~C and
Lemmas \ref{l:gram}, \ref{l:isotypic}, \ref{l:noeig}, \ref{l:string}, \ref{l:heckestring} in the bound for the denominator of the
projection to lifts; Theorem~D and Lemma~\ref{l:tamechar} for the heights of the ratios of $L$-values;
Theorem~E, with Lemmas \ref{l:sumemb}, \ref{l:newcyc} and~\ref{l:newsel}, to show that the crystalline losses above~$\ell$
are affordable. The definitions of §\ref{s:defs} are the vocabulary of both notes.

\sect{Notes on the formal text}{s:notes}
\begin{enumerate}[label=\textup{\arabic*.},leftmargin=1.6em,itemsep=5pt]
\item \emph{Letters.} The formal text writes {\leanfont Ω} for~\m{\Qb}, {\leanfont ιℓ} for~$\iota$, {\leanfont O}
both for a coefficient ring (here~$\cO$) and for an auxiliary pair (here~$O$), {\leanfont A} for a place
of~\m{\Qb} (here~$\fP$), {\leanfont P} for an auxiliary place (here~$\fq$) as well as for the parameters,
{\leanfont J} both for the jumps (here~$\HT$) and for sets of places, {\leanfont ι} for the index set of the rows
of a pencil (here~$\mathcal{I}$), {\leanfont star w} for~\m{\cj{w}}, {\leanfont conjGal g} for~\m{g^{\mathsf c}},
{\leanfont idOrConjInv b g} for~\m{g^{[b]}}, {\leanfont ellVal} for~\m{\vl}, {\leanfont Gp ℓ} for~\m{\Gl}, and
{\leanfont ℓ'}, {\leanfont E₁} for the prime and the field of Theorem~A.
\item \emph{The group \m{\mathbb{Z}/q}.} The homomorphisms $h$ of Definition~\ref{d:auxdata} take values in the
multiplicative copy of \m{\mathbb{Z}/q}; ``\m{h(\gamma)=1}'' above is the generator, ``\m{h(\phi)=0}'' the neutral
element.
\item \emph{Truncations.} In Definition~\ref{d:csd} the exponent is \m{t-1} in natural numbers, so \m{t=0} reads as
\m{t=1}. The rank difference \m{\#n-\#r}, the lower end \m{b-30h^3} in Definition~\ref{d:loss} and the differences in
Lemma~\ref{l:twistav} are truncated at~$0$. The valuation \m{\vl} is truncated: it is $0$ for \m{|x|\ge1}. In
Definitions \ref{d:cr} and~\ref{d:places} the number $\ell$ may be any natural number.
\item \emph{Places of~\m{\Qb}.} The formal type of places contains the trivial valuation ring \m{\Qb} itself; no
finite place lies under it.
\item \emph{Arguments.} The arguments of the declarations were read from the text of the files; for 35
declarations of this note the dossier gives a list reconstructed by hand (section variables declared twice,
variables that occur only as names of fields or bound names, statements that contain a local definition), and
no list has been confirmed by Lean's own answer, since the artifact was not compiled here. Three readings of
the program differ from the mathematics and were kept as slots or words: the element of the Galois group is
counted as an argument of \m{\blockRep{T}{i}{\cdot}} and of \m{\invAvatar{\vartheta}{\cdot}} (an arrow read as an argument);
the point character of Definition~\ref{d:point} counts the hypothesis \m{y_j^{q}=1} and the group element as
arguments, and is therefore written out; in Theorem~E the quantifier ``for all large~$\cO$'' is read as two
arguments. Instance arguments of definitions that the definition may not use are not listed.
\item \emph{Statements with local definitions.} Lemmas \ref{l:avconj}, \ref{l:avlabel}, \ref{l:inertia} and Theorems C
and~E contain a local definition in their formal statement, at which the census of the workspace cuts the
statement; they were read in full in the source (the second reading of this write-up lists their true ends).
\item \emph{Hypotheses with defaults, and hypotheses not used.} In Lemma~\ref{l:newsel} the last two hypotheses
have default proofs in the formal text (a tactic closes them at the uses). The symmetry of
\m{\tupleProj{O}{c}} and the fact that it fixes its image (Lemma~\ref{l:small3}(c)) are stated in the formal text
without the hypothesis \m{\IsApt{O}{c}{\vartheta}}, in its namespace. In Lemma~\ref{l:twistav} and
Lemma~\ref{l:prodlabel} the apt family $c'$ serves only to identify the chosen places of~$O$ with those of~$D$.
In Theorem~D the index set of the family is the product of a set with two elements and the group of signs.
\item \emph{Theorem~A.} The formal statement has two primes in scope; only the one of the conclusion matters.
The set $R$ is a set with a proof of finiteness, $S$ a finite set.
\item \emph{Assemblies and outlines.} Lemma~\ref{l:crys} and the proof of Theorem~E apply theorems of the unit
on lifting to the numbers of the data; Lemma~\ref{l:measure} in the case \m{z_p=1} and Theorem~C unfold
definitions of the units on forms. For these four, and for Lemma~\ref{l:corep}, the proofs above are outlines
of the formal proofs, as said at the head or at the end of each; the inequalities and the identities between adelic sums were not re-derived by hand.
\item \emph{Differences from the dossier.} The dossier's reading of ``from $\mathrm{U}(1)$'' says that the idele
character is trivial on the ideles that are units at \m{S\cup\cj{S}}; the formal definition asks triviality on the
ideles that are units \emph{outside} \m{S\cup\cj{S}} (and congruent to~$1$ at the modulus), which is the reading
of §2.2 and of the authors' comment (``trivial at $S$ and~$\ell$'').
\item \emph{What the final theorems use.} All the lemmas of this note are used by declarations of the two notes
that follow (§\ref{s:use}); whether those are reached by the final theorems is recorded with them.
\end{enumerate}

\subhead{How this note was written.} From the dossier of the unit (its first chunk): the statements in words and
the structured proofs, which a program compared line by line with the source (\UnitSteps{} leaf steps in the
\UnitProofs{} proofs of the whole unit). I read the fifteen Lean files of these declarations completely, and in
the background library and in other units the definitions that the statements rest on: Selmer modules,
relations and the characters from $\mathrm{U}(1)$; Galois lattices, classes, cup-orthogonals, jumps and
crystalline cocycles; tame characters, twists, pairs, ramified pairs and the ratio of $L$-values; the number
\m{\Mom{P}{\Dmu{D}}} and the divisibilities at the chosen places. The statements of this note were compared with the Lean
statements argument by argument, and the order of the quantifiers is that of the formal text. The proofs follow
the formal ones, with these differences: the proofs of Lemma~\ref{l:small2}(d) and of Lemma~\ref{l:noeig} are mine
(the formal proofs quote a lemma on roots of unity, and a count of the roots of a polynomial); Lemma~\ref{l:corep},
Theorems C and~E and Lemmas \ref{l:measure}, \ref{l:crys} are given as outlines. I did not open the proofs of the
inputs. The references and their places are those of the second reading of this write-up and of its catalogue;
I did not check them in the literature.

\head{Appendix A. Concordance with the Lean artifact}
{\fontsize{8.6}{11.6}\selectfont\renewcommand{\leanfont}{\ttfamily\fontsize{7.7}{10}\selectfont\color{leancol}}
\setlength{\LTleft}{0pt}\setlength{\LTright}{\fill}\setlength{\tabcolsep}{0pt}
\begin{longtable}{@{}p{30mm}p{56mm}p{43mm}r@{}}
\textbf{In this note} & \textbf{Lean declaration} & \textbf{File} & \textbf{Lines}\\\midrule\endhead
\ConcordanceRows
\end{longtable}}

\subhead{Three statements as they are written in the artifact.}
\quoted{existsisAnticyclotomicAt}
\quoted{isRow}
\quoted{SelmerLengthLe}
In these texts {\leanfont Gal(Ω/E₁)} is \m{\GE{E}}, {\leanfont AbsGal D} is \m{\AbsGal{D}}, {\leanfont placesDividing E ℓ}
is \m{\pd{E}{\ell}}, {\leanfont placesAboveG D} is \m{\placesAboveG{D}}, {\leanfont .conjClosure S} is
\m{S\cup\cj{S}}, {\leanfont inertiaAbove} is the inertia above a set of places, {\leanfont selmerModule ρ S L B} is
\m{\Sel{\varrho}{S}{L}{B}}, {\leanfont QuotSpan a} is \m{\cO/a}, {\leanfont crysRels} gives the crystalline relations,
{\leanfont normLength} is \m{\nlen_{\cO}}, and {\leanfont ∀ᶠ O in atTop} is ``for all large~$\cO$''.

\head{Notation}
\begingroup\fontsize{8.6}{10.9}\selectfont\raggedright
\setlength{\columnsep}{9mm}\setlength{\parskip}{0pt}
\begin{multicols}{2}
\nt{$\ell$, $\Ql$, $\Zl$, $\Qlb$, $\Zlb$, $\Cl$, $\Gl$}{a prime; the $\ell$-adic numbers, integers, an algebraic closure, its integers, its completion, \m{\Gal(\Qlb/\Ql)}}
\nt{$\iota$}{a fixed isomorphism \m{\CC\to\Qlb}}
\nt{$\Qb$, $\mathsf{c}$, $g^{\mathsf c}$, $g^{[b]}$}{an algebraic closure of~$\mathbb{Q}$; complex conjugation; \m{\mathsf{c}g\mathsf{c}}; $g$ or \m{(g^{\mathsf c})^{-1}}}
\nt{$E$, $E'$, $E_1$, $\GE{E}$, $f_E$}{CM fields as in the conventions; \m{\Gal(\Qb/E)}; the number of infinite places}
\nt{$P$, $D$, $N_0$, $d$}{parameters and admissible data; the two parameters}
\nt{$\ED{D}$, $\Dm{D}$, $\Dr{D}$, $\kD{D}$, $\tD{D}$, $\fD{D}$}{field, number of lifts, ranks, step, number of infinite places of the data}
\nt{$\Dg{i}{D}$, $\gD{D}$, $\lvl{D}$, $\pi_g$}{the two auxiliary primes, their product, the level, an integer of norm \m{\gD{D}}}
\nt{$\hhD{D}$, $\hD{D}$, $\bone{D}$, $\btwo{D}$, $\wt{D}{j}{\sigma}$}{half-gap, gap, the two rounded quotients, the weight of lift~$j$}
\nt{$\ND{D}$, $\Dmu{D}$, $\varsigma$, \m{\Mom{P}{\Dmu{D}}}, $\omega_D$}{order of the root of unity, exponent, the root of unity, the order and the family of residue characters}
\nt{$\ordch{D}$, $\thick{D}$, $\operatorname{cl}_m$}{the order of the chosen groups; the thick places; the class of a place modulo~$m$}
\nt{$w$, $v$, $\fq$, $\fq_0$, $\cj{w}$, $\Nm(w)$, $\F_w$, $\varpi_w$}{finite places, the conjugate, the norm, the residue field, a uniformiser}
\nt{$\fP$, $\fP_w$, $\fP^{\mathsf c}$, $G_{\fP}$, $I_{\fP}$, $\tau_{\fP}$, $\tau_w$}{places of~\m{\Qb}; the chosen one over~$w$; the conjugate; decomposition and inertia groups; tame generators}
\nt{$\pd{E}{M}$, $\inab{S}$, $S^{\complement}$}{finite places containing~$M$; inertia above~$S$; the complement of~$S$}
\nt{$\cyc$, $\cyc_M$, $\Art$, $\upsilon$}{cyclotomic characters; the Artin map; the Hecke character of the Weil data}
\nt{$\cO$, $\cO'$, $\cO_0$, $\varpi$, $\cO(-e)$}{coefficient rings; a uniformiser; a Tate twist}
\nt{$A$, $A'$, $\Lambda$, $B$, $\Gamma$, $v_A$, $\kappa_A$}{commutative rings, a module, a group; valuation and residue field of~$A$}
\nt{$\vl$, $v_l$, $\len$, $\nlen$, $\rk$, $\tr$, $\det$, $\diag$}{valuations, lengths, rank, trace, determinant, diagonal matrix}
\nt{$M_n(A)$, $\GL_n(A)$, $\#n$}{matrices and invertible matrices indexed by a finite set~$n$; its cardinality}
\nt{$\theta$, $\delta_g$, $\Ob{L}{B}$, $\Rel{Z}$, $\Coc{\varrho}$}{a relation; the indicator relation; functions obeying~$L$; relations of~$Z$; cocycle relations}
\nt{$\Sel{\varrho}{S}{L}{B}$}{the Selmer module}
\nt{$R$, $u$, $\varrho$, $T$, $\psi$, $\psi'$, $\chi$, $\eta$}{representations, lattices, characters, idele characters; $T$ also a faithful character and a set of places}
\nt{$\HT$, $\HT'$, $\HT_{\kappa,N}$, $\kappa$, $\lambda_i$}{types; the type of a weight $\kappa$ in rank~$N$; its partition}
\nt{$\Zf(L)$, $L^{\vee}(1)$, $L^{\perp\mathcal{Q}}$, $\mathcal{Q}(F,K,U)$}{crystalline cocycles, twisted dual, cup-orthogonal, a class of lattices}
\nt{$\mathfrak{M}$, $\varkappa$, $N'$, $S_u$, $c'_i$, $\mathsf{e}$, $k_3$}{a multiset (of jumps, of eigenvalues); the bound for the number of low jumps in Lemma~\ref{l:crys}; the integer of Lemma~\ref{l:measure}; the set of places and the polynomials of Theorem~C(3); the averaging idempotent in Lemma~\ref{l:aptaux}; the exponent of~$3$ in \m{\ND{D}}}
\nt{$\mu$, $\psi_w$, $e_v$, $\mu[k,j]$, $\av$}{a character of~$\mathrm{U}(1)$, its tame part and weight, the algebraic Hecke character, the avatar}
\nt{$\vartheta$, $\tau$, $Q(\vartheta)$, $\psi_\vartheta$}{tame characters and twists ($\tau$ also a family of tame generators); the set of places; the tame part}
\nt{$p$, $\nu_p$, $z_p$, $\varrho_p(\vartheta,s)$, $y_v(p,\vartheta)$}{a pair; its ratio of $L$-values; a value at a uniformiser}
\nt{$O$, $V$, $I_O$, $w_i$, $\mu_O$, $\psi_O$, $a_O$, $\mathcal{F}(O)$}{auxiliary pairs, chosen index set and places, characters, forms}
\nt{$O_D(\vartheta,j)$, $O_p(\vartheta)$, $\Xi_p(\vartheta)$}{the auxiliary pairs of the data and of a pair; the lift of the trivial character}
\nt{$c$, $c'$, $\tilde c$, $c_i$, \m{\tupleIdem{r}{c}}, $x_\chi$, $\xi$}{families of polynomials, the idempotent of a tuple of labels, the name of a character, a root}
\nt{$H$, $H_w(X)$, $H_x(w)$, $T_{w,i}$, $S$}{systems of Hecke polynomials and Hecke operators; $S$ also a set of places and a polynomial}
\nt{$x$, $\varrho_x$, $f$, $y$, $z$}{an isotypic space or eigensystem and its representation; forms; Hecke data; also elements}
\nt{\m{\obeyingCocycles{u}{\varrho}{a}}, \m{\RowBlockIndex{r}{i}}, \m{\blockRep{T}{i}{g}}}{Definitions \ref{d:obey}, \ref{d:blocks}}
\nt{\m{\coeffEmb{\cO}{g}}, \m{\complexEmbOf{e}{g}}, \m{\ZfAt{v}{u}}, \m{\crysRels{u}}, \m{\blockLoss{u}{w}{M}}}{Definitions \ref{d:coeffemb}--\ref{d:blockloss}}
\nt{\m{\cmSign{\sigma}}, \m{\IsRowType{D}{r}{t}{\HT}}, \m{\lambdaLoss{h}{b}{l}}, \m{\loss{D}{l}}}{Definitions \ref{d:cmsign}--\ref{d:loss}}
\nt{\m{\AbsGal{D}}, \m{\placesAboveG{D}}, \m{\badPlaces{D}{\ell}}, \m{\invAvatar{\vartheta}{g}}}{Definitions \ref{d:places}, \ref{d:invav}}
\nt{\m{\IsConjSelfDual{t}{R}}, \m{\crudeUnit{\ell}{D}}, \m{\SelmerLengthLe{\psi}{R}{Q}{m}{\gamma}{z}}}{Definitions \ref{d:csd}--\ref{d:short}}
\nt{\m{\IsColumn{D}{t}{\HT}{S}{R}}, \m{\IsRow{D}{t}{\HT}{R}}}{columns and rows (Definitions \ref{d:column}, \ref{d:row})}
\nt{\m{\chiOneInv{O}}, \m{\normalisedRep{O}{R}}, \m{\twistedLabel{O}{c}{i}}, \m{\tupleProj{O}{c}}}{Definitions \ref{d:chione}--\ref{d:tproj}}
\nt{\m{\IsApt{O}{c}{\vartheta}}, \m{\HasRow{O}{z}{R}}, \m{\IsLabelledColumn{O}{k}{\HT}{c}{S}{R}}}{Definitions \ref{d:apt}, \ref{d:hasrow}, \ref{d:labcol}}
\nt{\m{\AuxPrimeData{E}{q}}, \m{\IsAnticyclotomicAt{T}{\fq}{l}}, \m{\IsAuxiliary{S}{I}}, \m{\IsFamily{S}{\nu}{\mathrm{val}}}}{Definitions \ref{d:auxdata}--\ref{d:family}}
\nt{\m{\evenClasses{\cO}{S}{\nu}{e}}, \m{\IsNewPlaces{\nu}{w}{\tau}{\phi}}, \m{\withConjInv{\tau}{k}}}{Definitions \ref{d:even}, \ref{d:new}}
\nt{\m{\IsPrimeTriangle{\nu}{w}{\tau}{\phi}{a}{t}{\mathrm{ty}}{h}{c}}}{prime triangles (Definition~\ref{d:tri})}
\nt{\m{\IsAuxTwist{S}{Q}{\psi}{c}}, \m{\unionWithConj{S}{J}}}{Definitions \ref{d:auxtw}, \ref{d:union}}
\nt{\m{\twistAtPoint{h}{y}{e}{\varrho}{g}}, \m{\reduceModPow{l}{M}{\varrho}}}{Definition~\ref{d:point}}
\nt{\m{\IsBalanced{\ell}{S}{\varrho}{L}{\gamma}{\beta}{M}}}{balanced representations (Definition~\ref{d:bal})}
\nt{\m{\IsPencil{\ell}{S_0}{U}{W}{N}{\fq_0}{J}{x}}, \m{\tieRelations{J}{x}{i}}}{pencils and ties (Definitions \ref{d:pencil}, \ref{d:ties})}
\nt{\m{\matrixLattice{\varrho}}}{the lattice of a matrix representation (Definition~\ref{d:matlat})}
\nt{$l=(\fP,\gamma,\phi,h)$, $l_v$, $s_v$, $s_{v,i}$, $\mathcal{I}$}{data at an auxiliary prime; the slopes; the index set of the rows}
\nt{$\gamma$, $\phi$, $\sigma_0$, $g$, $g'$, $g_1$, $s$, $a$, $z$}{elements of Galois groups; $s$ also a sign, $a$ also an exponent or a bound}
\nt{$\sigma$, $\bar\sigma$, $[\sigma]$, $\mathfrak{S}_v$, $\sigma_{g'}$}{complex embeddings, the place, the CM type, sets of embeddings; $\sigma$ also a finite set and a family}
\nt{$j$, $j_v$, $j_{\fP}$, $e$, $e_v$, $\rho(g')$}{homomorphisms to~\m{\Gl}, embeddings into~\m{\Qlb}, a preimage; $j$, $e$ also integers}
\nt{$n$, $r$, $m$, $d$, $t$, $k$, $q$, $N$, $M$, $i$, $b$, $h$}{finite sets or natural numbers: index sets, ranks, exponents, moduli, signs}
\nt{$I$, $J$, $K$, $Q$, $S_0$, $T_0$, $W$, $X$, $\mathcal{K}$, $\mathcal{U}$}{sets of places, index sets, neighbourhoods; $X$ also the variable of polynomials, $W$ a window width, $Q$ a natural number in §\ref{s:loss}}
\nt{$\nu$, $\mathrm{val}$, $\alpha$, $c_a$, $g_a$, $F$, $u_a$}{a family of characters and of values; a presentation; an integral measure ($\nu$ also a natural number)}
\nt{$\mathbb{M}$, $\mathcal{L}$, $\mathcal{N}$, $\Lambda_0$, $Y$, $k_0$, $\theta^{g}$}{objects of the proof of Lemma~\ref{l:graph}; $Y$ also a matrix}
\nt{$V$, $W$, $\mathcal{A}$, $Z$, $Z'$, $\eta_a$, $x_b$, $\omega_k$, $l_{p,k}$, $\lambda_k$, $d_p$, $P_{b,k}$, $w_{p,b}$, $G_p$}{objects of Lemma~\ref{l:gram}; $Z$, $Z'$ also spaces of forms}
\nt{$T_0$, $\mathfrak{s}$, $y_u$, $u_0$, $\rho_0$, $f_0$, $\mathcal{R}$, $\rho_k$, $U$}{objects of the proofs of §\ref{s:aux}; $U$ also a subgroup and a function to finite sets}
\nt{$\gamma_0$, $g_{I,a}$, $c_{I,a}$, $u_{I,a}$, $y_{I,a}$, $u_{\mathrm f}$, $\Lambda(\vartheta)$, $\mathrm{den}$}{objects of §\ref{s:families}}
\nt{$\zeta$, $p_0$, $F_1$, $F_1'$, $F_2$, $F_3$, $\delta_i$, $\alpha_i$, $d_0$, $c_i(Y)$, $H_w(i)$}{roots of unity, a prime, polynomials, a sequence, eigenvalues, coefficients}
\nt{$G$, $\delta$, $\eta$, $\epsilon$, $\mathrm{th}$, $\beta$, $\beta_{g'}$, $\lambda$, $s_0$, $K_0$, $x_0$, $k_1$, $j_0$, $c_0$, $e_k$, $f_k$, $\rho$}{numbers and functions of §\ref{s:loss}; $G$ also a finite group or an operator, $\delta$ also the function of (T5)}
\nt{$\varepsilon$, $\pi$, $C$, $\top$, $\ast$, $\mathrm{ty}$, $\mathrm{U}(1)$}{a sign; the constant in \m{2\pi i}; a monoid; the trivial filter; the index of the row block; the map of a triangle; the unitary group of rank one}
\end{multicols}
\endgroup

\head{References}
{\fontsize{9.4}{12.6}\selectfont
The references are those named by the second reading of this write-up and by its catalogue of quoted results, with
the places given there; I did not open them.\par
\smallskip
S.~Bloch, K.~Kato, \emph{$L$-functions and Tamagawa numbers of motives}, The Grothendieck Festschrift~I (1990).\par
G.~Chenevier, M.~Harris, \emph{Construction of automorphic Galois representations, II}, Camb. J. Math. 1 (2013).\par
P.~Colmez, \emph{Théorie d'Iwasawa des représentations de de Rham d'un corps local}.\par
J.-M.~Fontaine, \emph{Représentations $p$-adiques semi-stables}, Astérisque 223 (1994).\par
M.~Harris, K.-W.~Lan, R.~Taylor, J.~Thorne, \emph{On the rigid cohomology of certain Shimura varieties}.\par
J.-P.~Labesse, \emph{Changement de base CM et séries discrètes} (2011).\par
J.~S.~Milne, \emph{Class Field Theory}, course notes, version 4.03 (2020).\par}

\end{document}
